%%% Exemplo de utilização da classe ITA
%%%																															
%%%   por        Fábio Fagundes Silveira   -  ffs [at] ita [dot] br
%%%              Benedito C. O. Maciel     -  bcmaciel [at] ita [dot] br
%%%              Giovani Volnei Meinertz   -  giovani [at] ita [dot] br
%%%    	         Hudson Alberto Bode       -  bode [at] ita [dot]br
%%%    	         P. I. Braga de Queiroz    -  pi [at] ita [dot] br
%%%    	         Jorge A. B. Gripp         -  gripp [at] ita [dot] br
%%%    	         Juliano Monte-Mor         -  jamontemor [at] yahoo [dot] com [dot] br
%%%    	         Tarcisio A. B. Gripp      -  tarcisio.gripp [at] gmail [dot] com
%%%    	         
%%%
%%%  IMPORTANTE: O texto contido neste exemplo nao significa absolutamente nada.  :-)
%%%              O intuito aqui eh demonstrar os comandos criados na classe e suas
%%%              respectivas utilizacoes.
%%%
%%%  Tese.tex  2016-08-25
%%%  $HeadURL: http://www.apgita.org.br/apgita/teses-e-latex.php $
%%%
%%% ITALUS
%%% Instituto Tecnológico de Aeronáutica --- ITA, Sao Jose dos Campos, Brasil
%%%                   http://groups.yahoo.com/group/italus/
%%% Discussion list: italus {at} yahoogroups.com
%%%
%++++++++++++++++++++++++++++++++++++++++++++++++++++++++++++++++++++++++++++++
% Para alterar o TIPO DE DOCUMENTO, preencher a linha abaixo \documentclass[?]{?}
%   \documentclass[tg]{ita}			= Trabalho de Graduacao
%   \documentclass[tgfem]{ita}	= Para Engenheiras
%   								msc     		= Dissertacao de Mestrado
%   								mscfem   		= Para Mestras
%   								dsc      		= Tese de Doutorado
%   								dscfem   		= Para Doutoras
%   								quali    		= Exame de Qualificacao
%   								qualifem 		= Exame de Qualificacao para Doutoras
% Para 'Draft Version'/'Versao Preliminar' com data no rodape, adicionar 'dv':
%   \documentclass[dsc, dv]{ita} 
% Para trabalhos em Inglês, adicionar 'eng':
%   \documentclass[dsc, eng]{ita}
%		\documentclass[dsc, eng, dv]{ita}
%++++++++++++++++++++++++++++++++++++++++++++++++++++++++++++++++++++++++++++++
%\PassOptionsToPackage{backref=page}{hyperref}

%\documentclass[dsc]{ita}    % ITA.cls based on standard book.cls
% Quando alterar a classe, por exemplo de [msc] para [msc, eng]) rode mais uma vez o botão BUILD OUTPUT caso haja erro

\documentclass[msc]{ita}    % ITA.cls based on standard book.cls

\usepackage[utf8]{inputenc} % Pacote para permitir o uso de acentos
% Ajustes de hifenização
\tolerance=1
\emergencystretch=\maxdimen
\hyphenpenalty=10000
\hbadness=10000
\usepackage{lipsum} % Pacote para gerar texto aleatório
\usepackage{ragged2e} % Pacote para mais opções de alinhamento
\usepackage{graphicx} % Pacote para inserir figuras jpeg, pdf e png
\usepackage{ae}
\usepackage{epsfig}
\usepackage{mathtools}
\usepackage[export]{adjustbox}
\usepackage{lineno}
\usepackage{amsthm}
\usepackage{commath}
\usepackage{color, colortbl}
\usepackage{bbm}
\usepackage{setspace} % Pacote para ajustar espaçamento entre linhas

% Instructions to prevent LaTeX from hyphenating words
\tolerance=1
\emergencystretch=\maxdimen
\hyphenpenalty=10000
\hbadness=10000

\graphicspath{{figures/}} % Indica o diretório no qual as imagens estão sendo salvas, evitando que precisamos digitar parte do caminho da pasta que contém as imagens
\usepackage[rightcaption]{sidecap} % Pacote para mais opções de figura, colocar, título da legenda da figura ao lado direito
\usepackage{url} % Pacote para citar referências de sites
\usepackage{makeidx} % Pacote para criar índice
\makeindex % Habilitar a criação do índice
\usepackage{longtable} % Pacote para usar tabelas longas

% ----------
% Pacotes para dar quebra de linha dentro da tabela:
%\usepackage{fourier} 
%\usepackage{array}
\usepackage{makecell}

\renewcommand\theadalign{bc}
\renewcommand\theadfont{\bfseries}
\renewcommand\theadgape{\Gape[4pt]}
\renewcommand\cellgape{\Gape[4pt]}
% ----------

\usepackage{wrapfig} % Pacote para alinhar figuras à esquerda ou à direita do texto
\usepackage{ae}
\usepackage{epsfig}
\usepackage{mathtools}
\usepackage[export]{adjustbox}
\usepackage{lineno}
\usepackage{amsfonts}
\usepackage{amsthm}
\usepackage{commath}
\usepackage{color, colortbl}
\usepackage{bbm}

% ----------
% Pacote para inserir cabeçalho e rodapé:
% Atenção! Se remover o cabeçalho e rodapé, o LaTeX coloca como cabeçalho o padrão dele, que é bem interessante também!
% \usepackage{fancyhdr} % Pacote para inserir cabeçalho e rodapé
% \pagestyle{fancy} % Estilo do cabeçalho e rodapé, neste caso apenas uma linha no cabeçalho
%\lhead{{\footnotesize Prandiano\textregistered}} % Inserir texto no canto esquerdo do cabeçalho
%\chead{centro} % Inserir texto no centro do cabeçalho
%\rhead{{\footnotesize Curso de Python}} % Inserir texto no canto direito do cabeçalho
%\lfoot{{\footnotesize Copyright \textcopyright 2018 Prandiano}} % Inserir texto no canto esquerdo do rodapé
%\cfoot{centro} % Inserir texto no centro do rodapé
%\rfoot{{\footnotesize Reprodução proibida: Lei 9.610 de 19/02/1998}} % Inserir texto no canto direito do rodapé
% O comando \renewcommand re-nova o parâmetro especificado
%\renewcommand{\footrulewidth}{0.4pt} % Inserir uma linha acima do rodapé
% ----------

% ----------
% Pacote para inserir múltiplas figuras:
\usepackage{subfigure} % Pacote para inserir múltiplas figuras
% \usepackage{subcaption} % Pacote incompatível com subfigure
%\usepackage{caption}
% ----------

\usepackage{multicol} % Pacote que permite usar multicolunas
\usepackage{float} % Pacote para usar o "H" do figure
\usepackage[dvipsnames]{xcolor} % Pacote de cores para textos e possivelmente outros
\usepackage{amsmath} % Pacote para usar o cases
\usepackage[makeroom]{cancel} % Pacote para usar o cancelar (em e embaixo)
\usepackage{enumitem} % Pacote para enumerar itens
\usepackage{amssymb} % Pacote para inserir os símbolos dos conjuntos numéricos N, Z, Q, R, C
\usepackage[T1]{fontenc}
\usepackage{enumitem}
\usepackage{textcomp}
\usepackage{pdfpages} % Pacote para inserir página em PDF
\usepackage{courier} % Pacote para usar a fonte courier
%\usepackage{matlab-prettifier}
\usepackage{lscape} % Pacote para mudar sentido da página
\newtheorem{theorem}{Teorema} % Pacote para inserir Teorema
\newtheorem{corollary}{Corollary} % Pacote para inserir Corolário
\newtheorem{remark}{Remark}
\newtheorem{mytheorem}{Definition}
\usepackage{hyperref} % Pacote para inserir hyperref
\usepackage{multirow} % Pacote para usar os comandos \multirow
\usepackage{pdfpages} % Pacote para incluir PDF's
\usepackage{listingsutf8} % Pacote para inserir arquivos externos tipo .txt
\usepackage{etoolbox}
\AtBeginEnvironment{verbatim}{\footnotesize}

\usepackage{standalone} % Pacote para inserir sub-arquivos LaTeX
\usepackage{bookmark} % Pacote para NÃO tratar o Índice como Apêndice e colocar letra

% Nome da listagem em português
\renewcommand{\lstlistingname}{Listagem}

% REFERÊNCIAS: ----------
%\usepackage[numbers, compress]{natbib} % Pacote para usar as referências bibliográficas
%\usepackage{usebib}
% \bibliographystyle{references} % Estilo das referências
%\setcitestyle{authoryear,aysep={},open={(},close={)}}
%\newcommand{\myand}{e}

% ----------


% PARÂMETROS DA FUNÇÃO LISTINGS: ----------
\lstset
{
    basicstyle=\small\ttfamily,
    commentstyle=\color{Green},
    keywordstyle=\color{Cerulean},
    frame=single,
    language=Python,
    morekeywords={True, False},
    numbers=left,
    numbersep=10pt,
    numberstyle=\footnotesize\color{Gray},
    showstringspaces=false,
    stringstyle=\color{Mulberry},
    tabsize=3,
    breaklines=true, % Este comando é para quebrar uma linha comprida
    postbreak=\mbox{\textcolor{red}{$\hookrightarrow$}\space}, % Este comando é para indicar com uma seta vermelha que a linha continua embaixo 
    extendedchars=true,
    literate=
        {á}{{\'a}}1 {Á}{{\'A}}1 {à}{{\`a}}1 {À}{{\`A}}1 {ã}{{\~a}}1 {Ã}{{\~A}}1 {â}{{\^a}}1 {Â}{{\^A}}1
        {ç}{{\c{c}}}1 {Ç}{{\c{C}}}1 {é}{{\'e}}1 {É}{{\'E}}1 {ê}{{\^e}}1 {Ê}{{\^E}}1 {í}{{\'\i}}1 {Í}{{\'I}}1
        {ó}{{\'o}}1 {Ó}{{\'O}}1 {ô}{{\^o}}1 {Ô}{{\^O}}1 {õ}{{\~o}}1 {Õ}{{\~O}}1 {ú}{{\'u}}1 {Ú}{{\'U}}1
        {ü}{{\"u}}1 {Ü}{{\"U}}1 {“}{{``}}1 {”}{{''}}1 {’}{{'}}1 {—}{{--}}1 {–}{{--}}1 {−}{{-}}1
        {±}{{$\pm$}}1 {σ}{{$\sigma$}}1 {→}{{$\to$}}1 {…}{{\ldots}}1
}

\lstset{inputencoding=utf8/latin1}

\lstdefinelanguage{Matlab}{
	language=Matlab,
	basicstyle=\ttfamily\small,
	keywordstyle=\color{blue}\bfseries,
	stringstyle=\color{red},
	commentstyle=\color{gray}\itshape,
	morekeywords={function, end, if, else, for, while, switch, case, otherwise},
	sensitive=true
}
% ----------

% ----------
\usepackage{nomencl} % Pacote para criar nomenclatura e lista de nomenclatura
\makenomenclature % Criar a nomenclatura
\renewcommand{\nomname}{Lista de Nomenclaturas} % Renomear o título da nomenclatura
% ----------

\usepackage{amsmath} % Pacote para inserir matrix com chaves, entre outras coisas

\usepackage{booktabs} % Pacote para usar \toprule, \midrule e \bottomrule


% EXECUTAR CÁLCULOS A PARTIR DA LINHA DE COMANDO: ---------
%\usepackage{siunitx, xfp} % Pacote para executar cálculos a partir da linha de comando
%\usepackage{l3regex}%unless you have very new expl3
%%\sisetup{locale=DE}
%%\let\oldnum\num
%%\ExplSyntaxOn
%\def\num#1{\def\tmp{#1}\regex_replace_all:nnN{,([0-9])}{.\1}{\tmp}\oldnum{\tmp}}
%\ExplSyntaxOff
% ----------


%\pagecolor{black} % Alterar a cor da página
%\color{white} % Alterar a cor do texto

% ----------
% Pacotes para inserir algoritmos e pseudo-códigos:
\setlength{\marginparwidth}{2.5cm}
\usepackage[colorinlistoftodos]{todonotes}
\usepackage{algorithm} % Pacote para inserir algoritmos
\usepackage{algpseudocode} % Pacote para inserir algoritmos e pseudo-códigos
\floatname{algorithm}{Algoritmo} % Renomear a palavra Algorithm para Procedimento
\renewcommand{\algorithmicrequire}{\textbf{entrada:}} % Renomear a palavra Input para Entrada
\renewcommand{\algorithmicensure}{\textbf{saída:}} % Renomear a palavra Output para Saída
\renewcommand{\algorithmicwhile}{\textbf{enquanto}} % Renomear a palavra While para Enquanto
\renewcommand{\algorithmicif}{\textbf{se}} % Renomear a palavra If para Se
\renewcommand{\algorithmicend}{\textbf{fim}} % Renomear a palavra End Output para Fim do
\renewcommand{\algorithmicprocedure}{\textbf{procedimento}} % Renomear a palavra Procedure para Procedimento
\renewcommand{\algorithmicfor}{\textbf{para}} % Renomear a palavra For para Para
\renewcommand{\algorithmicelse}{\textbf{senão}} % Renomear a palavra Else para Senão
\renewcommand{\algorithmicfunction}{\textbf{função}} % Renomear a palavra Function para Função 
\renewcommand{\algorithmicthen}{\textbf{então}} % Renomear a palavra Then para Então
\renewcommand{\algorithmicdo}{\textbf{fazer}} % Renomear a palavra Do para Fazer
% ----------

\newcommand{\quotes}[1]{``#1''} % Comando para inserir um texto dentro de aspas-dupla com espaço após o final

% FLUXOGRAMA: ----------
\usepackage{tikz}
\usetikzlibrary{positioning, shapes.geometric}
\tikzstyle{line}=[draw] % here

\tikzstyle{startstop} = [rectangle, rounded corners, minimum width=3cm,
minimum height=1cm,text centered, draw=black, fill=white!30]
\tikzstyle{arrow} = [thick, ->, >=stealth]
\tikzstyle{process} = [rectangle, minimum width=3cm, minimum height=1cm, text centered, draw=black, fill=white!30]
\tikzstyle{io} = [trapezium, trapezium left angle=70, trapezium right angle=110, minimum width=3cm, minimum height=1cm, text centered, draw=black, fill=white!30]
\tikzstyle{connector} = [draw, circle, fill=white!20]
\tikzstyle{decision} = [diamond, minimum width=3cm, minimum height=1cm, text centered, draw=black, fill=white!30, aspect=3]
% ----------


% CAIXA ENTORNO DAS FÓRMULAS: ----------
\usepackage{dashbox}
\newcommand\dboxed[1]{\dbox{\ensuremath{#1}}}

% ----------

%\usepackage{tabularx} % para colunas com largura ajustável

% -------- CONFIGURAÇÃO DE MARGENS --------
% \usepackage[top=2cm,bottom=2cm,left=3cm,right=2cm]{geometry}

% -------- PACOTE fancyhdr --------
\usepackage{fancyhdr}

\setlength{\marginparwidth}{2cm}

% Ajuste do tamanho do cabeçalho para acompanhar exatamente \textwidth
\setlength{\headwidth}{\textwidth}

% Define um headheight mínimo (para evitar avisos)
\setlength{\headheight}{15pt} 

% Limpa configurações padrão de cabeçalho/rodapé
\pagestyle{fancy}
\fancyhf{} 

% Marcações automáticas de seção para o cabeçalho (opcional)
\renewcommand{\sectionmark}[1]{%
	\markboth{\thesection\ #1}{}%
}

% Texto do cabeçalho à esquerda: seção e título
% Usamos \leftmark pois o \sectionmark acima grava no \leftmark
\fancyhead[L]{\small \leftmark}

% Texto do cabeçalho à direita: número de página
\fancyhead[R]{\small \thepage}

% Linha horizontal no cabeçalho (espessura 0.4pt)
\renewcommand{\headrulewidth}{0.4pt}

% Sobrescreve o título padrão “References” para “Referências”
%\renewcommand{\refname}{Referências}
% Tradução de autoref para o português
\def\figureautorefname{Figura}
\def\tableautorefname{Tabela}
\def\equationautorefname{Equa\c{c}\~ao}

% \hypersetup{backref=true}

\ifdefined\backref
        \renewcommand*{\backref}[1]{}
        \renewcommand*{\backrefentrycount}[2]{#1\ifnum#2>1 ~(#2)\fi}
\fi

\DeclarePairedDelimiter{\ceil}{\lceil}{\rceil}

\onehalfspacing

\usepackage{pgfgantt}

%++++++++++++++++++++++++++++++++++++++++++++++++++++++++++++++++++++++++++++++
% Identificacoes (se o trabalho for em inglês, insira os dados em inglês)
% Para entradas abreviadas de Professora (Profa.) em português escreva: Prof$^\textnormal{a}$.
%++++++++++++++++++++++++++++++++++++++++++++++++++++++++++++++++++++++++++++++
\course{Engenharia Aeronáutica e Mecânica} % Programa de PG ou Curso de Graduação
\area{Propulsão Aeroespacial e Energia} % Área de concentração na PG (Não utilizado no caso de TG)

% Autor do trabalho: Nome Sobrenome
\authorgender{masc} %sexo: masc ou fem
\author{\textbf{EDEN DENIS} Ferreira da Silva Lopes}{Santos}
\itaauthoraddress{Rua João Fonseca dos Santos, 20, Apto. 93}{12.230-088}{São José dos Campos--SP}

% Titulo da Tese/Dissertação
\title{Aplicação de Inteligência Artificial (AI) em Banco de Dados de Emissão de Motores}

% Orientador
\advisorgender{masc} % masc ou fem
\advisor{Prof.~Dr.}{Pedro Teixeira \textbf{LACAVA}}{ITA}

% Coorientador (Caso não haja coorientador, colocar ambas as variáveis \coadvisorgender e \coadvisor comentadas, com um % na frente)
\coadvisorgender{masc}
\coadvisor{}{}  % {Prof.~Dr.}{Carlos Alberto Rocha \textbf{PIMENTEL}}{ITA}

% Pró-reitor da Pós-graduação
\bossgender{fem}
\boss{Prof.~Dra.}{Emília \textbf{VILLANI}}

%Coordenador do curso no caso de TG
\bosscoursegender{masc}
\bosscourse{Prof.~Dr.}{ - } % masc ou fem

% Palavras-Chaves informadas pela Biblioteca -> utilizada na CIP
\kwcip{Engine taxonomy}
\kwcip{Emission classification}
\kwcip{Inteligência Artificial}
\kwcip{Artificial Intelligence}
\kwcip{Python}

% membros da banca examinadora
\examiner{Prof. Dr.}{Pedro Teixeira \textbf{LACAVA}}{President}{ITA}
\examiner{Prof. Dr.}{}{Internal Member}{ITA} % Vinícius \textbf{MALATESTA}}{Internal Member}{ITA}
\examiner{Prof. Dr.}{}{Internal Member}{ITA} % {Elisan dos Santos \textbf{MAGALHÃES}}{Internal Member}{ITA}
\examiner{Prof. Dr.}{}{Internal Member}{ITA} % {Carlos Alberto Rocha \textbf{PIMENTEL}}{External Member}{UFABC}
\examiner{Prof. Dr.}{}{Internal Member}{ITA} % {\textbf{KARL} Peter Burr}{External Member}{UFABC}

% Data da defesa (mês em maiúsculo, se trabalho em inglês, e minúsculo se trabalho em português) 
\date{14}{JULHO}{2025}

% Número CDU - (somente para TG)
\cdu{621.38}

% Folha de Registro do Documento
\newcommand{\advisorTitle}{Prof.~Dr.}
\newcommand{\advisorName}{Pedro Teixeira \textbf{LACAVA}}

\newcommand{\coadvisorTitle}{Prof.~Dr.}
\newcommand{\coadvisorName}{Carlos Alberto Rocha \textbf{PIMENTEL}}

\newcommand{\dateDay}{14}
\newcommand{\dateMonth}{JULHO}
\newcommand{\dateYear}{2025}

% Data de publicação e defesa (usada na folha de registro)
\newcommand{\frdDataPublicacao}{04 de agosto de 2025}
\newcommand{\frdNumeroRegistro}{DCTA/ITA/TD-029/2020}
\newcommand{\frdPalavrasAutor}{Big Data; Engine Emission; Emission Databank; Clustering; Artificial Intelligence}
\newcommand{\frdPalavrasResultado}{Taxonomia; Emissão de Motores; Planejamento de Dados; Inteligência Artificial; Processamento de Dados}
\newcommand{\frdApresentacao}{
	ITA, São José dos Campos. Curso de Mestrado. Programa de Pós-Graduação em Engenharia Aeronáutica e Mecânica. Área de Propulsão Aeroespacial e Energia.\\
	Orientador: \advisorTitle~\advisorName.\\
	Coorientador: \coadvisorTitle~\coadvisorName.\\
	Defesa em \dateDay/\dateMonth/\dateYear. Publicada em \frdDataPublicacao.
}

\newcommand{\PorcentagemDaLargura}{0.80}

\frontmatter
\makenoidxglossaries
\newacronym{acr:ai}{AI}{Artificial Intelligence (Inteligencia Artificial)}
\newacronym{acr:co}{CO}{Monoxido de carbono (Carbon Monoxide)}
\newacronym{acr:eedb}{EEDB}{Engine Emissions Databank}
\newacronym{acr:ei}{EI}{Emission Index}
\newacronym{acr:hc}{HC}{Hidrocarboneto (Hydrocarbon)}
\newacronym{acr:icao}{ICAO}{International Civil Aviation Organization}
\newacronym{acr:ita}{ITA}{Instituto Tecnologico de Aeronautica}
\newacronym{acr:lto}{LTO}{Landing and Take-Off}
\newacronym{acr:nox}{NOx}{Oxidos de nitrogenio (Nitrogen Oxides)}

\newglossaryentry{sym:adjacency_matrix}{
    name={$\mathcal{A} \in \{0,1\}^{N_{ob} \times N_{ob}}$},
    description={matriz de adjacencias entre os obstaculos do territorio},
    sort={a_adj}
}
\newglossaryentry{sym:connectivity_matrix}{
    name={$\mathcal{C} \in \{0,1\}^{N_{ob} \times N_{ob}}$},
    description={matriz de conectividade utilizada para verificar continuidade entre obstaculos},
    sort={c_conn}
}
\newglossaryentry{sym:cluster_centroid}{
    name={$C_{\ell} \in \mathbb{R}^{2}$},
    description={coordenadas do centroide (cluster $\ell$) obtido pelo particionamento},
    sort={c_cent}
}
\newglossaryentry{sym:distance_to_medoid}{
    name={$d_{i,\bar{m}_{k}}$},
    description={distancia do elemento $i$ ao medioide $\bar{m}_{k}$ do respectivo cluster},
    sort={d_med}
}
\newglossaryentry{sym:agent_obstacle_vector}{
    name={$\Delta r_{ao} \in \mathbb{R}^{N_{ob}}$},
    description={vetor de distancias agente-obstaculo},
    sort={delta_ao}
}
\newglossaryentry{sym:obstacle_distance_matrix}{
    name={$\Delta r_{o} \in \mathbb{R}^{N_{ob} \times N_{ob}}$},
    description={matriz de distancias entre obstaculos},
    sort={delta_o}
}
\newglossaryentry{sym:gamma_weight}{
    name={$\gamma \in \mathbb{R}_{+}$},
    description={peso associado ao termo de consumo de combustivel na funcao de custo},
    sort={gamma}
}
\newglossaryentry{sym:graph_obstacles}{
    name={$\mathcal{G} = (\mathcal{O}, \Omega)$},
    description={grafo nao direcionado que descreve as adjacencias do conjunto de obstaculos $\mathcal{O}$},
    sort={g_graph}
}
\newglossaryentry{sym:predicted_generic_variable}{
    name={$\hat{\odot}[k+\alpha \mid k]$},
    description={valor previsto da variavel generica $\odot$ para o instante $k+\alpha$ usando dados ate $k$},
    sort={hat_odo}
}
\newglossaryentry{sym:k_index}{
    name={$k \in \mathbb{N}$},
    description={indice da iteracao ou instante corrente},
    sort={k_index}
}
\newglossaryentry{sym:k_clusters}{
    name={$K \in \mathbb{N}$},
    description={número de clusters considerados no particionamento do K-means},
    sort={k_clusters}
}
\newglossaryentry{sym:n_observations}{
    name={$n \in \mathbb{N}$},
    description={quantidade total de observacoes disponiveis para agrupamento},
    sort={n_obs}
}
\newglossaryentry{sym:nc_clusters}{
    name={$N_c \in \mathbb{N}$},
    description={numero de clusters (medioides) selecionados apos a convergencia},
    sort={nc}
}
\newglossaryentry{sym:nob_obstacles}{
    name={$N_{ob} \in \mathbb{N}$},
    description={quantidade de obstaculos no territorio analisado},
    sort={nob}
}
\newglossaryentry{sym:terminal_set}{
    name={$\mathcal{Q} \subset \mathbb{R}^{4}$},
    description={conjunto admissivel de estados terminais da planta},
    sort={q_set}
}
\newglossaryentry{sym:agent_position}{
    name={$r = (r_x, r_y) \in \mathbb{R}^{2}$},
    description={posicao do agente no plano horizontal},
    sort={r_pos}
}
\newglossaryentry{sym:territory_region}{
    name={$\mathcal{T} \subset \mathbb{R}^{2}$},
    description={regiao do plano onde estao definidos os obstaculos e as trajetorias},
    sort={t_region}
}
\newglossaryentry{sym:sampling_period}{
    name={$T \in \mathbb{R}_{+}$},
    description={periodo de amostragem do controlador},
    sort={t_sampling}
}
\newglossaryentry{sym:control_input}{
    name={$u \in \mathcal{U}$},
    description={vetor de entrada de controle dentro do conjunto admissivel $\mathcal{U}$},
    sort={u_input}
}
\newglossaryentry{sym:control_set}{
    name={$\mathcal{U} \subset \mathbb{R}^{2}$},
    description={politopo das entradas de controle que contem a origem},
    sort={u_set}
}
\newglossaryentry{sym:vx_vy}{
    name={$v_x$, $v_y$},
    description={componentes da velocidade do agente nos eixos $r_x$ e $r_y$},
    sort={v_xy}
}
\newglossaryentry{sym:obstacle_width}{
    name={$w \in \mathbb{R}_{+}$},
    description={largura media de um obstaculo considerado na modelagem},
    sort={w_width}
}
\newglossaryentry{sym:plant_state}{
    name={$x \in \mathcal{X}$},
    description={estado instantaneo da planta controlada},
    sort={x_state}
}
\newglossaryentry{sym:sample_xi}{
    name={$x_i$},
    description={i-ésimo elemento (amostra) utilizado pelo algoritmo K-means},
    sort={xi}
}
\newglossaryentry{sym:state_set}{
    name={$\mathcal{X} \subset \mathbb{R}^{4}$},
    description={politopo dos estados admissiveis},
    sort={x_set}
}
\newglossaryentry{sym:medoid_mk}{
    name={$\bar{m}_{k} \in \mathbb{R}^{d}$},
    description={medioide do cluster $k$ retornado pelo algoritmo},
    sort={mbar}
}
\newglossaryentry{sym:within_sum}{
    name={$WS$},
    description={soma dos quadrados das distâncias intracluster (WSS) referente ao cluster analisado},
    sort={ws}
}
\newglossaryentry{sym:within_sum_total}{
    name={$WS_{Total}$},
    description={soma intra-cluster global usada na analise do Elbow Data},
    sort={ws_total}
}


%\hyphenation{param-e-ter-ization}
 
\begin{document}
	
	\thispagestyle{empty}
	\documentclass[12pt]{article} % Tamanho de letra 12, papel A4, Estilo de artigo

\usepackage[utf8]{inputenc} % Pacote para permitir o uso de acentos
% Ajustes de hifenização
\tolerance=1
\emergencystretch=\maxdimen
\hyphenpenalty=10000
\hbadness=10000
\usepackage{lipsum} % Pacote para gerar texto aleatório
\usepackage{ragged2e} % Pacote para mais opções de alinhamento
\usepackage{graphicx} % Pacote para inserir figuras jpeg, pdf e png
\usepackage{ae}
\usepackage{epsfig}
\usepackage{mathtools}
\usepackage[export]{adjustbox}
\usepackage{lineno}
\usepackage{amsthm}
\usepackage{commath}
\usepackage{color, colortbl}
\usepackage{bbm}
\usepackage{setspace} % Pacote para ajustar espaçamento entre linhas

% Instructions to prevent LaTeX from hyphenating words
\tolerance=1
\emergencystretch=\maxdimen
\hyphenpenalty=10000
\hbadness=10000

\graphicspath{{figures/}} % Indica o diretório no qual as imagens estão sendo salvas, evitando que precisamos digitar parte do caminho da pasta que contém as imagens
\usepackage[rightcaption]{sidecap} % Pacote para mais opções de figura, colocar, título da legenda da figura ao lado direito
\usepackage{url} % Pacote para citar referências de sites
\usepackage{makeidx} % Pacote para criar índice
\makeindex % Habilitar a criação do índice
\usepackage{longtable} % Pacote para usar tabelas longas

% ----------
% Pacotes para dar quebra de linha dentro da tabela:
%\usepackage{fourier} 
%\usepackage{array}
\usepackage{makecell}

\renewcommand\theadalign{bc}
\renewcommand\theadfont{\bfseries}
\renewcommand\theadgape{\Gape[4pt]}
\renewcommand\cellgape{\Gape[4pt]}
% ----------

\usepackage{wrapfig} % Pacote para alinhar figuras à esquerda ou à direita do texto
\usepackage{ae}
\usepackage{epsfig}
\usepackage{mathtools}
\usepackage[export]{adjustbox}
\usepackage{lineno}
\usepackage{amsfonts}
\usepackage{amsthm}
\usepackage{commath}
\usepackage{color, colortbl}
\usepackage{bbm}

% ----------
% Pacote para inserir cabeçalho e rodapé:
% Atenção! Se remover o cabeçalho e rodapé, o LaTeX coloca como cabeçalho o padrão dele, que é bem interessante também!
% \usepackage{fancyhdr} % Pacote para inserir cabeçalho e rodapé
% \pagestyle{fancy} % Estilo do cabeçalho e rodapé, neste caso apenas uma linha no cabeçalho
%\lhead{{\footnotesize Prandiano\textregistered}} % Inserir texto no canto esquerdo do cabeçalho
%\chead{centro} % Inserir texto no centro do cabeçalho
%\rhead{{\footnotesize Curso de Python}} % Inserir texto no canto direito do cabeçalho
%\lfoot{{\footnotesize Copyright \textcopyright 2018 Prandiano}} % Inserir texto no canto esquerdo do rodapé
%\cfoot{centro} % Inserir texto no centro do rodapé
%\rfoot{{\footnotesize Reprodução proibida: Lei 9.610 de 19/02/1998}} % Inserir texto no canto direito do rodapé
% O comando \renewcommand re-nova o parâmetro especificado
%\renewcommand{\footrulewidth}{0.4pt} % Inserir uma linha acima do rodapé
% ----------

% ----------
% Pacote para inserir múltiplas figuras:
\usepackage{subfigure} % Pacote para inserir múltiplas figuras
% \usepackage{subcaption} % Pacote incompatível com subfigure
%\usepackage{caption}
% ----------

\usepackage{multicol} % Pacote que permite usar multicolunas
\usepackage{float} % Pacote para usar o "H" do figure
\usepackage[dvipsnames]{xcolor} % Pacote de cores para textos e possivelmente outros
\usepackage{amsmath} % Pacote para usar o cases
\usepackage[makeroom]{cancel} % Pacote para usar o cancelar (em e embaixo)
\usepackage{enumitem} % Pacote para enumerar itens
\usepackage{amssymb} % Pacote para inserir os símbolos dos conjuntos numéricos N, Z, Q, R, C
\usepackage[T1]{fontenc}
\usepackage{enumitem}
\usepackage{textcomp}
\usepackage{pdfpages} % Pacote para inserir página em PDF
\usepackage{courier} % Pacote para usar a fonte courier
%\usepackage{matlab-prettifier}
\usepackage{lscape} % Pacote para mudar sentido da página
\newtheorem{theorem}{Teorema} % Pacote para inserir Teorema
\newtheorem{corollary}{Corollary} % Pacote para inserir Corolário
\newtheorem{remark}{Remark}
\newtheorem{mytheorem}{Definition}
\usepackage{hyperref} % Pacote para inserir hyperref
\usepackage{multirow} % Pacote para usar os comandos \multirow
\usepackage{pdfpages} % Pacote para incluir PDF's
\usepackage{listingsutf8} % Pacote para inserir arquivos externos tipo .txt
\usepackage{etoolbox}
\AtBeginEnvironment{verbatim}{\footnotesize}

\usepackage{standalone} % Pacote para inserir sub-arquivos LaTeX
\usepackage{bookmark} % Pacote para NÃO tratar o Índice como Apêndice e colocar letra

% Nome da listagem em português
\renewcommand{\lstlistingname}{Listagem}

% REFERÊNCIAS: ----------
%\usepackage[numbers, compress]{natbib} % Pacote para usar as referências bibliográficas
%\usepackage{usebib}
% \bibliographystyle{references} % Estilo das referências
%\setcitestyle{authoryear,aysep={},open={(},close={)}}
%\newcommand{\myand}{e}

% ----------


% PARÂMETROS DA FUNÇÃO LISTINGS: ----------
\lstset
{
    basicstyle=\small\ttfamily,
    commentstyle=\color{Green},
    keywordstyle=\color{Cerulean},
    frame=single,
    language=Python,
    morekeywords={True, False},
    numbers=left,
    numbersep=10pt,
    numberstyle=\footnotesize\color{Gray},
    showstringspaces=false,
    stringstyle=\color{Mulberry},
    tabsize=3,
    breaklines=true, % Este comando é para quebrar uma linha comprida
    postbreak=\mbox{\textcolor{red}{$\hookrightarrow$}\space}, % Este comando é para indicar com uma seta vermelha que a linha continua embaixo 
    extendedchars=true,
    literate=
        {á}{{\'a}}1 {Á}{{\'A}}1 {à}{{\`a}}1 {À}{{\`A}}1 {ã}{{\~a}}1 {Ã}{{\~A}}1 {â}{{\^a}}1 {Â}{{\^A}}1
        {ç}{{\c{c}}}1 {Ç}{{\c{C}}}1 {é}{{\'e}}1 {É}{{\'E}}1 {ê}{{\^e}}1 {Ê}{{\^E}}1 {í}{{\'\i}}1 {Í}{{\'I}}1
        {ó}{{\'o}}1 {Ó}{{\'O}}1 {ô}{{\^o}}1 {Ô}{{\^O}}1 {õ}{{\~o}}1 {Õ}{{\~O}}1 {ú}{{\'u}}1 {Ú}{{\'U}}1
        {ü}{{\"u}}1 {Ü}{{\"U}}1 {“}{{``}}1 {”}{{''}}1 {’}{{'}}1 {—}{{--}}1 {–}{{--}}1 {−}{{-}}1
        {±}{{$\pm$}}1 {σ}{{$\sigma$}}1 {→}{{$\to$}}1 {…}{{\ldots}}1
}

\lstset{inputencoding=utf8/latin1}

\lstdefinelanguage{Matlab}{
	language=Matlab,
	basicstyle=\ttfamily\small,
	keywordstyle=\color{blue}\bfseries,
	stringstyle=\color{red},
	commentstyle=\color{gray}\itshape,
	morekeywords={function, end, if, else, for, while, switch, case, otherwise},
	sensitive=true
}
% ----------

% ----------
\usepackage{nomencl} % Pacote para criar nomenclatura e lista de nomenclatura
\makenomenclature % Criar a nomenclatura
\renewcommand{\nomname}{Lista de Nomenclaturas} % Renomear o título da nomenclatura
% ----------

\usepackage{amsmath} % Pacote para inserir matrix com chaves, entre outras coisas

\usepackage{booktabs} % Pacote para usar \toprule, \midrule e \bottomrule


% EXECUTAR CÁLCULOS A PARTIR DA LINHA DE COMANDO: ---------
%\usepackage{siunitx, xfp} % Pacote para executar cálculos a partir da linha de comando
%\usepackage{l3regex}%unless you have very new expl3
%%\sisetup{locale=DE}
%%\let\oldnum\num
%%\ExplSyntaxOn
%\def\num#1{\def\tmp{#1}\regex_replace_all:nnN{,([0-9])}{.\1}{\tmp}\oldnum{\tmp}}
%\ExplSyntaxOff
% ----------


%\pagecolor{black} % Alterar a cor da página
%\color{white} % Alterar a cor do texto

% ----------
% Pacotes para inserir algoritmos e pseudo-códigos:
\setlength{\marginparwidth}{2.5cm}
\usepackage[colorinlistoftodos]{todonotes}
\usepackage{algorithm} % Pacote para inserir algoritmos
\usepackage{algpseudocode} % Pacote para inserir algoritmos e pseudo-códigos
\floatname{algorithm}{Algoritmo} % Renomear a palavra Algorithm para Procedimento
\renewcommand{\algorithmicrequire}{\textbf{entrada:}} % Renomear a palavra Input para Entrada
\renewcommand{\algorithmicensure}{\textbf{saída:}} % Renomear a palavra Output para Saída
\renewcommand{\algorithmicwhile}{\textbf{enquanto}} % Renomear a palavra While para Enquanto
\renewcommand{\algorithmicif}{\textbf{se}} % Renomear a palavra If para Se
\renewcommand{\algorithmicend}{\textbf{fim}} % Renomear a palavra End Output para Fim do
\renewcommand{\algorithmicprocedure}{\textbf{procedimento}} % Renomear a palavra Procedure para Procedimento
\renewcommand{\algorithmicfor}{\textbf{para}} % Renomear a palavra For para Para
\renewcommand{\algorithmicelse}{\textbf{senão}} % Renomear a palavra Else para Senão
\renewcommand{\algorithmicfunction}{\textbf{função}} % Renomear a palavra Function para Função 
\renewcommand{\algorithmicthen}{\textbf{então}} % Renomear a palavra Then para Então
\renewcommand{\algorithmicdo}{\textbf{fazer}} % Renomear a palavra Do para Fazer
% ----------

\newcommand{\quotes}[1]{``#1''} % Comando para inserir um texto dentro de aspas-dupla com espaço após o final

% FLUXOGRAMA: ----------
\usepackage{tikz}
\usetikzlibrary{positioning, shapes.geometric}
\tikzstyle{line}=[draw] % here

\tikzstyle{startstop} = [rectangle, rounded corners, minimum width=3cm,
minimum height=1cm,text centered, draw=black, fill=white!30]
\tikzstyle{arrow} = [thick, ->, >=stealth]
\tikzstyle{process} = [rectangle, minimum width=3cm, minimum height=1cm, text centered, draw=black, fill=white!30]
\tikzstyle{io} = [trapezium, trapezium left angle=70, trapezium right angle=110, minimum width=3cm, minimum height=1cm, text centered, draw=black, fill=white!30]
\tikzstyle{connector} = [draw, circle, fill=white!20]
\tikzstyle{decision} = [diamond, minimum width=3cm, minimum height=1cm, text centered, draw=black, fill=white!30, aspect=3]
% ----------


% CAIXA ENTORNO DAS FÓRMULAS: ----------
\usepackage{dashbox}
\newcommand\dboxed[1]{\dbox{\ensuremath{#1}}}

% ----------

%\usepackage{tabularx} % para colunas com largura ajustável

% -------- CONFIGURAÇÃO DE MARGENS --------
% \usepackage[top=2cm,bottom=2cm,left=3cm,right=2cm]{geometry}

% -------- PACOTE fancyhdr --------
\usepackage{fancyhdr}

\setlength{\marginparwidth}{2cm}

% Ajuste do tamanho do cabeçalho para acompanhar exatamente \textwidth
\setlength{\headwidth}{\textwidth}

% Define um headheight mínimo (para evitar avisos)
\setlength{\headheight}{15pt} 

% Limpa configurações padrão de cabeçalho/rodapé
\pagestyle{fancy}
\fancyhf{} 

% Marcações automáticas de seção para o cabeçalho (opcional)
\renewcommand{\sectionmark}[1]{%
	\markboth{\thesection\ #1}{}%
}

% Texto do cabeçalho à esquerda: seção e título
% Usamos \leftmark pois o \sectionmark acima grava no \leftmark
\fancyhead[L]{\small \leftmark}

% Texto do cabeçalho à direita: número de página
\fancyhead[R]{\small \thepage}

% Linha horizontal no cabeçalho (espessura 0.4pt)
\renewcommand{\headrulewidth}{0.4pt}

% Sobrescreve o título padrão “References” para “Referências”
%\renewcommand{\refname}{Referências}
% Tradução de autoref para o português
\def\figureautorefname{Figura}
\def\tableautorefname{Tabela}
\def\equationautorefname{Equa\c{c}\~ao}

% \hypersetup{backref=true}

\ifdefined\backref
        \renewcommand*{\backref}[1]{}
        \renewcommand*{\backrefentrycount}[2]{#1\ifnum#2>1 ~(#2)\fi}
\fi

\DeclarePairedDelimiter{\ceil}{\lceil}{\rceil}

\onehalfspacing

\usepackage{pgfgantt}

%++++++++++++++++++++++++++++++++++++++++++++++++++++++++++++++++++++++++++++++
% Identificacoes (se o trabalho for em inglês, insira os dados em inglês)
% Para entradas abreviadas de Professora (Profa.) em português escreva: Prof$^\textnormal{a}$.
%++++++++++++++++++++++++++++++++++++++++++++++++++++++++++++++++++++++++++++++
\course{Engenharia Aeronáutica e Mecânica} % Programa de PG ou Curso de Graduação
\area{Propulsão Aeroespacial e Energia} % Área de concentração na PG (Não utilizado no caso de TG)

% Autor do trabalho: Nome Sobrenome
\authorgender{masc} %sexo: masc ou fem
\author{\textbf{EDEN DENIS} Ferreira da Silva Lopes}{Santos}
\itaauthoraddress{Rua João Fonseca dos Santos, 20, Apto. 93}{12.230-088}{São José dos Campos--SP}

% Titulo da Tese/Dissertação
\title{Aplicação de Inteligência Artificial (AI) em Banco de Dados de Emissão de Motores}

% Orientador
\advisorgender{masc} % masc ou fem
\advisor{Prof.~Dr.}{Pedro Teixeira \textbf{LACAVA}}{ITA}

% Coorientador (Caso não haja coorientador, colocar ambas as variáveis \coadvisorgender e \coadvisor comentadas, com um % na frente)
\coadvisorgender{masc}
\coadvisor{}{}  % {Prof.~Dr.}{Carlos Alberto Rocha \textbf{PIMENTEL}}{ITA}

% Pró-reitor da Pós-graduação
\bossgender{fem}
\boss{Prof.~Dra.}{Emília \textbf{VILLANI}}

%Coordenador do curso no caso de TG
\bosscoursegender{masc}
\bosscourse{Prof.~Dr.}{ - } % masc ou fem

% Palavras-Chaves informadas pela Biblioteca -> utilizada na CIP
\kwcip{Engine taxonomy}
\kwcip{Emission classification}
\kwcip{Inteligência Artificial}
\kwcip{Artificial Intelligence}
\kwcip{Python}

% membros da banca examinadora
\examiner{Prof. Dr.}{Pedro Teixeira \textbf{LACAVA}}{President}{ITA}
\examiner{Prof. Dr.}{}{Internal Member}{ITA} % Vinícius \textbf{MALATESTA}}{Internal Member}{ITA}
\examiner{Prof. Dr.}{}{Internal Member}{ITA} % {Elisan dos Santos \textbf{MAGALHÃES}}{Internal Member}{ITA}
\examiner{Prof. Dr.}{}{Internal Member}{ITA} % {Carlos Alberto Rocha \textbf{PIMENTEL}}{External Member}{UFABC}
\examiner{Prof. Dr.}{}{Internal Member}{ITA} % {\textbf{KARL} Peter Burr}{External Member}{UFABC}

% Data da defesa (mês em maiúsculo, se trabalho em inglês, e minúsculo se trabalho em português) 
\date{14}{JULHO}{2025}

% Número CDU - (somente para TG)
\cdu{621.38}

% Folha de Registro do Documento
\newcommand{\advisorTitle}{Prof.~Dr.}
\newcommand{\advisorName}{Pedro Teixeira \textbf{LACAVA}}

\newcommand{\coadvisorTitle}{Prof.~Dr.}
\newcommand{\coadvisorName}{Carlos Alberto Rocha \textbf{PIMENTEL}}

\newcommand{\dateDay}{14}
\newcommand{\dateMonth}{JULHO}
\newcommand{\dateYear}{2025}

% Data de publicação e defesa (usada na folha de registro)
\newcommand{\frdDataPublicacao}{04 de agosto de 2025}
\newcommand{\frdNumeroRegistro}{DCTA/ITA/TD-029/2020}
\newcommand{\frdPalavrasAutor}{Big Data; Engine Emission; Emission Databank; Clustering; Artificial Intelligence}
\newcommand{\frdPalavrasResultado}{Taxonomia; Emissão de Motores; Planejamento de Dados; Inteligência Artificial; Processamento de Dados}
\newcommand{\frdApresentacao}{
	ITA, São José dos Campos. Curso de Mestrado. Programa de Pós-Graduação em Engenharia Aeronáutica e Mecânica. Área de Propulsão Aeroespacial e Energia.\\
	Orientador: \advisorTitle~\advisorName.\\
	Coorientador: \coadvisorTitle~\coadvisorName.\\
	Defesa em \dateDay/\dateMonth/\dateYear. Publicada em \frdDataPublicacao.
}

\newcommand{\PorcentagemDaLargura}{0.80}

\begin{document}
	
\end{document} \clearpage  % [x] Revisão do texto
	
	% Folha de Rosto e Capa para o caso do TG
	\maketitle
	
	% Dedicatoria:
	\documentclass[12pt]{article} % Tamanho de letra 12, papel A4, Estilo de artigo

\usepackage[utf8]{inputenc} % Pacote para permitir o uso de acentos
% Ajustes de hifenização
\tolerance=1
\emergencystretch=\maxdimen
\hyphenpenalty=10000
\hbadness=10000
\usepackage{lipsum} % Pacote para gerar texto aleatório
\usepackage{ragged2e} % Pacote para mais opções de alinhamento
\usepackage{graphicx} % Pacote para inserir figuras jpeg, pdf e png
\usepackage{ae}
\usepackage{epsfig}
\usepackage{mathtools}
\usepackage[export]{adjustbox}
\usepackage{lineno}
\usepackage{amsthm}
\usepackage{commath}
\usepackage{color, colortbl}
\usepackage{bbm}
\usepackage{setspace} % Pacote para ajustar espaçamento entre linhas

% Instructions to prevent LaTeX from hyphenating words
\tolerance=1
\emergencystretch=\maxdimen
\hyphenpenalty=10000
\hbadness=10000

\graphicspath{{figures/}} % Indica o diretório no qual as imagens estão sendo salvas, evitando que precisamos digitar parte do caminho da pasta que contém as imagens
\usepackage[rightcaption]{sidecap} % Pacote para mais opções de figura, colocar, título da legenda da figura ao lado direito
\usepackage{url} % Pacote para citar referências de sites
\usepackage{makeidx} % Pacote para criar índice
\makeindex % Habilitar a criação do índice
\usepackage{longtable} % Pacote para usar tabelas longas

% ----------
% Pacotes para dar quebra de linha dentro da tabela:
%\usepackage{fourier} 
%\usepackage{array}
\usepackage{makecell}

\renewcommand\theadalign{bc}
\renewcommand\theadfont{\bfseries}
\renewcommand\theadgape{\Gape[4pt]}
\renewcommand\cellgape{\Gape[4pt]}
% ----------

\usepackage{wrapfig} % Pacote para alinhar figuras à esquerda ou à direita do texto
\usepackage{ae}
\usepackage{epsfig}
\usepackage{mathtools}
\usepackage[export]{adjustbox}
\usepackage{lineno}
\usepackage{amsfonts}
\usepackage{amsthm}
\usepackage{commath}
\usepackage{color, colortbl}
\usepackage{bbm}

% ----------
% Pacote para inserir cabeçalho e rodapé:
% Atenção! Se remover o cabeçalho e rodapé, o LaTeX coloca como cabeçalho o padrão dele, que é bem interessante também!
% \usepackage{fancyhdr} % Pacote para inserir cabeçalho e rodapé
% \pagestyle{fancy} % Estilo do cabeçalho e rodapé, neste caso apenas uma linha no cabeçalho
%\lhead{{\footnotesize Prandiano\textregistered}} % Inserir texto no canto esquerdo do cabeçalho
%\chead{centro} % Inserir texto no centro do cabeçalho
%\rhead{{\footnotesize Curso de Python}} % Inserir texto no canto direito do cabeçalho
%\lfoot{{\footnotesize Copyright \textcopyright 2018 Prandiano}} % Inserir texto no canto esquerdo do rodapé
%\cfoot{centro} % Inserir texto no centro do rodapé
%\rfoot{{\footnotesize Reprodução proibida: Lei 9.610 de 19/02/1998}} % Inserir texto no canto direito do rodapé
% O comando \renewcommand re-nova o parâmetro especificado
%\renewcommand{\footrulewidth}{0.4pt} % Inserir uma linha acima do rodapé
% ----------

% ----------
% Pacote para inserir múltiplas figuras:
\usepackage{subfigure} % Pacote para inserir múltiplas figuras
% \usepackage{subcaption} % Pacote incompatível com subfigure
%\usepackage{caption}
% ----------

\usepackage{multicol} % Pacote que permite usar multicolunas
\usepackage{float} % Pacote para usar o "H" do figure
\usepackage[dvipsnames]{xcolor} % Pacote de cores para textos e possivelmente outros
\usepackage{amsmath} % Pacote para usar o cases
\usepackage[makeroom]{cancel} % Pacote para usar o cancelar (em e embaixo)
\usepackage{enumitem} % Pacote para enumerar itens
\usepackage{amssymb} % Pacote para inserir os símbolos dos conjuntos numéricos N, Z, Q, R, C
\usepackage[T1]{fontenc}
\usepackage{enumitem}
\usepackage{textcomp}
\usepackage{pdfpages} % Pacote para inserir página em PDF
\usepackage{courier} % Pacote para usar a fonte courier
%\usepackage{matlab-prettifier}
\usepackage{lscape} % Pacote para mudar sentido da página
\newtheorem{theorem}{Teorema} % Pacote para inserir Teorema
\newtheorem{corollary}{Corollary} % Pacote para inserir Corolário
\newtheorem{remark}{Remark}
\newtheorem{mytheorem}{Definition}
\usepackage{hyperref} % Pacote para inserir hyperref
\usepackage{multirow} % Pacote para usar os comandos \multirow
\usepackage{pdfpages} % Pacote para incluir PDF's
\usepackage{listingsutf8} % Pacote para inserir arquivos externos tipo .txt
\usepackage{etoolbox}
\AtBeginEnvironment{verbatim}{\footnotesize}

\usepackage{standalone} % Pacote para inserir sub-arquivos LaTeX
\usepackage{bookmark} % Pacote para NÃO tratar o Índice como Apêndice e colocar letra

% Nome da listagem em português
\renewcommand{\lstlistingname}{Listagem}

% REFERÊNCIAS: ----------
%\usepackage[numbers, compress]{natbib} % Pacote para usar as referências bibliográficas
%\usepackage{usebib}
% \bibliographystyle{references} % Estilo das referências
%\setcitestyle{authoryear,aysep={},open={(},close={)}}
%\newcommand{\myand}{e}

% ----------


% PARÂMETROS DA FUNÇÃO LISTINGS: ----------
\lstset
{
    basicstyle=\small\ttfamily,
    commentstyle=\color{Green},
    keywordstyle=\color{Cerulean},
    frame=single,
    language=Python,
    morekeywords={True, False},
    numbers=left,
    numbersep=10pt,
    numberstyle=\footnotesize\color{Gray},
    showstringspaces=false,
    stringstyle=\color{Mulberry},
    tabsize=3,
    breaklines=true, % Este comando é para quebrar uma linha comprida
    postbreak=\mbox{\textcolor{red}{$\hookrightarrow$}\space}, % Este comando é para indicar com uma seta vermelha que a linha continua embaixo 
    extendedchars=true,
    literate=
        {á}{{\'a}}1 {Á}{{\'A}}1 {à}{{\`a}}1 {À}{{\`A}}1 {ã}{{\~a}}1 {Ã}{{\~A}}1 {â}{{\^a}}1 {Â}{{\^A}}1
        {ç}{{\c{c}}}1 {Ç}{{\c{C}}}1 {é}{{\'e}}1 {É}{{\'E}}1 {ê}{{\^e}}1 {Ê}{{\^E}}1 {í}{{\'\i}}1 {Í}{{\'I}}1
        {ó}{{\'o}}1 {Ó}{{\'O}}1 {ô}{{\^o}}1 {Ô}{{\^O}}1 {õ}{{\~o}}1 {Õ}{{\~O}}1 {ú}{{\'u}}1 {Ú}{{\'U}}1
        {ü}{{\"u}}1 {Ü}{{\"U}}1 {“}{{``}}1 {”}{{''}}1 {’}{{'}}1 {—}{{--}}1 {–}{{--}}1 {−}{{-}}1
        {±}{{$\pm$}}1 {σ}{{$\sigma$}}1 {→}{{$\to$}}1 {…}{{\ldots}}1
}

\lstset{inputencoding=utf8/latin1}

\lstdefinelanguage{Matlab}{
	language=Matlab,
	basicstyle=\ttfamily\small,
	keywordstyle=\color{blue}\bfseries,
	stringstyle=\color{red},
	commentstyle=\color{gray}\itshape,
	morekeywords={function, end, if, else, for, while, switch, case, otherwise},
	sensitive=true
}
% ----------

% ----------
\usepackage{nomencl} % Pacote para criar nomenclatura e lista de nomenclatura
\makenomenclature % Criar a nomenclatura
\renewcommand{\nomname}{Lista de Nomenclaturas} % Renomear o título da nomenclatura
% ----------

\usepackage{amsmath} % Pacote para inserir matrix com chaves, entre outras coisas

\usepackage{booktabs} % Pacote para usar \toprule, \midrule e \bottomrule


% EXECUTAR CÁLCULOS A PARTIR DA LINHA DE COMANDO: ---------
%\usepackage{siunitx, xfp} % Pacote para executar cálculos a partir da linha de comando
%\usepackage{l3regex}%unless you have very new expl3
%%\sisetup{locale=DE}
%%\let\oldnum\num
%%\ExplSyntaxOn
%\def\num#1{\def\tmp{#1}\regex_replace_all:nnN{,([0-9])}{.\1}{\tmp}\oldnum{\tmp}}
%\ExplSyntaxOff
% ----------


%\pagecolor{black} % Alterar a cor da página
%\color{white} % Alterar a cor do texto

% ----------
% Pacotes para inserir algoritmos e pseudo-códigos:
\setlength{\marginparwidth}{2.5cm}
\usepackage[colorinlistoftodos]{todonotes}
\usepackage{algorithm} % Pacote para inserir algoritmos
\usepackage{algpseudocode} % Pacote para inserir algoritmos e pseudo-códigos
\floatname{algorithm}{Algoritmo} % Renomear a palavra Algorithm para Procedimento
\renewcommand{\algorithmicrequire}{\textbf{entrada:}} % Renomear a palavra Input para Entrada
\renewcommand{\algorithmicensure}{\textbf{saída:}} % Renomear a palavra Output para Saída
\renewcommand{\algorithmicwhile}{\textbf{enquanto}} % Renomear a palavra While para Enquanto
\renewcommand{\algorithmicif}{\textbf{se}} % Renomear a palavra If para Se
\renewcommand{\algorithmicend}{\textbf{fim}} % Renomear a palavra End Output para Fim do
\renewcommand{\algorithmicprocedure}{\textbf{procedimento}} % Renomear a palavra Procedure para Procedimento
\renewcommand{\algorithmicfor}{\textbf{para}} % Renomear a palavra For para Para
\renewcommand{\algorithmicelse}{\textbf{senão}} % Renomear a palavra Else para Senão
\renewcommand{\algorithmicfunction}{\textbf{função}} % Renomear a palavra Function para Função 
\renewcommand{\algorithmicthen}{\textbf{então}} % Renomear a palavra Then para Então
\renewcommand{\algorithmicdo}{\textbf{fazer}} % Renomear a palavra Do para Fazer
% ----------

\newcommand{\quotes}[1]{``#1''} % Comando para inserir um texto dentro de aspas-dupla com espaço após o final

% FLUXOGRAMA: ----------
\usepackage{tikz}
\usetikzlibrary{positioning, shapes.geometric}
\tikzstyle{line}=[draw] % here

\tikzstyle{startstop} = [rectangle, rounded corners, minimum width=3cm,
minimum height=1cm,text centered, draw=black, fill=white!30]
\tikzstyle{arrow} = [thick, ->, >=stealth]
\tikzstyle{process} = [rectangle, minimum width=3cm, minimum height=1cm, text centered, draw=black, fill=white!30]
\tikzstyle{io} = [trapezium, trapezium left angle=70, trapezium right angle=110, minimum width=3cm, minimum height=1cm, text centered, draw=black, fill=white!30]
\tikzstyle{connector} = [draw, circle, fill=white!20]
\tikzstyle{decision} = [diamond, minimum width=3cm, minimum height=1cm, text centered, draw=black, fill=white!30, aspect=3]
% ----------


% CAIXA ENTORNO DAS FÓRMULAS: ----------
\usepackage{dashbox}
\newcommand\dboxed[1]{\dbox{\ensuremath{#1}}}

% ----------

%\usepackage{tabularx} % para colunas com largura ajustável

% -------- CONFIGURAÇÃO DE MARGENS --------
% \usepackage[top=2cm,bottom=2cm,left=3cm,right=2cm]{geometry}

% -------- PACOTE fancyhdr --------
\usepackage{fancyhdr}

\setlength{\marginparwidth}{2cm}

% Ajuste do tamanho do cabeçalho para acompanhar exatamente \textwidth
\setlength{\headwidth}{\textwidth}

% Define um headheight mínimo (para evitar avisos)
\setlength{\headheight}{15pt} 

% Limpa configurações padrão de cabeçalho/rodapé
\pagestyle{fancy}
\fancyhf{} 

% Marcações automáticas de seção para o cabeçalho (opcional)
\renewcommand{\sectionmark}[1]{%
	\markboth{\thesection\ #1}{}%
}

% Texto do cabeçalho à esquerda: seção e título
% Usamos \leftmark pois o \sectionmark acima grava no \leftmark
\fancyhead[L]{\small \leftmark}

% Texto do cabeçalho à direita: número de página
\fancyhead[R]{\small \thepage}

% Linha horizontal no cabeçalho (espessura 0.4pt)
\renewcommand{\headrulewidth}{0.4pt}

% Sobrescreve o título padrão “References” para “Referências”
%\renewcommand{\refname}{Referências}
% Tradução de autoref para o português
\def\figureautorefname{Figura}
\def\tableautorefname{Tabela}
\def\equationautorefname{Equa\c{c}\~ao}

% \hypersetup{backref=true}

\ifdefined\backref
        \renewcommand*{\backref}[1]{}
        \renewcommand*{\backrefentrycount}[2]{#1\ifnum#2>1 ~(#2)\fi}
\fi

\DeclarePairedDelimiter{\ceil}{\lceil}{\rceil}

\onehalfspacing

\usepackage{pgfgantt}

%++++++++++++++++++++++++++++++++++++++++++++++++++++++++++++++++++++++++++++++
% Identificacoes (se o trabalho for em inglês, insira os dados em inglês)
% Para entradas abreviadas de Professora (Profa.) em português escreva: Prof$^\textnormal{a}$.
%++++++++++++++++++++++++++++++++++++++++++++++++++++++++++++++++++++++++++++++
\course{Engenharia Aeronáutica e Mecânica} % Programa de PG ou Curso de Graduação
\area{Propulsão Aeroespacial e Energia} % Área de concentração na PG (Não utilizado no caso de TG)

% Autor do trabalho: Nome Sobrenome
\authorgender{masc} %sexo: masc ou fem
\author{\textbf{EDEN DENIS} Ferreira da Silva Lopes}{Santos}
\itaauthoraddress{Rua João Fonseca dos Santos, 20, Apto. 93}{12.230-088}{São José dos Campos--SP}

% Titulo da Tese/Dissertação
\title{Aplicação de Inteligência Artificial (AI) em Banco de Dados de Emissão de Motores}

% Orientador
\advisorgender{masc} % masc ou fem
\advisor{Prof.~Dr.}{Pedro Teixeira \textbf{LACAVA}}{ITA}

% Coorientador (Caso não haja coorientador, colocar ambas as variáveis \coadvisorgender e \coadvisor comentadas, com um % na frente)
\coadvisorgender{masc}
\coadvisor{}{}  % {Prof.~Dr.}{Carlos Alberto Rocha \textbf{PIMENTEL}}{ITA}

% Pró-reitor da Pós-graduação
\bossgender{fem}
\boss{Prof.~Dra.}{Emília \textbf{VILLANI}}

%Coordenador do curso no caso de TG
\bosscoursegender{masc}
\bosscourse{Prof.~Dr.}{ - } % masc ou fem

% Palavras-Chaves informadas pela Biblioteca -> utilizada na CIP
\kwcip{Engine taxonomy}
\kwcip{Emission classification}
\kwcip{Inteligência Artificial}
\kwcip{Artificial Intelligence}
\kwcip{Python}

% membros da banca examinadora
\examiner{Prof. Dr.}{Pedro Teixeira \textbf{LACAVA}}{President}{ITA}
\examiner{Prof. Dr.}{}{Internal Member}{ITA} % Vinícius \textbf{MALATESTA}}{Internal Member}{ITA}
\examiner{Prof. Dr.}{}{Internal Member}{ITA} % {Elisan dos Santos \textbf{MAGALHÃES}}{Internal Member}{ITA}
\examiner{Prof. Dr.}{}{Internal Member}{ITA} % {Carlos Alberto Rocha \textbf{PIMENTEL}}{External Member}{UFABC}
\examiner{Prof. Dr.}{}{Internal Member}{ITA} % {\textbf{KARL} Peter Burr}{External Member}{UFABC}

% Data da defesa (mês em maiúsculo, se trabalho em inglês, e minúsculo se trabalho em português) 
\date{14}{JULHO}{2025}

% Número CDU - (somente para TG)
\cdu{621.38}

% Folha de Registro do Documento
\newcommand{\advisorTitle}{Prof.~Dr.}
\newcommand{\advisorName}{Pedro Teixeira \textbf{LACAVA}}

\newcommand{\coadvisorTitle}{Prof.~Dr.}
\newcommand{\coadvisorName}{Carlos Alberto Rocha \textbf{PIMENTEL}}

\newcommand{\dateDay}{14}
\newcommand{\dateMonth}{JULHO}
\newcommand{\dateYear}{2025}

% Data de publicação e defesa (usada na folha de registro)
\newcommand{\frdDataPublicacao}{04 de agosto de 2025}
\newcommand{\frdNumeroRegistro}{DCTA/ITA/TD-029/2020}
\newcommand{\frdPalavrasAutor}{Big Data; Engine Emission; Emission Databank; Clustering; Artificial Intelligence}
\newcommand{\frdPalavrasResultado}{Taxonomia; Emissão de Motores; Planejamento de Dados; Inteligência Artificial; Processamento de Dados}
\newcommand{\frdApresentacao}{
	ITA, São José dos Campos. Curso de Mestrado. Programa de Pós-Graduação em Engenharia Aeronáutica e Mecânica. Área de Propulsão Aeroespacial e Energia.\\
	Orientador: \advisorTitle~\advisorName.\\
	Coorientador: \coadvisorTitle~\coadvisorName.\\
	Defesa em \dateDay/\dateMonth/\dateYear. Publicada em \frdDataPublicacao.
}

\newcommand{\PorcentagemDaLargura}{0.80}

\begin{document}

	% Dedicatoria: Nao esqueca essa secao  ... :-)
	\begin{itadedication}
		À todas as formas de inteligência.
	\end{itadedication}

\end{document}  % [x] Revisão do texto
	
	% Agradecimentos:
	\documentclass[12pt]{article} % Tamanho de letra 12, papel A4, Estilo de artigo

\usepackage[utf8]{inputenc} % Pacote para permitir o uso de acentos
% Ajustes de hifenização
\tolerance=1
\emergencystretch=\maxdimen
\hyphenpenalty=10000
\hbadness=10000
\usepackage{lipsum} % Pacote para gerar texto aleatório
\usepackage{ragged2e} % Pacote para mais opções de alinhamento
\usepackage{graphicx} % Pacote para inserir figuras jpeg, pdf e png
\usepackage{ae}
\usepackage{epsfig}
\usepackage{mathtools}
\usepackage[export]{adjustbox}
\usepackage{lineno}
\usepackage{amsthm}
\usepackage{commath}
\usepackage{color, colortbl}
\usepackage{bbm}
\usepackage{setspace} % Pacote para ajustar espaçamento entre linhas

% Instructions to prevent LaTeX from hyphenating words
\tolerance=1
\emergencystretch=\maxdimen
\hyphenpenalty=10000
\hbadness=10000

\graphicspath{{figures/}} % Indica o diretório no qual as imagens estão sendo salvas, evitando que precisamos digitar parte do caminho da pasta que contém as imagens
\usepackage[rightcaption]{sidecap} % Pacote para mais opções de figura, colocar, título da legenda da figura ao lado direito
\usepackage{url} % Pacote para citar referências de sites
\usepackage{makeidx} % Pacote para criar índice
\makeindex % Habilitar a criação do índice
\usepackage{longtable} % Pacote para usar tabelas longas

% ----------
% Pacotes para dar quebra de linha dentro da tabela:
%\usepackage{fourier} 
%\usepackage{array}
\usepackage{makecell}

\renewcommand\theadalign{bc}
\renewcommand\theadfont{\bfseries}
\renewcommand\theadgape{\Gape[4pt]}
\renewcommand\cellgape{\Gape[4pt]}
% ----------

\usepackage{wrapfig} % Pacote para alinhar figuras à esquerda ou à direita do texto
\usepackage{ae}
\usepackage{epsfig}
\usepackage{mathtools}
\usepackage[export]{adjustbox}
\usepackage{lineno}
\usepackage{amsfonts}
\usepackage{amsthm}
\usepackage{commath}
\usepackage{color, colortbl}
\usepackage{bbm}

% ----------
% Pacote para inserir cabeçalho e rodapé:
% Atenção! Se remover o cabeçalho e rodapé, o LaTeX coloca como cabeçalho o padrão dele, que é bem interessante também!
% \usepackage{fancyhdr} % Pacote para inserir cabeçalho e rodapé
% \pagestyle{fancy} % Estilo do cabeçalho e rodapé, neste caso apenas uma linha no cabeçalho
%\lhead{{\footnotesize Prandiano\textregistered}} % Inserir texto no canto esquerdo do cabeçalho
%\chead{centro} % Inserir texto no centro do cabeçalho
%\rhead{{\footnotesize Curso de Python}} % Inserir texto no canto direito do cabeçalho
%\lfoot{{\footnotesize Copyright \textcopyright 2018 Prandiano}} % Inserir texto no canto esquerdo do rodapé
%\cfoot{centro} % Inserir texto no centro do rodapé
%\rfoot{{\footnotesize Reprodução proibida: Lei 9.610 de 19/02/1998}} % Inserir texto no canto direito do rodapé
% O comando \renewcommand re-nova o parâmetro especificado
%\renewcommand{\footrulewidth}{0.4pt} % Inserir uma linha acima do rodapé
% ----------

% ----------
% Pacote para inserir múltiplas figuras:
\usepackage{subfigure} % Pacote para inserir múltiplas figuras
% \usepackage{subcaption} % Pacote incompatível com subfigure
%\usepackage{caption}
% ----------

\usepackage{multicol} % Pacote que permite usar multicolunas
\usepackage{float} % Pacote para usar o "H" do figure
\usepackage[dvipsnames]{xcolor} % Pacote de cores para textos e possivelmente outros
\usepackage{amsmath} % Pacote para usar o cases
\usepackage[makeroom]{cancel} % Pacote para usar o cancelar (em e embaixo)
\usepackage{enumitem} % Pacote para enumerar itens
\usepackage{amssymb} % Pacote para inserir os símbolos dos conjuntos numéricos N, Z, Q, R, C
\usepackage[T1]{fontenc}
\usepackage{enumitem}
\usepackage{textcomp}
\usepackage{pdfpages} % Pacote para inserir página em PDF
\usepackage{courier} % Pacote para usar a fonte courier
%\usepackage{matlab-prettifier}
\usepackage{lscape} % Pacote para mudar sentido da página
\newtheorem{theorem}{Teorema} % Pacote para inserir Teorema
\newtheorem{corollary}{Corollary} % Pacote para inserir Corolário
\newtheorem{remark}{Remark}
\newtheorem{mytheorem}{Definition}
\usepackage{hyperref} % Pacote para inserir hyperref
\usepackage{multirow} % Pacote para usar os comandos \multirow
\usepackage{pdfpages} % Pacote para incluir PDF's
\usepackage{listingsutf8} % Pacote para inserir arquivos externos tipo .txt
\usepackage{etoolbox}
\AtBeginEnvironment{verbatim}{\footnotesize}

\usepackage{standalone} % Pacote para inserir sub-arquivos LaTeX
\usepackage{bookmark} % Pacote para NÃO tratar o Índice como Apêndice e colocar letra

% Nome da listagem em português
\renewcommand{\lstlistingname}{Listagem}

% REFERÊNCIAS: ----------
%\usepackage[numbers, compress]{natbib} % Pacote para usar as referências bibliográficas
%\usepackage{usebib}
% \bibliographystyle{references} % Estilo das referências
%\setcitestyle{authoryear,aysep={},open={(},close={)}}
%\newcommand{\myand}{e}

% ----------


% PARÂMETROS DA FUNÇÃO LISTINGS: ----------
\lstset
{
    basicstyle=\small\ttfamily,
    commentstyle=\color{Green},
    keywordstyle=\color{Cerulean},
    frame=single,
    language=Python,
    morekeywords={True, False},
    numbers=left,
    numbersep=10pt,
    numberstyle=\footnotesize\color{Gray},
    showstringspaces=false,
    stringstyle=\color{Mulberry},
    tabsize=3,
    breaklines=true, % Este comando é para quebrar uma linha comprida
    postbreak=\mbox{\textcolor{red}{$\hookrightarrow$}\space}, % Este comando é para indicar com uma seta vermelha que a linha continua embaixo 
    extendedchars=true,
    literate=
        {á}{{\'a}}1 {Á}{{\'A}}1 {à}{{\`a}}1 {À}{{\`A}}1 {ã}{{\~a}}1 {Ã}{{\~A}}1 {â}{{\^a}}1 {Â}{{\^A}}1
        {ç}{{\c{c}}}1 {Ç}{{\c{C}}}1 {é}{{\'e}}1 {É}{{\'E}}1 {ê}{{\^e}}1 {Ê}{{\^E}}1 {í}{{\'\i}}1 {Í}{{\'I}}1
        {ó}{{\'o}}1 {Ó}{{\'O}}1 {ô}{{\^o}}1 {Ô}{{\^O}}1 {õ}{{\~o}}1 {Õ}{{\~O}}1 {ú}{{\'u}}1 {Ú}{{\'U}}1
        {ü}{{\"u}}1 {Ü}{{\"U}}1 {“}{{``}}1 {”}{{''}}1 {’}{{'}}1 {—}{{--}}1 {–}{{--}}1 {−}{{-}}1
        {±}{{$\pm$}}1 {σ}{{$\sigma$}}1 {→}{{$\to$}}1 {…}{{\ldots}}1
}

\lstset{inputencoding=utf8/latin1}

\lstdefinelanguage{Matlab}{
	language=Matlab,
	basicstyle=\ttfamily\small,
	keywordstyle=\color{blue}\bfseries,
	stringstyle=\color{red},
	commentstyle=\color{gray}\itshape,
	morekeywords={function, end, if, else, for, while, switch, case, otherwise},
	sensitive=true
}
% ----------

% ----------
\usepackage{nomencl} % Pacote para criar nomenclatura e lista de nomenclatura
\makenomenclature % Criar a nomenclatura
\renewcommand{\nomname}{Lista de Nomenclaturas} % Renomear o título da nomenclatura
% ----------

\usepackage{amsmath} % Pacote para inserir matrix com chaves, entre outras coisas

\usepackage{booktabs} % Pacote para usar \toprule, \midrule e \bottomrule


% EXECUTAR CÁLCULOS A PARTIR DA LINHA DE COMANDO: ---------
%\usepackage{siunitx, xfp} % Pacote para executar cálculos a partir da linha de comando
%\usepackage{l3regex}%unless you have very new expl3
%%\sisetup{locale=DE}
%%\let\oldnum\num
%%\ExplSyntaxOn
%\def\num#1{\def\tmp{#1}\regex_replace_all:nnN{,([0-9])}{.\1}{\tmp}\oldnum{\tmp}}
%\ExplSyntaxOff
% ----------


%\pagecolor{black} % Alterar a cor da página
%\color{white} % Alterar a cor do texto

% ----------
% Pacotes para inserir algoritmos e pseudo-códigos:
\setlength{\marginparwidth}{2.5cm}
\usepackage[colorinlistoftodos]{todonotes}
\usepackage{algorithm} % Pacote para inserir algoritmos
\usepackage{algpseudocode} % Pacote para inserir algoritmos e pseudo-códigos
\floatname{algorithm}{Algoritmo} % Renomear a palavra Algorithm para Procedimento
\renewcommand{\algorithmicrequire}{\textbf{entrada:}} % Renomear a palavra Input para Entrada
\renewcommand{\algorithmicensure}{\textbf{saída:}} % Renomear a palavra Output para Saída
\renewcommand{\algorithmicwhile}{\textbf{enquanto}} % Renomear a palavra While para Enquanto
\renewcommand{\algorithmicif}{\textbf{se}} % Renomear a palavra If para Se
\renewcommand{\algorithmicend}{\textbf{fim}} % Renomear a palavra End Output para Fim do
\renewcommand{\algorithmicprocedure}{\textbf{procedimento}} % Renomear a palavra Procedure para Procedimento
\renewcommand{\algorithmicfor}{\textbf{para}} % Renomear a palavra For para Para
\renewcommand{\algorithmicelse}{\textbf{senão}} % Renomear a palavra Else para Senão
\renewcommand{\algorithmicfunction}{\textbf{função}} % Renomear a palavra Function para Função 
\renewcommand{\algorithmicthen}{\textbf{então}} % Renomear a palavra Then para Então
\renewcommand{\algorithmicdo}{\textbf{fazer}} % Renomear a palavra Do para Fazer
% ----------

\newcommand{\quotes}[1]{``#1''} % Comando para inserir um texto dentro de aspas-dupla com espaço após o final

% FLUXOGRAMA: ----------
\usepackage{tikz}
\usetikzlibrary{positioning, shapes.geometric}
\tikzstyle{line}=[draw] % here

\tikzstyle{startstop} = [rectangle, rounded corners, minimum width=3cm,
minimum height=1cm,text centered, draw=black, fill=white!30]
\tikzstyle{arrow} = [thick, ->, >=stealth]
\tikzstyle{process} = [rectangle, minimum width=3cm, minimum height=1cm, text centered, draw=black, fill=white!30]
\tikzstyle{io} = [trapezium, trapezium left angle=70, trapezium right angle=110, minimum width=3cm, minimum height=1cm, text centered, draw=black, fill=white!30]
\tikzstyle{connector} = [draw, circle, fill=white!20]
\tikzstyle{decision} = [diamond, minimum width=3cm, minimum height=1cm, text centered, draw=black, fill=white!30, aspect=3]
% ----------


% CAIXA ENTORNO DAS FÓRMULAS: ----------
\usepackage{dashbox}
\newcommand\dboxed[1]{\dbox{\ensuremath{#1}}}

% ----------

%\usepackage{tabularx} % para colunas com largura ajustável

% -------- CONFIGURAÇÃO DE MARGENS --------
% \usepackage[top=2cm,bottom=2cm,left=3cm,right=2cm]{geometry}

% -------- PACOTE fancyhdr --------
\usepackage{fancyhdr}

\setlength{\marginparwidth}{2cm}

% Ajuste do tamanho do cabeçalho para acompanhar exatamente \textwidth
\setlength{\headwidth}{\textwidth}

% Define um headheight mínimo (para evitar avisos)
\setlength{\headheight}{15pt} 

% Limpa configurações padrão de cabeçalho/rodapé
\pagestyle{fancy}
\fancyhf{} 

% Marcações automáticas de seção para o cabeçalho (opcional)
\renewcommand{\sectionmark}[1]{%
	\markboth{\thesection\ #1}{}%
}

% Texto do cabeçalho à esquerda: seção e título
% Usamos \leftmark pois o \sectionmark acima grava no \leftmark
\fancyhead[L]{\small \leftmark}

% Texto do cabeçalho à direita: número de página
\fancyhead[R]{\small \thepage}

% Linha horizontal no cabeçalho (espessura 0.4pt)
\renewcommand{\headrulewidth}{0.4pt}

% Sobrescreve o título padrão “References” para “Referências”
%\renewcommand{\refname}{Referências}
% Tradução de autoref para o português
\def\figureautorefname{Figura}
\def\tableautorefname{Tabela}
\def\equationautorefname{Equa\c{c}\~ao}

% \hypersetup{backref=true}

\ifdefined\backref
        \renewcommand*{\backref}[1]{}
        \renewcommand*{\backrefentrycount}[2]{#1\ifnum#2>1 ~(#2)\fi}
\fi

\DeclarePairedDelimiter{\ceil}{\lceil}{\rceil}

\onehalfspacing

\usepackage{pgfgantt}

%++++++++++++++++++++++++++++++++++++++++++++++++++++++++++++++++++++++++++++++
% Identificacoes (se o trabalho for em inglês, insira os dados em inglês)
% Para entradas abreviadas de Professora (Profa.) em português escreva: Prof$^\textnormal{a}$.
%++++++++++++++++++++++++++++++++++++++++++++++++++++++++++++++++++++++++++++++
\course{Engenharia Aeronáutica e Mecânica} % Programa de PG ou Curso de Graduação
\area{Propulsão Aeroespacial e Energia} % Área de concentração na PG (Não utilizado no caso de TG)

% Autor do trabalho: Nome Sobrenome
\authorgender{masc} %sexo: masc ou fem
\author{\textbf{EDEN DENIS} Ferreira da Silva Lopes}{Santos}
\itaauthoraddress{Rua João Fonseca dos Santos, 20, Apto. 93}{12.230-088}{São José dos Campos--SP}

% Titulo da Tese/Dissertação
\title{Aplicação de Inteligência Artificial (AI) em Banco de Dados de Emissão de Motores}

% Orientador
\advisorgender{masc} % masc ou fem
\advisor{Prof.~Dr.}{Pedro Teixeira \textbf{LACAVA}}{ITA}

% Coorientador (Caso não haja coorientador, colocar ambas as variáveis \coadvisorgender e \coadvisor comentadas, com um % na frente)
\coadvisorgender{masc}
\coadvisor{}{}  % {Prof.~Dr.}{Carlos Alberto Rocha \textbf{PIMENTEL}}{ITA}

% Pró-reitor da Pós-graduação
\bossgender{fem}
\boss{Prof.~Dra.}{Emília \textbf{VILLANI}}

%Coordenador do curso no caso de TG
\bosscoursegender{masc}
\bosscourse{Prof.~Dr.}{ - } % masc ou fem

% Palavras-Chaves informadas pela Biblioteca -> utilizada na CIP
\kwcip{Engine taxonomy}
\kwcip{Emission classification}
\kwcip{Inteligência Artificial}
\kwcip{Artificial Intelligence}
\kwcip{Python}

% membros da banca examinadora
\examiner{Prof. Dr.}{Pedro Teixeira \textbf{LACAVA}}{President}{ITA}
\examiner{Prof. Dr.}{}{Internal Member}{ITA} % Vinícius \textbf{MALATESTA}}{Internal Member}{ITA}
\examiner{Prof. Dr.}{}{Internal Member}{ITA} % {Elisan dos Santos \textbf{MAGALHÃES}}{Internal Member}{ITA}
\examiner{Prof. Dr.}{}{Internal Member}{ITA} % {Carlos Alberto Rocha \textbf{PIMENTEL}}{External Member}{UFABC}
\examiner{Prof. Dr.}{}{Internal Member}{ITA} % {\textbf{KARL} Peter Burr}{External Member}{UFABC}

% Data da defesa (mês em maiúsculo, se trabalho em inglês, e minúsculo se trabalho em português) 
\date{14}{JULHO}{2025}

% Número CDU - (somente para TG)
\cdu{621.38}

% Folha de Registro do Documento
\newcommand{\advisorTitle}{Prof.~Dr.}
\newcommand{\advisorName}{Pedro Teixeira \textbf{LACAVA}}

\newcommand{\coadvisorTitle}{Prof.~Dr.}
\newcommand{\coadvisorName}{Carlos Alberto Rocha \textbf{PIMENTEL}}

\newcommand{\dateDay}{14}
\newcommand{\dateMonth}{JULHO}
\newcommand{\dateYear}{2025}

% Data de publicação e defesa (usada na folha de registro)
\newcommand{\frdDataPublicacao}{04 de agosto de 2025}
\newcommand{\frdNumeroRegistro}{DCTA/ITA/TD-029/2020}
\newcommand{\frdPalavrasAutor}{Big Data; Engine Emission; Emission Databank; Clustering; Artificial Intelligence}
\newcommand{\frdPalavrasResultado}{Taxonomia; Emissão de Motores; Planejamento de Dados; Inteligência Artificial; Processamento de Dados}
\newcommand{\frdApresentacao}{
	ITA, São José dos Campos. Curso de Mestrado. Programa de Pós-Graduação em Engenharia Aeronáutica e Mecânica. Área de Propulsão Aeroespacial e Energia.\\
	Orientador: \advisorTitle~\advisorName.\\
	Coorientador: \coadvisorTitle~\coadvisorName.\\
	Defesa em \dateDay/\dateMonth/\dateYear. Publicada em \frdDataPublicacao.
}

\newcommand{\PorcentagemDaLargura}{0.80}

\begin{document}

    % Agradecimentos
    \begin{itathanks}
        \textcolor{black}{À minha família, em particular à minha querida avó Severina Lira da Silva, que sempre esteve ao meu lado, por todo o amor, inspiração, motivação, cuja sabedoria simples, sempre me incentivou aos estudos o que me trouxe até esta grande realização.}

        \textcolor{black}{Ao meu avô, Antônio Ferreira da Silva, cuja sabedoria e força, nos trouxe até aqui.}
        
        \textcolor{black}{À minha mãe que, embora tenha tido muitos momentos de ansiedade e dúvidas, nunca duvidou de mim e sempre me ergueu quando mais precisei.}

        \textcolor{black}{Ao meu filho, que tem se tornado fonte de esperança e propósito, por iluminar cada dia e me lembrar do impacto que este trabalho terá nas próximas gerações.}

        \textcolor{black}{Ao meu orientador, \itaadvisortitle~\itaadvisorname, pela confiança, liberdade, compreensão, que tanto me impulsionaram a crescer como pessoa em primeiro lugar, mas também como profissional.}
    \end{itathanks}

\end{document}
  % [x] Revisão do texto
	
	% Epigrafo:
	\documentclass[12pt]{article} % Tamanho de letra 12, papel A4, Estilo de artigo

\usepackage[utf8]{inputenc} % Pacote para permitir o uso de acentos
% Ajustes de hifenização
\tolerance=1
\emergencystretch=\maxdimen
\hyphenpenalty=10000
\hbadness=10000
\usepackage{lipsum} % Pacote para gerar texto aleatório
\usepackage{ragged2e} % Pacote para mais opções de alinhamento
\usepackage{graphicx} % Pacote para inserir figuras jpeg, pdf e png
\usepackage{ae}
\usepackage{epsfig}
\usepackage{mathtools}
\usepackage[export]{adjustbox}
\usepackage{lineno}
\usepackage{amsthm}
\usepackage{commath}
\usepackage{color, colortbl}
\usepackage{bbm}
\usepackage{setspace} % Pacote para ajustar espaçamento entre linhas

% Instructions to prevent LaTeX from hyphenating words
\tolerance=1
\emergencystretch=\maxdimen
\hyphenpenalty=10000
\hbadness=10000

\graphicspath{{figures/}} % Indica o diretório no qual as imagens estão sendo salvas, evitando que precisamos digitar parte do caminho da pasta que contém as imagens
\usepackage[rightcaption]{sidecap} % Pacote para mais opções de figura, colocar, título da legenda da figura ao lado direito
\usepackage{url} % Pacote para citar referências de sites
\usepackage{makeidx} % Pacote para criar índice
\makeindex % Habilitar a criação do índice
\usepackage{longtable} % Pacote para usar tabelas longas

% ----------
% Pacotes para dar quebra de linha dentro da tabela:
%\usepackage{fourier} 
%\usepackage{array}
\usepackage{makecell}

\renewcommand\theadalign{bc}
\renewcommand\theadfont{\bfseries}
\renewcommand\theadgape{\Gape[4pt]}
\renewcommand\cellgape{\Gape[4pt]}
% ----------

\usepackage{wrapfig} % Pacote para alinhar figuras à esquerda ou à direita do texto
\usepackage{ae}
\usepackage{epsfig}
\usepackage{mathtools}
\usepackage[export]{adjustbox}
\usepackage{lineno}
\usepackage{amsfonts}
\usepackage{amsthm}
\usepackage{commath}
\usepackage{color, colortbl}
\usepackage{bbm}

% ----------
% Pacote para inserir cabeçalho e rodapé:
% Atenção! Se remover o cabeçalho e rodapé, o LaTeX coloca como cabeçalho o padrão dele, que é bem interessante também!
% \usepackage{fancyhdr} % Pacote para inserir cabeçalho e rodapé
% \pagestyle{fancy} % Estilo do cabeçalho e rodapé, neste caso apenas uma linha no cabeçalho
%\lhead{{\footnotesize Prandiano\textregistered}} % Inserir texto no canto esquerdo do cabeçalho
%\chead{centro} % Inserir texto no centro do cabeçalho
%\rhead{{\footnotesize Curso de Python}} % Inserir texto no canto direito do cabeçalho
%\lfoot{{\footnotesize Copyright \textcopyright 2018 Prandiano}} % Inserir texto no canto esquerdo do rodapé
%\cfoot{centro} % Inserir texto no centro do rodapé
%\rfoot{{\footnotesize Reprodução proibida: Lei 9.610 de 19/02/1998}} % Inserir texto no canto direito do rodapé
% O comando \renewcommand re-nova o parâmetro especificado
%\renewcommand{\footrulewidth}{0.4pt} % Inserir uma linha acima do rodapé
% ----------

% ----------
% Pacote para inserir múltiplas figuras:
\usepackage{subfigure} % Pacote para inserir múltiplas figuras
% \usepackage{subcaption} % Pacote incompatível com subfigure
%\usepackage{caption}
% ----------

\usepackage{multicol} % Pacote que permite usar multicolunas
\usepackage{float} % Pacote para usar o "H" do figure
\usepackage[dvipsnames]{xcolor} % Pacote de cores para textos e possivelmente outros
\usepackage{amsmath} % Pacote para usar o cases
\usepackage[makeroom]{cancel} % Pacote para usar o cancelar (em e embaixo)
\usepackage{enumitem} % Pacote para enumerar itens
\usepackage{amssymb} % Pacote para inserir os símbolos dos conjuntos numéricos N, Z, Q, R, C
\usepackage[T1]{fontenc}
\usepackage{enumitem}
\usepackage{textcomp}
\usepackage{pdfpages} % Pacote para inserir página em PDF
\usepackage{courier} % Pacote para usar a fonte courier
%\usepackage{matlab-prettifier}
\usepackage{lscape} % Pacote para mudar sentido da página
\newtheorem{theorem}{Teorema} % Pacote para inserir Teorema
\newtheorem{corollary}{Corollary} % Pacote para inserir Corolário
\newtheorem{remark}{Remark}
\newtheorem{mytheorem}{Definition}
\usepackage{hyperref} % Pacote para inserir hyperref
\usepackage{multirow} % Pacote para usar os comandos \multirow
\usepackage{pdfpages} % Pacote para incluir PDF's
\usepackage{listingsutf8} % Pacote para inserir arquivos externos tipo .txt
\usepackage{etoolbox}
\AtBeginEnvironment{verbatim}{\footnotesize}

\usepackage{standalone} % Pacote para inserir sub-arquivos LaTeX
\usepackage{bookmark} % Pacote para NÃO tratar o Índice como Apêndice e colocar letra

% Nome da listagem em português
\renewcommand{\lstlistingname}{Listagem}

% REFERÊNCIAS: ----------
%\usepackage[numbers, compress]{natbib} % Pacote para usar as referências bibliográficas
%\usepackage{usebib}
% \bibliographystyle{references} % Estilo das referências
%\setcitestyle{authoryear,aysep={},open={(},close={)}}
%\newcommand{\myand}{e}

% ----------


% PARÂMETROS DA FUNÇÃO LISTINGS: ----------
\lstset
{
    basicstyle=\small\ttfamily,
    commentstyle=\color{Green},
    keywordstyle=\color{Cerulean},
    frame=single,
    language=Python,
    morekeywords={True, False},
    numbers=left,
    numbersep=10pt,
    numberstyle=\footnotesize\color{Gray},
    showstringspaces=false,
    stringstyle=\color{Mulberry},
    tabsize=3,
    breaklines=true, % Este comando é para quebrar uma linha comprida
    postbreak=\mbox{\textcolor{red}{$\hookrightarrow$}\space}, % Este comando é para indicar com uma seta vermelha que a linha continua embaixo 
    extendedchars=true,
    literate=
        {á}{{\'a}}1 {Á}{{\'A}}1 {à}{{\`a}}1 {À}{{\`A}}1 {ã}{{\~a}}1 {Ã}{{\~A}}1 {â}{{\^a}}1 {Â}{{\^A}}1
        {ç}{{\c{c}}}1 {Ç}{{\c{C}}}1 {é}{{\'e}}1 {É}{{\'E}}1 {ê}{{\^e}}1 {Ê}{{\^E}}1 {í}{{\'\i}}1 {Í}{{\'I}}1
        {ó}{{\'o}}1 {Ó}{{\'O}}1 {ô}{{\^o}}1 {Ô}{{\^O}}1 {õ}{{\~o}}1 {Õ}{{\~O}}1 {ú}{{\'u}}1 {Ú}{{\'U}}1
        {ü}{{\"u}}1 {Ü}{{\"U}}1 {“}{{``}}1 {”}{{''}}1 {’}{{'}}1 {—}{{--}}1 {–}{{--}}1 {−}{{-}}1
        {±}{{$\pm$}}1 {σ}{{$\sigma$}}1 {→}{{$\to$}}1 {…}{{\ldots}}1
}

\lstset{inputencoding=utf8/latin1}

\lstdefinelanguage{Matlab}{
	language=Matlab,
	basicstyle=\ttfamily\small,
	keywordstyle=\color{blue}\bfseries,
	stringstyle=\color{red},
	commentstyle=\color{gray}\itshape,
	morekeywords={function, end, if, else, for, while, switch, case, otherwise},
	sensitive=true
}
% ----------

% ----------
\usepackage{nomencl} % Pacote para criar nomenclatura e lista de nomenclatura
\makenomenclature % Criar a nomenclatura
\renewcommand{\nomname}{Lista de Nomenclaturas} % Renomear o título da nomenclatura
% ----------

\usepackage{amsmath} % Pacote para inserir matrix com chaves, entre outras coisas

\usepackage{booktabs} % Pacote para usar \toprule, \midrule e \bottomrule


% EXECUTAR CÁLCULOS A PARTIR DA LINHA DE COMANDO: ---------
%\usepackage{siunitx, xfp} % Pacote para executar cálculos a partir da linha de comando
%\usepackage{l3regex}%unless you have very new expl3
%%\sisetup{locale=DE}
%%\let\oldnum\num
%%\ExplSyntaxOn
%\def\num#1{\def\tmp{#1}\regex_replace_all:nnN{,([0-9])}{.\1}{\tmp}\oldnum{\tmp}}
%\ExplSyntaxOff
% ----------


%\pagecolor{black} % Alterar a cor da página
%\color{white} % Alterar a cor do texto

% ----------
% Pacotes para inserir algoritmos e pseudo-códigos:
\setlength{\marginparwidth}{2.5cm}
\usepackage[colorinlistoftodos]{todonotes}
\usepackage{algorithm} % Pacote para inserir algoritmos
\usepackage{algpseudocode} % Pacote para inserir algoritmos e pseudo-códigos
\floatname{algorithm}{Algoritmo} % Renomear a palavra Algorithm para Procedimento
\renewcommand{\algorithmicrequire}{\textbf{entrada:}} % Renomear a palavra Input para Entrada
\renewcommand{\algorithmicensure}{\textbf{saída:}} % Renomear a palavra Output para Saída
\renewcommand{\algorithmicwhile}{\textbf{enquanto}} % Renomear a palavra While para Enquanto
\renewcommand{\algorithmicif}{\textbf{se}} % Renomear a palavra If para Se
\renewcommand{\algorithmicend}{\textbf{fim}} % Renomear a palavra End Output para Fim do
\renewcommand{\algorithmicprocedure}{\textbf{procedimento}} % Renomear a palavra Procedure para Procedimento
\renewcommand{\algorithmicfor}{\textbf{para}} % Renomear a palavra For para Para
\renewcommand{\algorithmicelse}{\textbf{senão}} % Renomear a palavra Else para Senão
\renewcommand{\algorithmicfunction}{\textbf{função}} % Renomear a palavra Function para Função 
\renewcommand{\algorithmicthen}{\textbf{então}} % Renomear a palavra Then para Então
\renewcommand{\algorithmicdo}{\textbf{fazer}} % Renomear a palavra Do para Fazer
% ----------

\newcommand{\quotes}[1]{``#1''} % Comando para inserir um texto dentro de aspas-dupla com espaço após o final

% FLUXOGRAMA: ----------
\usepackage{tikz}
\usetikzlibrary{positioning, shapes.geometric}
\tikzstyle{line}=[draw] % here

\tikzstyle{startstop} = [rectangle, rounded corners, minimum width=3cm,
minimum height=1cm,text centered, draw=black, fill=white!30]
\tikzstyle{arrow} = [thick, ->, >=stealth]
\tikzstyle{process} = [rectangle, minimum width=3cm, minimum height=1cm, text centered, draw=black, fill=white!30]
\tikzstyle{io} = [trapezium, trapezium left angle=70, trapezium right angle=110, minimum width=3cm, minimum height=1cm, text centered, draw=black, fill=white!30]
\tikzstyle{connector} = [draw, circle, fill=white!20]
\tikzstyle{decision} = [diamond, minimum width=3cm, minimum height=1cm, text centered, draw=black, fill=white!30, aspect=3]
% ----------


% CAIXA ENTORNO DAS FÓRMULAS: ----------
\usepackage{dashbox}
\newcommand\dboxed[1]{\dbox{\ensuremath{#1}}}

% ----------

%\usepackage{tabularx} % para colunas com largura ajustável

% -------- CONFIGURAÇÃO DE MARGENS --------
% \usepackage[top=2cm,bottom=2cm,left=3cm,right=2cm]{geometry}

% -------- PACOTE fancyhdr --------
\usepackage{fancyhdr}

\setlength{\marginparwidth}{2cm}

% Ajuste do tamanho do cabeçalho para acompanhar exatamente \textwidth
\setlength{\headwidth}{\textwidth}

% Define um headheight mínimo (para evitar avisos)
\setlength{\headheight}{15pt} 

% Limpa configurações padrão de cabeçalho/rodapé
\pagestyle{fancy}
\fancyhf{} 

% Marcações automáticas de seção para o cabeçalho (opcional)
\renewcommand{\sectionmark}[1]{%
	\markboth{\thesection\ #1}{}%
}

% Texto do cabeçalho à esquerda: seção e título
% Usamos \leftmark pois o \sectionmark acima grava no \leftmark
\fancyhead[L]{\small \leftmark}

% Texto do cabeçalho à direita: número de página
\fancyhead[R]{\small \thepage}

% Linha horizontal no cabeçalho (espessura 0.4pt)
\renewcommand{\headrulewidth}{0.4pt}

% Sobrescreve o título padrão “References” para “Referências”
%\renewcommand{\refname}{Referências}
% Tradução de autoref para o português
\def\figureautorefname{Figura}
\def\tableautorefname{Tabela}
\def\equationautorefname{Equa\c{c}\~ao}

% \hypersetup{backref=true}

\ifdefined\backref
        \renewcommand*{\backref}[1]{}
        \renewcommand*{\backrefentrycount}[2]{#1\ifnum#2>1 ~(#2)\fi}
\fi

\DeclarePairedDelimiter{\ceil}{\lceil}{\rceil}

\onehalfspacing

\usepackage{pgfgantt}

%++++++++++++++++++++++++++++++++++++++++++++++++++++++++++++++++++++++++++++++
% Identificacoes (se o trabalho for em inglês, insira os dados em inglês)
% Para entradas abreviadas de Professora (Profa.) em português escreva: Prof$^\textnormal{a}$.
%++++++++++++++++++++++++++++++++++++++++++++++++++++++++++++++++++++++++++++++
\course{Engenharia Aeronáutica e Mecânica} % Programa de PG ou Curso de Graduação
\area{Propulsão Aeroespacial e Energia} % Área de concentração na PG (Não utilizado no caso de TG)

% Autor do trabalho: Nome Sobrenome
\authorgender{masc} %sexo: masc ou fem
\author{\textbf{EDEN DENIS} Ferreira da Silva Lopes}{Santos}
\itaauthoraddress{Rua João Fonseca dos Santos, 20, Apto. 93}{12.230-088}{São José dos Campos--SP}

% Titulo da Tese/Dissertação
\title{Aplicação de Inteligência Artificial (AI) em Banco de Dados de Emissão de Motores}

% Orientador
\advisorgender{masc} % masc ou fem
\advisor{Prof.~Dr.}{Pedro Teixeira \textbf{LACAVA}}{ITA}

% Coorientador (Caso não haja coorientador, colocar ambas as variáveis \coadvisorgender e \coadvisor comentadas, com um % na frente)
\coadvisorgender{masc}
\coadvisor{}{}  % {Prof.~Dr.}{Carlos Alberto Rocha \textbf{PIMENTEL}}{ITA}

% Pró-reitor da Pós-graduação
\bossgender{fem}
\boss{Prof.~Dra.}{Emília \textbf{VILLANI}}

%Coordenador do curso no caso de TG
\bosscoursegender{masc}
\bosscourse{Prof.~Dr.}{ - } % masc ou fem

% Palavras-Chaves informadas pela Biblioteca -> utilizada na CIP
\kwcip{Engine taxonomy}
\kwcip{Emission classification}
\kwcip{Inteligência Artificial}
\kwcip{Artificial Intelligence}
\kwcip{Python}

% membros da banca examinadora
\examiner{Prof. Dr.}{Pedro Teixeira \textbf{LACAVA}}{President}{ITA}
\examiner{Prof. Dr.}{}{Internal Member}{ITA} % Vinícius \textbf{MALATESTA}}{Internal Member}{ITA}
\examiner{Prof. Dr.}{}{Internal Member}{ITA} % {Elisan dos Santos \textbf{MAGALHÃES}}{Internal Member}{ITA}
\examiner{Prof. Dr.}{}{Internal Member}{ITA} % {Carlos Alberto Rocha \textbf{PIMENTEL}}{External Member}{UFABC}
\examiner{Prof. Dr.}{}{Internal Member}{ITA} % {\textbf{KARL} Peter Burr}{External Member}{UFABC}

% Data da defesa (mês em maiúsculo, se trabalho em inglês, e minúsculo se trabalho em português) 
\date{14}{JULHO}{2025}

% Número CDU - (somente para TG)
\cdu{621.38}

% Folha de Registro do Documento
\newcommand{\advisorTitle}{Prof.~Dr.}
\newcommand{\advisorName}{Pedro Teixeira \textbf{LACAVA}}

\newcommand{\coadvisorTitle}{Prof.~Dr.}
\newcommand{\coadvisorName}{Carlos Alberto Rocha \textbf{PIMENTEL}}

\newcommand{\dateDay}{14}
\newcommand{\dateMonth}{JULHO}
\newcommand{\dateYear}{2025}

% Data de publicação e defesa (usada na folha de registro)
\newcommand{\frdDataPublicacao}{04 de agosto de 2025}
\newcommand{\frdNumeroRegistro}{DCTA/ITA/TD-029/2020}
\newcommand{\frdPalavrasAutor}{Big Data; Engine Emission; Emission Databank; Clustering; Artificial Intelligence}
\newcommand{\frdPalavrasResultado}{Taxonomia; Emissão de Motores; Planejamento de Dados; Inteligência Artificial; Processamento de Dados}
\newcommand{\frdApresentacao}{
	ITA, São José dos Campos. Curso de Mestrado. Programa de Pós-Graduação em Engenharia Aeronáutica e Mecânica. Área de Propulsão Aeroespacial e Energia.\\
	Orientador: \advisorTitle~\advisorName.\\
	Coorientador: \coadvisorTitle~\coadvisorName.\\
	Defesa em \dateDay/\dateMonth/\dateYear. Publicada em \frdDataPublicacao.
}

\newcommand{\PorcentagemDaLargura}{0.80}

\begin{document}
	
	% Epígrafe:
	\thispagestyle{empty}
	\ifhyperref\pdfbookmark[0]{\nameepigraphe}{epigrafe}\fi
	\begin{flushright}
		\begin{spacing}{1}
			\mbox{}\vfill
			\vspace*{1.0cm}
			{\sffamily\itshape
				\textcolor{black}{``Bendita seja todas as formas de inteligência.''}}. \\
			--- \textcolor{black}{\textsc{Atribuída a Clarice Linspector (1920 - 1977); Frase citada no filme Animatrix (2003)}} \\
			\vspace*{1.0cm}
		\end{spacing}
	\end{flushright}

\end{document}
  % [x] Revisão do texto
	
	% Resumo
	\documentclass[12pt]{article} % Tamanho de letra 12, papel A4, Estilo de artigo

\usepackage[utf8]{inputenc} % Pacote para permitir o uso de acentos
% Ajustes de hifenização
\tolerance=1
\emergencystretch=\maxdimen
\hyphenpenalty=10000
\hbadness=10000
\usepackage{lipsum} % Pacote para gerar texto aleatório
\usepackage{ragged2e} % Pacote para mais opções de alinhamento
\usepackage{graphicx} % Pacote para inserir figuras jpeg, pdf e png
\usepackage{ae}
\usepackage{epsfig}
\usepackage{mathtools}
\usepackage[export]{adjustbox}
\usepackage{lineno}
\usepackage{amsthm}
\usepackage{commath}
\usepackage{color, colortbl}
\usepackage{bbm}
\usepackage{setspace} % Pacote para ajustar espaçamento entre linhas

% Instructions to prevent LaTeX from hyphenating words
\tolerance=1
\emergencystretch=\maxdimen
\hyphenpenalty=10000
\hbadness=10000

\graphicspath{{figures/}} % Indica o diretório no qual as imagens estão sendo salvas, evitando que precisamos digitar parte do caminho da pasta que contém as imagens
\usepackage[rightcaption]{sidecap} % Pacote para mais opções de figura, colocar, título da legenda da figura ao lado direito
\usepackage{url} % Pacote para citar referências de sites
\usepackage{makeidx} % Pacote para criar índice
\makeindex % Habilitar a criação do índice
\usepackage{longtable} % Pacote para usar tabelas longas

% ----------
% Pacotes para dar quebra de linha dentro da tabela:
%\usepackage{fourier} 
%\usepackage{array}
\usepackage{makecell}

\renewcommand\theadalign{bc}
\renewcommand\theadfont{\bfseries}
\renewcommand\theadgape{\Gape[4pt]}
\renewcommand\cellgape{\Gape[4pt]}
% ----------

\usepackage{wrapfig} % Pacote para alinhar figuras à esquerda ou à direita do texto
\usepackage{ae}
\usepackage{epsfig}
\usepackage{mathtools}
\usepackage[export]{adjustbox}
\usepackage{lineno}
\usepackage{amsfonts}
\usepackage{amsthm}
\usepackage{commath}
\usepackage{color, colortbl}
\usepackage{bbm}

% ----------
% Pacote para inserir cabeçalho e rodapé:
% Atenção! Se remover o cabeçalho e rodapé, o LaTeX coloca como cabeçalho o padrão dele, que é bem interessante também!
% \usepackage{fancyhdr} % Pacote para inserir cabeçalho e rodapé
% \pagestyle{fancy} % Estilo do cabeçalho e rodapé, neste caso apenas uma linha no cabeçalho
%\lhead{{\footnotesize Prandiano\textregistered}} % Inserir texto no canto esquerdo do cabeçalho
%\chead{centro} % Inserir texto no centro do cabeçalho
%\rhead{{\footnotesize Curso de Python}} % Inserir texto no canto direito do cabeçalho
%\lfoot{{\footnotesize Copyright \textcopyright 2018 Prandiano}} % Inserir texto no canto esquerdo do rodapé
%\cfoot{centro} % Inserir texto no centro do rodapé
%\rfoot{{\footnotesize Reprodução proibida: Lei 9.610 de 19/02/1998}} % Inserir texto no canto direito do rodapé
% O comando \renewcommand re-nova o parâmetro especificado
%\renewcommand{\footrulewidth}{0.4pt} % Inserir uma linha acima do rodapé
% ----------

% ----------
% Pacote para inserir múltiplas figuras:
\usepackage{subfigure} % Pacote para inserir múltiplas figuras
% \usepackage{subcaption} % Pacote incompatível com subfigure
%\usepackage{caption}
% ----------

\usepackage{multicol} % Pacote que permite usar multicolunas
\usepackage{float} % Pacote para usar o "H" do figure
\usepackage[dvipsnames]{xcolor} % Pacote de cores para textos e possivelmente outros
\usepackage{amsmath} % Pacote para usar o cases
\usepackage[makeroom]{cancel} % Pacote para usar o cancelar (em e embaixo)
\usepackage{enumitem} % Pacote para enumerar itens
\usepackage{amssymb} % Pacote para inserir os símbolos dos conjuntos numéricos N, Z, Q, R, C
\usepackage[T1]{fontenc}
\usepackage{enumitem}
\usepackage{textcomp}
\usepackage{pdfpages} % Pacote para inserir página em PDF
\usepackage{courier} % Pacote para usar a fonte courier
%\usepackage{matlab-prettifier}
\usepackage{lscape} % Pacote para mudar sentido da página
\newtheorem{theorem}{Teorema} % Pacote para inserir Teorema
\newtheorem{corollary}{Corollary} % Pacote para inserir Corolário
\newtheorem{remark}{Remark}
\newtheorem{mytheorem}{Definition}
\usepackage{hyperref} % Pacote para inserir hyperref
\usepackage{multirow} % Pacote para usar os comandos \multirow
\usepackage{pdfpages} % Pacote para incluir PDF's
\usepackage{listingsutf8} % Pacote para inserir arquivos externos tipo .txt
\usepackage{etoolbox}
\AtBeginEnvironment{verbatim}{\footnotesize}

\usepackage{standalone} % Pacote para inserir sub-arquivos LaTeX
\usepackage{bookmark} % Pacote para NÃO tratar o Índice como Apêndice e colocar letra

% Nome da listagem em português
\renewcommand{\lstlistingname}{Listagem}

% REFERÊNCIAS: ----------
%\usepackage[numbers, compress]{natbib} % Pacote para usar as referências bibliográficas
%\usepackage{usebib}
% \bibliographystyle{references} % Estilo das referências
%\setcitestyle{authoryear,aysep={},open={(},close={)}}
%\newcommand{\myand}{e}

% ----------


% PARÂMETROS DA FUNÇÃO LISTINGS: ----------
\lstset
{
    basicstyle=\small\ttfamily,
    commentstyle=\color{Green},
    keywordstyle=\color{Cerulean},
    frame=single,
    language=Python,
    morekeywords={True, False},
    numbers=left,
    numbersep=10pt,
    numberstyle=\footnotesize\color{Gray},
    showstringspaces=false,
    stringstyle=\color{Mulberry},
    tabsize=3,
    breaklines=true, % Este comando é para quebrar uma linha comprida
    postbreak=\mbox{\textcolor{red}{$\hookrightarrow$}\space}, % Este comando é para indicar com uma seta vermelha que a linha continua embaixo 
    extendedchars=true,
    literate=
        {á}{{\'a}}1 {Á}{{\'A}}1 {à}{{\`a}}1 {À}{{\`A}}1 {ã}{{\~a}}1 {Ã}{{\~A}}1 {â}{{\^a}}1 {Â}{{\^A}}1
        {ç}{{\c{c}}}1 {Ç}{{\c{C}}}1 {é}{{\'e}}1 {É}{{\'E}}1 {ê}{{\^e}}1 {Ê}{{\^E}}1 {í}{{\'\i}}1 {Í}{{\'I}}1
        {ó}{{\'o}}1 {Ó}{{\'O}}1 {ô}{{\^o}}1 {Ô}{{\^O}}1 {õ}{{\~o}}1 {Õ}{{\~O}}1 {ú}{{\'u}}1 {Ú}{{\'U}}1
        {ü}{{\"u}}1 {Ü}{{\"U}}1 {“}{{``}}1 {”}{{''}}1 {’}{{'}}1 {—}{{--}}1 {–}{{--}}1 {−}{{-}}1
        {±}{{$\pm$}}1 {σ}{{$\sigma$}}1 {→}{{$\to$}}1 {…}{{\ldots}}1
}

\lstset{inputencoding=utf8/latin1}

\lstdefinelanguage{Matlab}{
	language=Matlab,
	basicstyle=\ttfamily\small,
	keywordstyle=\color{blue}\bfseries,
	stringstyle=\color{red},
	commentstyle=\color{gray}\itshape,
	morekeywords={function, end, if, else, for, while, switch, case, otherwise},
	sensitive=true
}
% ----------

% ----------
\usepackage{nomencl} % Pacote para criar nomenclatura e lista de nomenclatura
\makenomenclature % Criar a nomenclatura
\renewcommand{\nomname}{Lista de Nomenclaturas} % Renomear o título da nomenclatura
% ----------

\usepackage{amsmath} % Pacote para inserir matrix com chaves, entre outras coisas

\usepackage{booktabs} % Pacote para usar \toprule, \midrule e \bottomrule


% EXECUTAR CÁLCULOS A PARTIR DA LINHA DE COMANDO: ---------
%\usepackage{siunitx, xfp} % Pacote para executar cálculos a partir da linha de comando
%\usepackage{l3regex}%unless you have very new expl3
%%\sisetup{locale=DE}
%%\let\oldnum\num
%%\ExplSyntaxOn
%\def\num#1{\def\tmp{#1}\regex_replace_all:nnN{,([0-9])}{.\1}{\tmp}\oldnum{\tmp}}
%\ExplSyntaxOff
% ----------


%\pagecolor{black} % Alterar a cor da página
%\color{white} % Alterar a cor do texto

% ----------
% Pacotes para inserir algoritmos e pseudo-códigos:
\setlength{\marginparwidth}{2.5cm}
\usepackage[colorinlistoftodos]{todonotes}
\usepackage{algorithm} % Pacote para inserir algoritmos
\usepackage{algpseudocode} % Pacote para inserir algoritmos e pseudo-códigos
\floatname{algorithm}{Algoritmo} % Renomear a palavra Algorithm para Procedimento
\renewcommand{\algorithmicrequire}{\textbf{entrada:}} % Renomear a palavra Input para Entrada
\renewcommand{\algorithmicensure}{\textbf{saída:}} % Renomear a palavra Output para Saída
\renewcommand{\algorithmicwhile}{\textbf{enquanto}} % Renomear a palavra While para Enquanto
\renewcommand{\algorithmicif}{\textbf{se}} % Renomear a palavra If para Se
\renewcommand{\algorithmicend}{\textbf{fim}} % Renomear a palavra End Output para Fim do
\renewcommand{\algorithmicprocedure}{\textbf{procedimento}} % Renomear a palavra Procedure para Procedimento
\renewcommand{\algorithmicfor}{\textbf{para}} % Renomear a palavra For para Para
\renewcommand{\algorithmicelse}{\textbf{senão}} % Renomear a palavra Else para Senão
\renewcommand{\algorithmicfunction}{\textbf{função}} % Renomear a palavra Function para Função 
\renewcommand{\algorithmicthen}{\textbf{então}} % Renomear a palavra Then para Então
\renewcommand{\algorithmicdo}{\textbf{fazer}} % Renomear a palavra Do para Fazer
% ----------

\newcommand{\quotes}[1]{``#1''} % Comando para inserir um texto dentro de aspas-dupla com espaço após o final

% FLUXOGRAMA: ----------
\usepackage{tikz}
\usetikzlibrary{positioning, shapes.geometric}
\tikzstyle{line}=[draw] % here

\tikzstyle{startstop} = [rectangle, rounded corners, minimum width=3cm,
minimum height=1cm,text centered, draw=black, fill=white!30]
\tikzstyle{arrow} = [thick, ->, >=stealth]
\tikzstyle{process} = [rectangle, minimum width=3cm, minimum height=1cm, text centered, draw=black, fill=white!30]
\tikzstyle{io} = [trapezium, trapezium left angle=70, trapezium right angle=110, minimum width=3cm, minimum height=1cm, text centered, draw=black, fill=white!30]
\tikzstyle{connector} = [draw, circle, fill=white!20]
\tikzstyle{decision} = [diamond, minimum width=3cm, minimum height=1cm, text centered, draw=black, fill=white!30, aspect=3]
% ----------


% CAIXA ENTORNO DAS FÓRMULAS: ----------
\usepackage{dashbox}
\newcommand\dboxed[1]{\dbox{\ensuremath{#1}}}

% ----------

%\usepackage{tabularx} % para colunas com largura ajustável

% -------- CONFIGURAÇÃO DE MARGENS --------
% \usepackage[top=2cm,bottom=2cm,left=3cm,right=2cm]{geometry}

% -------- PACOTE fancyhdr --------
\usepackage{fancyhdr}

\setlength{\marginparwidth}{2cm}

% Ajuste do tamanho do cabeçalho para acompanhar exatamente \textwidth
\setlength{\headwidth}{\textwidth}

% Define um headheight mínimo (para evitar avisos)
\setlength{\headheight}{15pt} 

% Limpa configurações padrão de cabeçalho/rodapé
\pagestyle{fancy}
\fancyhf{} 

% Marcações automáticas de seção para o cabeçalho (opcional)
\renewcommand{\sectionmark}[1]{%
	\markboth{\thesection\ #1}{}%
}

% Texto do cabeçalho à esquerda: seção e título
% Usamos \leftmark pois o \sectionmark acima grava no \leftmark
\fancyhead[L]{\small \leftmark}

% Texto do cabeçalho à direita: número de página
\fancyhead[R]{\small \thepage}

% Linha horizontal no cabeçalho (espessura 0.4pt)
\renewcommand{\headrulewidth}{0.4pt}

% Sobrescreve o título padrão “References” para “Referências”
%\renewcommand{\refname}{Referências}
% Tradução de autoref para o português
\def\figureautorefname{Figura}
\def\tableautorefname{Tabela}
\def\equationautorefname{Equa\c{c}\~ao}

% \hypersetup{backref=true}

\ifdefined\backref
        \renewcommand*{\backref}[1]{}
        \renewcommand*{\backrefentrycount}[2]{#1\ifnum#2>1 ~(#2)\fi}
\fi

\DeclarePairedDelimiter{\ceil}{\lceil}{\rceil}

\onehalfspacing

\usepackage{pgfgantt}

%++++++++++++++++++++++++++++++++++++++++++++++++++++++++++++++++++++++++++++++
% Identificacoes (se o trabalho for em inglês, insira os dados em inglês)
% Para entradas abreviadas de Professora (Profa.) em português escreva: Prof$^\textnormal{a}$.
%++++++++++++++++++++++++++++++++++++++++++++++++++++++++++++++++++++++++++++++
\course{Engenharia Aeronáutica e Mecânica} % Programa de PG ou Curso de Graduação
\area{Propulsão Aeroespacial e Energia} % Área de concentração na PG (Não utilizado no caso de TG)

% Autor do trabalho: Nome Sobrenome
\authorgender{masc} %sexo: masc ou fem
\author{\textbf{EDEN DENIS} Ferreira da Silva Lopes}{Santos}
\itaauthoraddress{Rua João Fonseca dos Santos, 20, Apto. 93}{12.230-088}{São José dos Campos--SP}

% Titulo da Tese/Dissertação
\title{Aplicação de Inteligência Artificial (AI) em Banco de Dados de Emissão de Motores}

% Orientador
\advisorgender{masc} % masc ou fem
\advisor{Prof.~Dr.}{Pedro Teixeira \textbf{LACAVA}}{ITA}

% Coorientador (Caso não haja coorientador, colocar ambas as variáveis \coadvisorgender e \coadvisor comentadas, com um % na frente)
\coadvisorgender{masc}
\coadvisor{}{}  % {Prof.~Dr.}{Carlos Alberto Rocha \textbf{PIMENTEL}}{ITA}

% Pró-reitor da Pós-graduação
\bossgender{fem}
\boss{Prof.~Dra.}{Emília \textbf{VILLANI}}

%Coordenador do curso no caso de TG
\bosscoursegender{masc}
\bosscourse{Prof.~Dr.}{ - } % masc ou fem

% Palavras-Chaves informadas pela Biblioteca -> utilizada na CIP
\kwcip{Engine taxonomy}
\kwcip{Emission classification}
\kwcip{Inteligência Artificial}
\kwcip{Artificial Intelligence}
\kwcip{Python}

% membros da banca examinadora
\examiner{Prof. Dr.}{Pedro Teixeira \textbf{LACAVA}}{President}{ITA}
\examiner{Prof. Dr.}{}{Internal Member}{ITA} % Vinícius \textbf{MALATESTA}}{Internal Member}{ITA}
\examiner{Prof. Dr.}{}{Internal Member}{ITA} % {Elisan dos Santos \textbf{MAGALHÃES}}{Internal Member}{ITA}
\examiner{Prof. Dr.}{}{Internal Member}{ITA} % {Carlos Alberto Rocha \textbf{PIMENTEL}}{External Member}{UFABC}
\examiner{Prof. Dr.}{}{Internal Member}{ITA} % {\textbf{KARL} Peter Burr}{External Member}{UFABC}

% Data da defesa (mês em maiúsculo, se trabalho em inglês, e minúsculo se trabalho em português) 
\date{14}{JULHO}{2025}

% Número CDU - (somente para TG)
\cdu{621.38}

% Folha de Registro do Documento
\newcommand{\advisorTitle}{Prof.~Dr.}
\newcommand{\advisorName}{Pedro Teixeira \textbf{LACAVA}}

\newcommand{\coadvisorTitle}{Prof.~Dr.}
\newcommand{\coadvisorName}{Carlos Alberto Rocha \textbf{PIMENTEL}}

\newcommand{\dateDay}{14}
\newcommand{\dateMonth}{JULHO}
\newcommand{\dateYear}{2025}

% Data de publicação e defesa (usada na folha de registro)
\newcommand{\frdDataPublicacao}{04 de agosto de 2025}
\newcommand{\frdNumeroRegistro}{DCTA/ITA/TD-029/2020}
\newcommand{\frdPalavrasAutor}{Big Data; Engine Emission; Emission Databank; Clustering; Artificial Intelligence}
\newcommand{\frdPalavrasResultado}{Taxonomia; Emissão de Motores; Planejamento de Dados; Inteligência Artificial; Processamento de Dados}
\newcommand{\frdApresentacao}{
	ITA, São José dos Campos. Curso de Mestrado. Programa de Pós-Graduação em Engenharia Aeronáutica e Mecânica. Área de Propulsão Aeroespacial e Energia.\\
	Orientador: \advisorTitle~\advisorName.\\
	Coorientador: \coadvisorTitle~\coadvisorName.\\
	Defesa em \dateDay/\dateMonth/\dateYear. Publicada em \frdDataPublicacao.
}

\newcommand{\PorcentagemDaLargura}{0.80}

\begin{document}
% \chapter{Resumo}
	
%	\newpage
	\pagestyle{fancy}
	\thispagestyle{empty}
	%\shipout\null % Inserir página em branco
	\begin{abstract}
		\noindent
		
                Este trabalho apresenta uma metodologia baseada em dados para classificar motores de aeronaves com base em características de projeto e termodinâmicas relacionadas à eficiência e emissões. A abordagem proposta visa identificar similaridades entre diferentes configurações de motores utilizando técnicas de aprendizado de máquina não supervisionado, permitindo uma comparação sistemática de tecnologias de propulsão e apoiando estudos de avaliação tecnológica. A análise é baseada em dados de motores disponíveis publicamente, fornecidos pela Agência Europeia para a Segurança da Aviação (EASA). Um conjunto de parâmetros-chave de projeto e desempenho, incluindo a taxa de compressão total, a temperatura de entrada da turbina, a taxa de derivação e outras variáveis ??relevantes, é selecionado para descrever cada motor em um espaço de características multidimensional. Antes do agrupamento, o conjunto de dados é pré-processado por meio de normalização, verificações de consistência e tratamento de entradas incompletas, a fim de garantir robustez e comparabilidade entre as amostras. O algoritmo de agrupamento K-means é aplicado para agrupar motores com características semelhantes. Diversas configurações de agrupamento são avaliadas variando-se o número de clusters e aplicando-se métricas de validação padrão para identificar uma partição adequada do conjunto de dados. Os agrupamentos resultantes são examinados para identificar padrões tecnológicos, distinguir diferentes filosofias de projeto e destacar configurações de motores que apresentam eficiência relativamente maior ou características relacionadas a emissões mais favoráveis ??dentro de seus respectivos grupos. Os resultados do agrupamento indicam que motores com características termodinâmicas e de projeto comparáveis ??tendem a formar grupos bem definidos, refletindo similaridades no nível de desempenho e maturidade tecnológica. Essa classificação permite uma comparação mais clara entre famílias de motores e auxilia na identificação de configurações representativas dentro de cada agrupamento. Além disso, a análise fornece informações sobre a importância relativa das variáveis ??selecionadas no processo de agrupamento, contribuindo para uma melhor compreensão de como os parâmetros de projeto influenciam o agrupamento de motores. A metodologia proposta não se baseia em classificações predefinidas e permite que os padrões emerjam diretamente dos dados. Dessa forma, ela fornece uma estrutura flexível e reproduzível para a análise de tecnologias de motores aeronáuticos. Os resultados podem apoiar estudos futuros sobre avaliação de emissões, benchmarking de motores e a seleção de parâmetros representativos para o desenvolvimento de sistemas de propulsão aeronáutica mais eficientes e ambientalmente sustentáveis.
		
	\end{abstract}
	
    \textbf{Palavras-chave:} taxonomia de motores de aeronaves, classificação de motores para pré-projeto, Big Data de emissões, dendrograma, K-means, Python.
	
\end{document}
  % [x] Revisão do texto
	
	% Abstract
	\documentclass[12pt]{article} % Tamanho de letra 12, papel A4, Estilo de artigo

\usepackage[utf8]{inputenc} % Pacote para permitir o uso de acentos
% Ajustes de hifenização
\tolerance=1
\emergencystretch=\maxdimen
\hyphenpenalty=10000
\hbadness=10000
\usepackage{lipsum} % Pacote para gerar texto aleatório
\usepackage{ragged2e} % Pacote para mais opções de alinhamento
\usepackage{graphicx} % Pacote para inserir figuras jpeg, pdf e png
\usepackage{ae}
\usepackage{epsfig}
\usepackage{mathtools}
\usepackage[export]{adjustbox}
\usepackage{lineno}
\usepackage{amsthm}
\usepackage{commath}
\usepackage{color, colortbl}
\usepackage{bbm}
\usepackage{setspace} % Pacote para ajustar espaçamento entre linhas

% Instructions to prevent LaTeX from hyphenating words
\tolerance=1
\emergencystretch=\maxdimen
\hyphenpenalty=10000
\hbadness=10000

\graphicspath{{figures/}} % Indica o diretório no qual as imagens estão sendo salvas, evitando que precisamos digitar parte do caminho da pasta que contém as imagens
\usepackage[rightcaption]{sidecap} % Pacote para mais opções de figura, colocar, título da legenda da figura ao lado direito
\usepackage{url} % Pacote para citar referências de sites
\usepackage{makeidx} % Pacote para criar índice
\makeindex % Habilitar a criação do índice
\usepackage{longtable} % Pacote para usar tabelas longas

% ----------
% Pacotes para dar quebra de linha dentro da tabela:
%\usepackage{fourier} 
%\usepackage{array}
\usepackage{makecell}

\renewcommand\theadalign{bc}
\renewcommand\theadfont{\bfseries}
\renewcommand\theadgape{\Gape[4pt]}
\renewcommand\cellgape{\Gape[4pt]}
% ----------

\usepackage{wrapfig} % Pacote para alinhar figuras à esquerda ou à direita do texto
\usepackage{ae}
\usepackage{epsfig}
\usepackage{mathtools}
\usepackage[export]{adjustbox}
\usepackage{lineno}
\usepackage{amsfonts}
\usepackage{amsthm}
\usepackage{commath}
\usepackage{color, colortbl}
\usepackage{bbm}

% ----------
% Pacote para inserir cabeçalho e rodapé:
% Atenção! Se remover o cabeçalho e rodapé, o LaTeX coloca como cabeçalho o padrão dele, que é bem interessante também!
% \usepackage{fancyhdr} % Pacote para inserir cabeçalho e rodapé
% \pagestyle{fancy} % Estilo do cabeçalho e rodapé, neste caso apenas uma linha no cabeçalho
%\lhead{{\footnotesize Prandiano\textregistered}} % Inserir texto no canto esquerdo do cabeçalho
%\chead{centro} % Inserir texto no centro do cabeçalho
%\rhead{{\footnotesize Curso de Python}} % Inserir texto no canto direito do cabeçalho
%\lfoot{{\footnotesize Copyright \textcopyright 2018 Prandiano}} % Inserir texto no canto esquerdo do rodapé
%\cfoot{centro} % Inserir texto no centro do rodapé
%\rfoot{{\footnotesize Reprodução proibida: Lei 9.610 de 19/02/1998}} % Inserir texto no canto direito do rodapé
% O comando \renewcommand re-nova o parâmetro especificado
%\renewcommand{\footrulewidth}{0.4pt} % Inserir uma linha acima do rodapé
% ----------

% ----------
% Pacote para inserir múltiplas figuras:
\usepackage{subfigure} % Pacote para inserir múltiplas figuras
% \usepackage{subcaption} % Pacote incompatível com subfigure
%\usepackage{caption}
% ----------

\usepackage{multicol} % Pacote que permite usar multicolunas
\usepackage{float} % Pacote para usar o "H" do figure
\usepackage[dvipsnames]{xcolor} % Pacote de cores para textos e possivelmente outros
\usepackage{amsmath} % Pacote para usar o cases
\usepackage[makeroom]{cancel} % Pacote para usar o cancelar (em e embaixo)
\usepackage{enumitem} % Pacote para enumerar itens
\usepackage{amssymb} % Pacote para inserir os símbolos dos conjuntos numéricos N, Z, Q, R, C
\usepackage[T1]{fontenc}
\usepackage{enumitem}
\usepackage{textcomp}
\usepackage{pdfpages} % Pacote para inserir página em PDF
\usepackage{courier} % Pacote para usar a fonte courier
%\usepackage{matlab-prettifier}
\usepackage{lscape} % Pacote para mudar sentido da página
\newtheorem{theorem}{Teorema} % Pacote para inserir Teorema
\newtheorem{corollary}{Corollary} % Pacote para inserir Corolário
\newtheorem{remark}{Remark}
\newtheorem{mytheorem}{Definition}
\usepackage{hyperref} % Pacote para inserir hyperref
\usepackage{multirow} % Pacote para usar os comandos \multirow
\usepackage{pdfpages} % Pacote para incluir PDF's
\usepackage{listingsutf8} % Pacote para inserir arquivos externos tipo .txt
\usepackage{etoolbox}
\AtBeginEnvironment{verbatim}{\footnotesize}

\usepackage{standalone} % Pacote para inserir sub-arquivos LaTeX
\usepackage{bookmark} % Pacote para NÃO tratar o Índice como Apêndice e colocar letra

% Nome da listagem em português
\renewcommand{\lstlistingname}{Listagem}

% REFERÊNCIAS: ----------
%\usepackage[numbers, compress]{natbib} % Pacote para usar as referências bibliográficas
%\usepackage{usebib}
% \bibliographystyle{references} % Estilo das referências
%\setcitestyle{authoryear,aysep={},open={(},close={)}}
%\newcommand{\myand}{e}

% ----------


% PARÂMETROS DA FUNÇÃO LISTINGS: ----------
\lstset
{
    basicstyle=\small\ttfamily,
    commentstyle=\color{Green},
    keywordstyle=\color{Cerulean},
    frame=single,
    language=Python,
    morekeywords={True, False},
    numbers=left,
    numbersep=10pt,
    numberstyle=\footnotesize\color{Gray},
    showstringspaces=false,
    stringstyle=\color{Mulberry},
    tabsize=3,
    breaklines=true, % Este comando é para quebrar uma linha comprida
    postbreak=\mbox{\textcolor{red}{$\hookrightarrow$}\space}, % Este comando é para indicar com uma seta vermelha que a linha continua embaixo 
    extendedchars=true,
    literate=
        {á}{{\'a}}1 {Á}{{\'A}}1 {à}{{\`a}}1 {À}{{\`A}}1 {ã}{{\~a}}1 {Ã}{{\~A}}1 {â}{{\^a}}1 {Â}{{\^A}}1
        {ç}{{\c{c}}}1 {Ç}{{\c{C}}}1 {é}{{\'e}}1 {É}{{\'E}}1 {ê}{{\^e}}1 {Ê}{{\^E}}1 {í}{{\'\i}}1 {Í}{{\'I}}1
        {ó}{{\'o}}1 {Ó}{{\'O}}1 {ô}{{\^o}}1 {Ô}{{\^O}}1 {õ}{{\~o}}1 {Õ}{{\~O}}1 {ú}{{\'u}}1 {Ú}{{\'U}}1
        {ü}{{\"u}}1 {Ü}{{\"U}}1 {“}{{``}}1 {”}{{''}}1 {’}{{'}}1 {—}{{--}}1 {–}{{--}}1 {−}{{-}}1
        {±}{{$\pm$}}1 {σ}{{$\sigma$}}1 {→}{{$\to$}}1 {…}{{\ldots}}1
}

\lstset{inputencoding=utf8/latin1}

\lstdefinelanguage{Matlab}{
	language=Matlab,
	basicstyle=\ttfamily\small,
	keywordstyle=\color{blue}\bfseries,
	stringstyle=\color{red},
	commentstyle=\color{gray}\itshape,
	morekeywords={function, end, if, else, for, while, switch, case, otherwise},
	sensitive=true
}
% ----------

% ----------
\usepackage{nomencl} % Pacote para criar nomenclatura e lista de nomenclatura
\makenomenclature % Criar a nomenclatura
\renewcommand{\nomname}{Lista de Nomenclaturas} % Renomear o título da nomenclatura
% ----------

\usepackage{amsmath} % Pacote para inserir matrix com chaves, entre outras coisas

\usepackage{booktabs} % Pacote para usar \toprule, \midrule e \bottomrule


% EXECUTAR CÁLCULOS A PARTIR DA LINHA DE COMANDO: ---------
%\usepackage{siunitx, xfp} % Pacote para executar cálculos a partir da linha de comando
%\usepackage{l3regex}%unless you have very new expl3
%%\sisetup{locale=DE}
%%\let\oldnum\num
%%\ExplSyntaxOn
%\def\num#1{\def\tmp{#1}\regex_replace_all:nnN{,([0-9])}{.\1}{\tmp}\oldnum{\tmp}}
%\ExplSyntaxOff
% ----------


%\pagecolor{black} % Alterar a cor da página
%\color{white} % Alterar a cor do texto

% ----------
% Pacotes para inserir algoritmos e pseudo-códigos:
\setlength{\marginparwidth}{2.5cm}
\usepackage[colorinlistoftodos]{todonotes}
\usepackage{algorithm} % Pacote para inserir algoritmos
\usepackage{algpseudocode} % Pacote para inserir algoritmos e pseudo-códigos
\floatname{algorithm}{Algoritmo} % Renomear a palavra Algorithm para Procedimento
\renewcommand{\algorithmicrequire}{\textbf{entrada:}} % Renomear a palavra Input para Entrada
\renewcommand{\algorithmicensure}{\textbf{saída:}} % Renomear a palavra Output para Saída
\renewcommand{\algorithmicwhile}{\textbf{enquanto}} % Renomear a palavra While para Enquanto
\renewcommand{\algorithmicif}{\textbf{se}} % Renomear a palavra If para Se
\renewcommand{\algorithmicend}{\textbf{fim}} % Renomear a palavra End Output para Fim do
\renewcommand{\algorithmicprocedure}{\textbf{procedimento}} % Renomear a palavra Procedure para Procedimento
\renewcommand{\algorithmicfor}{\textbf{para}} % Renomear a palavra For para Para
\renewcommand{\algorithmicelse}{\textbf{senão}} % Renomear a palavra Else para Senão
\renewcommand{\algorithmicfunction}{\textbf{função}} % Renomear a palavra Function para Função 
\renewcommand{\algorithmicthen}{\textbf{então}} % Renomear a palavra Then para Então
\renewcommand{\algorithmicdo}{\textbf{fazer}} % Renomear a palavra Do para Fazer
% ----------

\newcommand{\quotes}[1]{``#1''} % Comando para inserir um texto dentro de aspas-dupla com espaço após o final

% FLUXOGRAMA: ----------
\usepackage{tikz}
\usetikzlibrary{positioning, shapes.geometric}
\tikzstyle{line}=[draw] % here

\tikzstyle{startstop} = [rectangle, rounded corners, minimum width=3cm,
minimum height=1cm,text centered, draw=black, fill=white!30]
\tikzstyle{arrow} = [thick, ->, >=stealth]
\tikzstyle{process} = [rectangle, minimum width=3cm, minimum height=1cm, text centered, draw=black, fill=white!30]
\tikzstyle{io} = [trapezium, trapezium left angle=70, trapezium right angle=110, minimum width=3cm, minimum height=1cm, text centered, draw=black, fill=white!30]
\tikzstyle{connector} = [draw, circle, fill=white!20]
\tikzstyle{decision} = [diamond, minimum width=3cm, minimum height=1cm, text centered, draw=black, fill=white!30, aspect=3]
% ----------


% CAIXA ENTORNO DAS FÓRMULAS: ----------
\usepackage{dashbox}
\newcommand\dboxed[1]{\dbox{\ensuremath{#1}}}

% ----------

%\usepackage{tabularx} % para colunas com largura ajustável

% -------- CONFIGURAÇÃO DE MARGENS --------
% \usepackage[top=2cm,bottom=2cm,left=3cm,right=2cm]{geometry}

% -------- PACOTE fancyhdr --------
\usepackage{fancyhdr}

\setlength{\marginparwidth}{2cm}

% Ajuste do tamanho do cabeçalho para acompanhar exatamente \textwidth
\setlength{\headwidth}{\textwidth}

% Define um headheight mínimo (para evitar avisos)
\setlength{\headheight}{15pt} 

% Limpa configurações padrão de cabeçalho/rodapé
\pagestyle{fancy}
\fancyhf{} 

% Marcações automáticas de seção para o cabeçalho (opcional)
\renewcommand{\sectionmark}[1]{%
	\markboth{\thesection\ #1}{}%
}

% Texto do cabeçalho à esquerda: seção e título
% Usamos \leftmark pois o \sectionmark acima grava no \leftmark
\fancyhead[L]{\small \leftmark}

% Texto do cabeçalho à direita: número de página
\fancyhead[R]{\small \thepage}

% Linha horizontal no cabeçalho (espessura 0.4pt)
\renewcommand{\headrulewidth}{0.4pt}

% Sobrescreve o título padrão “References” para “Referências”
%\renewcommand{\refname}{Referências}
% Tradução de autoref para o português
\def\figureautorefname{Figura}
\def\tableautorefname{Tabela}
\def\equationautorefname{Equa\c{c}\~ao}

% \hypersetup{backref=true}

\ifdefined\backref
        \renewcommand*{\backref}[1]{}
        \renewcommand*{\backrefentrycount}[2]{#1\ifnum#2>1 ~(#2)\fi}
\fi

\DeclarePairedDelimiter{\ceil}{\lceil}{\rceil}

\onehalfspacing

\usepackage{pgfgantt}

%++++++++++++++++++++++++++++++++++++++++++++++++++++++++++++++++++++++++++++++
% Identificacoes (se o trabalho for em inglês, insira os dados em inglês)
% Para entradas abreviadas de Professora (Profa.) em português escreva: Prof$^\textnormal{a}$.
%++++++++++++++++++++++++++++++++++++++++++++++++++++++++++++++++++++++++++++++
\course{Engenharia Aeronáutica e Mecânica} % Programa de PG ou Curso de Graduação
\area{Propulsão Aeroespacial e Energia} % Área de concentração na PG (Não utilizado no caso de TG)

% Autor do trabalho: Nome Sobrenome
\authorgender{masc} %sexo: masc ou fem
\author{\textbf{EDEN DENIS} Ferreira da Silva Lopes}{Santos}
\itaauthoraddress{Rua João Fonseca dos Santos, 20, Apto. 93}{12.230-088}{São José dos Campos--SP}

% Titulo da Tese/Dissertação
\title{Aplicação de Inteligência Artificial (AI) em Banco de Dados de Emissão de Motores}

% Orientador
\advisorgender{masc} % masc ou fem
\advisor{Prof.~Dr.}{Pedro Teixeira \textbf{LACAVA}}{ITA}

% Coorientador (Caso não haja coorientador, colocar ambas as variáveis \coadvisorgender e \coadvisor comentadas, com um % na frente)
\coadvisorgender{masc}
\coadvisor{}{}  % {Prof.~Dr.}{Carlos Alberto Rocha \textbf{PIMENTEL}}{ITA}

% Pró-reitor da Pós-graduação
\bossgender{fem}
\boss{Prof.~Dra.}{Emília \textbf{VILLANI}}

%Coordenador do curso no caso de TG
\bosscoursegender{masc}
\bosscourse{Prof.~Dr.}{ - } % masc ou fem

% Palavras-Chaves informadas pela Biblioteca -> utilizada na CIP
\kwcip{Engine taxonomy}
\kwcip{Emission classification}
\kwcip{Inteligência Artificial}
\kwcip{Artificial Intelligence}
\kwcip{Python}

% membros da banca examinadora
\examiner{Prof. Dr.}{Pedro Teixeira \textbf{LACAVA}}{President}{ITA}
\examiner{Prof. Dr.}{}{Internal Member}{ITA} % Vinícius \textbf{MALATESTA}}{Internal Member}{ITA}
\examiner{Prof. Dr.}{}{Internal Member}{ITA} % {Elisan dos Santos \textbf{MAGALHÃES}}{Internal Member}{ITA}
\examiner{Prof. Dr.}{}{Internal Member}{ITA} % {Carlos Alberto Rocha \textbf{PIMENTEL}}{External Member}{UFABC}
\examiner{Prof. Dr.}{}{Internal Member}{ITA} % {\textbf{KARL} Peter Burr}{External Member}{UFABC}

% Data da defesa (mês em maiúsculo, se trabalho em inglês, e minúsculo se trabalho em português) 
\date{14}{JULHO}{2025}

% Número CDU - (somente para TG)
\cdu{621.38}

% Folha de Registro do Documento
\newcommand{\advisorTitle}{Prof.~Dr.}
\newcommand{\advisorName}{Pedro Teixeira \textbf{LACAVA}}

\newcommand{\coadvisorTitle}{Prof.~Dr.}
\newcommand{\coadvisorName}{Carlos Alberto Rocha \textbf{PIMENTEL}}

\newcommand{\dateDay}{14}
\newcommand{\dateMonth}{JULHO}
\newcommand{\dateYear}{2025}

% Data de publicação e defesa (usada na folha de registro)
\newcommand{\frdDataPublicacao}{04 de agosto de 2025}
\newcommand{\frdNumeroRegistro}{DCTA/ITA/TD-029/2020}
\newcommand{\frdPalavrasAutor}{Big Data; Engine Emission; Emission Databank; Clustering; Artificial Intelligence}
\newcommand{\frdPalavrasResultado}{Taxonomia; Emissão de Motores; Planejamento de Dados; Inteligência Artificial; Processamento de Dados}
\newcommand{\frdApresentacao}{
	ITA, São José dos Campos. Curso de Mestrado. Programa de Pós-Graduação em Engenharia Aeronáutica e Mecânica. Área de Propulsão Aeroespacial e Energia.\\
	Orientador: \advisorTitle~\advisorName.\\
	Coorientador: \coadvisorTitle~\coadvisorName.\\
	Defesa em \dateDay/\dateMonth/\dateYear. Publicada em \frdDataPublicacao.
}

\newcommand{\PorcentagemDaLargura}{0.80}

\begin{document}
% \chapter{Abstract}
	
	\begin{englishabstract}
		\noindent
		
                This work presents a data-driven methodology to classify aircraft engines based on design and thermodynamic characteristics related to efficiency and emissions. The proposed approach aims to identify similarities among different engine configurations using unsupervised machine learning techniques, enabling a systematic comparison of propulsion technologies and supporting technology assessment studies. The analysis is based on publicly available engine data provided by the European Union Aviation Safety Agency (EASA). A set of key design and performance parameters, including overall pressure ratio, turbine inlet temperature, bypass ratio, and other relevant variables, is selected to describe each engine in a multidimensional feature space. Prior to clustering, the dataset is preprocessed through normalization, consistency checks, and the treatment of incomplete entries in order to ensure robustness and comparability among samples. The K-means clustering algorithm is applied to group engines with similar characteristics. Several clustering configurations are evaluated by varying the number of clusters and applying standard validation metrics to identify an adequate partition of the dataset. The resulting clusters are examined to identify technological patterns, distinguish different design philosophies, and highlight engine configurations that exhibit relatively higher efficiency or more favorable emission-related characteristics within their respective groups. The clustering results indicate that engines with comparable thermodynamic and design characteristics tend to form well-defined groups, reflecting similarities in performance level and technological maturity. This classification allows a clearer comparison between engine families and supports the identification of representative configurations within each cluster. In addition, the analysis provides insights into the relative importance of the selected variables in the clustering process, contributing to a better understanding of how design parameters influence engine grouping. The proposed methodology does not rely on predefined classifications and allows patterns to emerge directly from the data. As such, it provides a flexible and reproducible framework for the analysis of aircraft engine technologies. The results may support future studies on emissions assessment, engine benchmarking, and the selection of representative parameters for the development of more efficient and environmentally sustainable aircraft propulsion systems.
	\end{englishabstract}
	
    \textbf{Keywords:} aircraft engine taxonomy, emission classification, Big Data, dendrogram, K-means, Python.
	
\end{document}
  % [x] Revisão do texto
	
	% Listas:
	\glsaddall

	% Lista de abreviaturas
	\printnoidxglossary[type=\acronymtype,
	    title=\listofabbreviationsname,
	    sort=standard,
	    nonumberlist]
	\pdfbookmark[0]{\listofabbreviationsname}{listaabreviaturas}

	% Lista de simbolos
	\printnoidxglossary[
	    title=\listofsymbolsname,
	    sort=use,
	    nonumberlist]
	\pdfbookmark[0]{\listofsymbolsname}{listasimbolos}

	\documentclass[12pt]{article} % Tamanho de letra 12, papel A4, Estilo de artigo

\usepackage[utf8]{inputenc} % Pacote para permitir o uso de acentos
% Ajustes de hifenização
\tolerance=1
\emergencystretch=\maxdimen
\hyphenpenalty=10000
\hbadness=10000
\usepackage{lipsum} % Pacote para gerar texto aleatório
\usepackage{ragged2e} % Pacote para mais opções de alinhamento
\usepackage{graphicx} % Pacote para inserir figuras jpeg, pdf e png
\usepackage{ae}
\usepackage{epsfig}
\usepackage{mathtools}
\usepackage[export]{adjustbox}
\usepackage{lineno}
\usepackage{amsthm}
\usepackage{commath}
\usepackage{color, colortbl}
\usepackage{bbm}
\usepackage{setspace} % Pacote para ajustar espaçamento entre linhas

% Instructions to prevent LaTeX from hyphenating words
\tolerance=1
\emergencystretch=\maxdimen
\hyphenpenalty=10000
\hbadness=10000

\graphicspath{{figures/}} % Indica o diretório no qual as imagens estão sendo salvas, evitando que precisamos digitar parte do caminho da pasta que contém as imagens
\usepackage[rightcaption]{sidecap} % Pacote para mais opções de figura, colocar, título da legenda da figura ao lado direito
\usepackage{url} % Pacote para citar referências de sites
\usepackage{makeidx} % Pacote para criar índice
\makeindex % Habilitar a criação do índice
\usepackage{longtable} % Pacote para usar tabelas longas

% ----------
% Pacotes para dar quebra de linha dentro da tabela:
%\usepackage{fourier} 
%\usepackage{array}
\usepackage{makecell}

\renewcommand\theadalign{bc}
\renewcommand\theadfont{\bfseries}
\renewcommand\theadgape{\Gape[4pt]}
\renewcommand\cellgape{\Gape[4pt]}
% ----------

\usepackage{wrapfig} % Pacote para alinhar figuras à esquerda ou à direita do texto
\usepackage{ae}
\usepackage{epsfig}
\usepackage{mathtools}
\usepackage[export]{adjustbox}
\usepackage{lineno}
\usepackage{amsfonts}
\usepackage{amsthm}
\usepackage{commath}
\usepackage{color, colortbl}
\usepackage{bbm}

% ----------
% Pacote para inserir cabeçalho e rodapé:
% Atenção! Se remover o cabeçalho e rodapé, o LaTeX coloca como cabeçalho o padrão dele, que é bem interessante também!
% \usepackage{fancyhdr} % Pacote para inserir cabeçalho e rodapé
% \pagestyle{fancy} % Estilo do cabeçalho e rodapé, neste caso apenas uma linha no cabeçalho
%\lhead{{\footnotesize Prandiano\textregistered}} % Inserir texto no canto esquerdo do cabeçalho
%\chead{centro} % Inserir texto no centro do cabeçalho
%\rhead{{\footnotesize Curso de Python}} % Inserir texto no canto direito do cabeçalho
%\lfoot{{\footnotesize Copyright \textcopyright 2018 Prandiano}} % Inserir texto no canto esquerdo do rodapé
%\cfoot{centro} % Inserir texto no centro do rodapé
%\rfoot{{\footnotesize Reprodução proibida: Lei 9.610 de 19/02/1998}} % Inserir texto no canto direito do rodapé
% O comando \renewcommand re-nova o parâmetro especificado
%\renewcommand{\footrulewidth}{0.4pt} % Inserir uma linha acima do rodapé
% ----------

% ----------
% Pacote para inserir múltiplas figuras:
\usepackage{subfigure} % Pacote para inserir múltiplas figuras
% \usepackage{subcaption} % Pacote incompatível com subfigure
%\usepackage{caption}
% ----------

\usepackage{multicol} % Pacote que permite usar multicolunas
\usepackage{float} % Pacote para usar o "H" do figure
\usepackage[dvipsnames]{xcolor} % Pacote de cores para textos e possivelmente outros
\usepackage{amsmath} % Pacote para usar o cases
\usepackage[makeroom]{cancel} % Pacote para usar o cancelar (em e embaixo)
\usepackage{enumitem} % Pacote para enumerar itens
\usepackage{amssymb} % Pacote para inserir os símbolos dos conjuntos numéricos N, Z, Q, R, C
\usepackage[T1]{fontenc}
\usepackage{enumitem}
\usepackage{textcomp}
\usepackage{pdfpages} % Pacote para inserir página em PDF
\usepackage{courier} % Pacote para usar a fonte courier
%\usepackage{matlab-prettifier}
\usepackage{lscape} % Pacote para mudar sentido da página
\newtheorem{theorem}{Teorema} % Pacote para inserir Teorema
\newtheorem{corollary}{Corollary} % Pacote para inserir Corolário
\newtheorem{remark}{Remark}
\newtheorem{mytheorem}{Definition}
\usepackage{hyperref} % Pacote para inserir hyperref
\usepackage{multirow} % Pacote para usar os comandos \multirow
\usepackage{pdfpages} % Pacote para incluir PDF's
\usepackage{listingsutf8} % Pacote para inserir arquivos externos tipo .txt
\usepackage{etoolbox}
\AtBeginEnvironment{verbatim}{\footnotesize}

\usepackage{standalone} % Pacote para inserir sub-arquivos LaTeX
\usepackage{bookmark} % Pacote para NÃO tratar o Índice como Apêndice e colocar letra

% Nome da listagem em português
\renewcommand{\lstlistingname}{Listagem}

% REFERÊNCIAS: ----------
%\usepackage[numbers, compress]{natbib} % Pacote para usar as referências bibliográficas
%\usepackage{usebib}
% \bibliographystyle{references} % Estilo das referências
%\setcitestyle{authoryear,aysep={},open={(},close={)}}
%\newcommand{\myand}{e}

% ----------


% PARÂMETROS DA FUNÇÃO LISTINGS: ----------
\lstset
{
    basicstyle=\small\ttfamily,
    commentstyle=\color{Green},
    keywordstyle=\color{Cerulean},
    frame=single,
    language=Python,
    morekeywords={True, False},
    numbers=left,
    numbersep=10pt,
    numberstyle=\footnotesize\color{Gray},
    showstringspaces=false,
    stringstyle=\color{Mulberry},
    tabsize=3,
    breaklines=true, % Este comando é para quebrar uma linha comprida
    postbreak=\mbox{\textcolor{red}{$\hookrightarrow$}\space}, % Este comando é para indicar com uma seta vermelha que a linha continua embaixo 
    extendedchars=true,
    literate=
        {á}{{\'a}}1 {Á}{{\'A}}1 {à}{{\`a}}1 {À}{{\`A}}1 {ã}{{\~a}}1 {Ã}{{\~A}}1 {â}{{\^a}}1 {Â}{{\^A}}1
        {ç}{{\c{c}}}1 {Ç}{{\c{C}}}1 {é}{{\'e}}1 {É}{{\'E}}1 {ê}{{\^e}}1 {Ê}{{\^E}}1 {í}{{\'\i}}1 {Í}{{\'I}}1
        {ó}{{\'o}}1 {Ó}{{\'O}}1 {ô}{{\^o}}1 {Ô}{{\^O}}1 {õ}{{\~o}}1 {Õ}{{\~O}}1 {ú}{{\'u}}1 {Ú}{{\'U}}1
        {ü}{{\"u}}1 {Ü}{{\"U}}1 {“}{{``}}1 {”}{{''}}1 {’}{{'}}1 {—}{{--}}1 {–}{{--}}1 {−}{{-}}1
        {±}{{$\pm$}}1 {σ}{{$\sigma$}}1 {→}{{$\to$}}1 {…}{{\ldots}}1
}

\lstset{inputencoding=utf8/latin1}

\lstdefinelanguage{Matlab}{
	language=Matlab,
	basicstyle=\ttfamily\small,
	keywordstyle=\color{blue}\bfseries,
	stringstyle=\color{red},
	commentstyle=\color{gray}\itshape,
	morekeywords={function, end, if, else, for, while, switch, case, otherwise},
	sensitive=true
}
% ----------

% ----------
\usepackage{nomencl} % Pacote para criar nomenclatura e lista de nomenclatura
\makenomenclature % Criar a nomenclatura
\renewcommand{\nomname}{Lista de Nomenclaturas} % Renomear o título da nomenclatura
% ----------

\usepackage{amsmath} % Pacote para inserir matrix com chaves, entre outras coisas

\usepackage{booktabs} % Pacote para usar \toprule, \midrule e \bottomrule


% EXECUTAR CÁLCULOS A PARTIR DA LINHA DE COMANDO: ---------
%\usepackage{siunitx, xfp} % Pacote para executar cálculos a partir da linha de comando
%\usepackage{l3regex}%unless you have very new expl3
%%\sisetup{locale=DE}
%%\let\oldnum\num
%%\ExplSyntaxOn
%\def\num#1{\def\tmp{#1}\regex_replace_all:nnN{,([0-9])}{.\1}{\tmp}\oldnum{\tmp}}
%\ExplSyntaxOff
% ----------


%\pagecolor{black} % Alterar a cor da página
%\color{white} % Alterar a cor do texto

% ----------
% Pacotes para inserir algoritmos e pseudo-códigos:
\setlength{\marginparwidth}{2.5cm}
\usepackage[colorinlistoftodos]{todonotes}
\usepackage{algorithm} % Pacote para inserir algoritmos
\usepackage{algpseudocode} % Pacote para inserir algoritmos e pseudo-códigos
\floatname{algorithm}{Algoritmo} % Renomear a palavra Algorithm para Procedimento
\renewcommand{\algorithmicrequire}{\textbf{entrada:}} % Renomear a palavra Input para Entrada
\renewcommand{\algorithmicensure}{\textbf{saída:}} % Renomear a palavra Output para Saída
\renewcommand{\algorithmicwhile}{\textbf{enquanto}} % Renomear a palavra While para Enquanto
\renewcommand{\algorithmicif}{\textbf{se}} % Renomear a palavra If para Se
\renewcommand{\algorithmicend}{\textbf{fim}} % Renomear a palavra End Output para Fim do
\renewcommand{\algorithmicprocedure}{\textbf{procedimento}} % Renomear a palavra Procedure para Procedimento
\renewcommand{\algorithmicfor}{\textbf{para}} % Renomear a palavra For para Para
\renewcommand{\algorithmicelse}{\textbf{senão}} % Renomear a palavra Else para Senão
\renewcommand{\algorithmicfunction}{\textbf{função}} % Renomear a palavra Function para Função 
\renewcommand{\algorithmicthen}{\textbf{então}} % Renomear a palavra Then para Então
\renewcommand{\algorithmicdo}{\textbf{fazer}} % Renomear a palavra Do para Fazer
% ----------

\newcommand{\quotes}[1]{``#1''} % Comando para inserir um texto dentro de aspas-dupla com espaço após o final

% FLUXOGRAMA: ----------
\usepackage{tikz}
\usetikzlibrary{positioning, shapes.geometric}
\tikzstyle{line}=[draw] % here

\tikzstyle{startstop} = [rectangle, rounded corners, minimum width=3cm,
minimum height=1cm,text centered, draw=black, fill=white!30]
\tikzstyle{arrow} = [thick, ->, >=stealth]
\tikzstyle{process} = [rectangle, minimum width=3cm, minimum height=1cm, text centered, draw=black, fill=white!30]
\tikzstyle{io} = [trapezium, trapezium left angle=70, trapezium right angle=110, minimum width=3cm, minimum height=1cm, text centered, draw=black, fill=white!30]
\tikzstyle{connector} = [draw, circle, fill=white!20]
\tikzstyle{decision} = [diamond, minimum width=3cm, minimum height=1cm, text centered, draw=black, fill=white!30, aspect=3]
% ----------


% CAIXA ENTORNO DAS FÓRMULAS: ----------
\usepackage{dashbox}
\newcommand\dboxed[1]{\dbox{\ensuremath{#1}}}

% ----------

%\usepackage{tabularx} % para colunas com largura ajustável

% -------- CONFIGURAÇÃO DE MARGENS --------
% \usepackage[top=2cm,bottom=2cm,left=3cm,right=2cm]{geometry}

% -------- PACOTE fancyhdr --------
\usepackage{fancyhdr}

\setlength{\marginparwidth}{2cm}

% Ajuste do tamanho do cabeçalho para acompanhar exatamente \textwidth
\setlength{\headwidth}{\textwidth}

% Define um headheight mínimo (para evitar avisos)
\setlength{\headheight}{15pt} 

% Limpa configurações padrão de cabeçalho/rodapé
\pagestyle{fancy}
\fancyhf{} 

% Marcações automáticas de seção para o cabeçalho (opcional)
\renewcommand{\sectionmark}[1]{%
	\markboth{\thesection\ #1}{}%
}

% Texto do cabeçalho à esquerda: seção e título
% Usamos \leftmark pois o \sectionmark acima grava no \leftmark
\fancyhead[L]{\small \leftmark}

% Texto do cabeçalho à direita: número de página
\fancyhead[R]{\small \thepage}

% Linha horizontal no cabeçalho (espessura 0.4pt)
\renewcommand{\headrulewidth}{0.4pt}

% Sobrescreve o título padrão “References” para “Referências”
%\renewcommand{\refname}{Referências}
% Tradução de autoref para o português
\def\figureautorefname{Figura}
\def\tableautorefname{Tabela}
\def\equationautorefname{Equa\c{c}\~ao}

% \hypersetup{backref=true}

\ifdefined\backref
        \renewcommand*{\backref}[1]{}
        \renewcommand*{\backrefentrycount}[2]{#1\ifnum#2>1 ~(#2)\fi}
\fi

\DeclarePairedDelimiter{\ceil}{\lceil}{\rceil}

\onehalfspacing

\usepackage{pgfgantt}

%++++++++++++++++++++++++++++++++++++++++++++++++++++++++++++++++++++++++++++++
% Identificacoes (se o trabalho for em inglês, insira os dados em inglês)
% Para entradas abreviadas de Professora (Profa.) em português escreva: Prof$^\textnormal{a}$.
%++++++++++++++++++++++++++++++++++++++++++++++++++++++++++++++++++++++++++++++
\course{Engenharia Aeronáutica e Mecânica} % Programa de PG ou Curso de Graduação
\area{Propulsão Aeroespacial e Energia} % Área de concentração na PG (Não utilizado no caso de TG)

% Autor do trabalho: Nome Sobrenome
\authorgender{masc} %sexo: masc ou fem
\author{\textbf{EDEN DENIS} Ferreira da Silva Lopes}{Santos}
\itaauthoraddress{Rua João Fonseca dos Santos, 20, Apto. 93}{12.230-088}{São José dos Campos--SP}

% Titulo da Tese/Dissertação
\title{Aplicação de Inteligência Artificial (AI) em Banco de Dados de Emissão de Motores}

% Orientador
\advisorgender{masc} % masc ou fem
\advisor{Prof.~Dr.}{Pedro Teixeira \textbf{LACAVA}}{ITA}

% Coorientador (Caso não haja coorientador, colocar ambas as variáveis \coadvisorgender e \coadvisor comentadas, com um % na frente)
\coadvisorgender{masc}
\coadvisor{}{}  % {Prof.~Dr.}{Carlos Alberto Rocha \textbf{PIMENTEL}}{ITA}

% Pró-reitor da Pós-graduação
\bossgender{fem}
\boss{Prof.~Dra.}{Emília \textbf{VILLANI}}

%Coordenador do curso no caso de TG
\bosscoursegender{masc}
\bosscourse{Prof.~Dr.}{ - } % masc ou fem

% Palavras-Chaves informadas pela Biblioteca -> utilizada na CIP
\kwcip{Engine taxonomy}
\kwcip{Emission classification}
\kwcip{Inteligência Artificial}
\kwcip{Artificial Intelligence}
\kwcip{Python}

% membros da banca examinadora
\examiner{Prof. Dr.}{Pedro Teixeira \textbf{LACAVA}}{President}{ITA}
\examiner{Prof. Dr.}{}{Internal Member}{ITA} % Vinícius \textbf{MALATESTA}}{Internal Member}{ITA}
\examiner{Prof. Dr.}{}{Internal Member}{ITA} % {Elisan dos Santos \textbf{MAGALHÃES}}{Internal Member}{ITA}
\examiner{Prof. Dr.}{}{Internal Member}{ITA} % {Carlos Alberto Rocha \textbf{PIMENTEL}}{External Member}{UFABC}
\examiner{Prof. Dr.}{}{Internal Member}{ITA} % {\textbf{KARL} Peter Burr}{External Member}{UFABC}

% Data da defesa (mês em maiúsculo, se trabalho em inglês, e minúsculo se trabalho em português) 
\date{14}{JULHO}{2025}

% Número CDU - (somente para TG)
\cdu{621.38}

% Folha de Registro do Documento
\newcommand{\advisorTitle}{Prof.~Dr.}
\newcommand{\advisorName}{Pedro Teixeira \textbf{LACAVA}}

\newcommand{\coadvisorTitle}{Prof.~Dr.}
\newcommand{\coadvisorName}{Carlos Alberto Rocha \textbf{PIMENTEL}}

\newcommand{\dateDay}{14}
\newcommand{\dateMonth}{JULHO}
\newcommand{\dateYear}{2025}

% Data de publicação e defesa (usada na folha de registro)
\newcommand{\frdDataPublicacao}{04 de agosto de 2025}
\newcommand{\frdNumeroRegistro}{DCTA/ITA/TD-029/2020}
\newcommand{\frdPalavrasAutor}{Big Data; Engine Emission; Emission Databank; Clustering; Artificial Intelligence}
\newcommand{\frdPalavrasResultado}{Taxonomia; Emissão de Motores; Planejamento de Dados; Inteligência Artificial; Processamento de Dados}
\newcommand{\frdApresentacao}{
	ITA, São José dos Campos. Curso de Mestrado. Programa de Pós-Graduação em Engenharia Aeronáutica e Mecânica. Área de Propulsão Aeroespacial e Energia.\\
	Orientador: \advisorTitle~\advisorName.\\
	Coorientador: \coadvisorTitle~\coadvisorName.\\
	Defesa em \dateDay/\dateMonth/\dateYear. Publicada em \frdDataPublicacao.
}

\newcommand{\PorcentagemDaLargura}{0.80}

\begin{document}
	
	\newpage
	\pagestyle{fancy}
	\thispagestyle{empty}
	\section*{Lista de Definições}
	\addcontentsline{toc}{section}{Lista de Defini\c{c}\~oes}
    \begin{description}
        \item[EEDB] \textit{Engine Emissions Databank}: base oficial da ICAO que reúne os dados de emissões dos motores certificados.
        
        \item[Índice de Emissão (EI)] Métrica de emissões normalizada pelo consumo de combustível, usualmente em $ \mathrm{g/kg} $ para comparar motores em diferentes modos do ciclo LTO.
        
        \item[Bypass Ratio (B/P Ratio)] Razão entre a vazão de ar que contorna o núcleo do motor e a vazão que passa pelo núcleo, parâmetro central para caracterização de turbofans.
        
        \item[Pressure Ratio] Razão de pressão total de entrada e saída do compressor principal, relacionada ao desempenho termodinâmico do motor.
        
        \item[Ciclo LTO] Sequência de modos de operação (taxi/idle, approach, climb-out, take-off) usada para quantificar emissões regulamentares em solo e baixa altitude.
        
        \item[Dendrograma] Representação hierárquica que mostra as distâncias entre grupos formados pelo algoritmo de clusterização, guiando o corte em $ K $ clusters.
        
        \item[Elbow Data Chart (EDC)] Gráfico que relaciona o número de clusters $ K $ ao WSS (soma dos quadrados das distâncias intracluster), apoiando a identificação de uma região de redução marginal menor; o cotovelo é um diagnóstico, não uma escolha automática.
        
        \item[Kneedle] Heurística que determina o $ K $ ótimo ao maximizar a diferença entre a curva do EDC e a diagonal $ y = x $, destacando o primeiro ganho expressivo de $ WS $.
        
        \item[Distância à corda] Método que mede o afastamento perpendicular entre cada ponto do EDC e a reta que conecta os extremos, escolhendo o $ K $ mais distante dessa corda.
        
        \item[Regressão segmentada] Ajuste de duas retas ao EDC para estimar o ponto de quebra que minimiza o erro quadrático total e define o cotovelo.
        
        \item[Curvatura discreta] Cálculo do maior ângulo formado por três pontos consecutivos do EDC, sinalizando a inflexão de maior concavidade.
        
        \item[One-hot encoding] Estratégia que transforma variáveis categóricas em vetores binários (\textit{dummies}), eliminando relações de ordem à custa do aumento de dimensionalidade.
        
        \item[\textit{Label encoding}] Codificação que atribui inteiros sequenciais a cada categoria, preservando o número de atributos e permitindo reprodutibilidade, ainda que introduza um viés ordinal.
    \end{description}
	
\end{document}

	\documentclass[12pt]{article} % Tamanho de letra 12, papel A4, Estilo de artigo

\usepackage[utf8]{inputenc} % Pacote para permitir o uso de acentos
% Ajustes de hifenização
\tolerance=1
\emergencystretch=\maxdimen
\hyphenpenalty=10000
\hbadness=10000
\usepackage{lipsum} % Pacote para gerar texto aleatório
\usepackage{ragged2e} % Pacote para mais opções de alinhamento
\usepackage{graphicx} % Pacote para inserir figuras jpeg, pdf e png
\usepackage{ae}
\usepackage{epsfig}
\usepackage{mathtools}
\usepackage[export]{adjustbox}
\usepackage{lineno}
\usepackage{amsthm}
\usepackage{commath}
\usepackage{color, colortbl}
\usepackage{bbm}
\usepackage{setspace} % Pacote para ajustar espaçamento entre linhas

% Instructions to prevent LaTeX from hyphenating words
\tolerance=1
\emergencystretch=\maxdimen
\hyphenpenalty=10000
\hbadness=10000

\graphicspath{{figures/}} % Indica o diretório no qual as imagens estão sendo salvas, evitando que precisamos digitar parte do caminho da pasta que contém as imagens
\usepackage[rightcaption]{sidecap} % Pacote para mais opções de figura, colocar, título da legenda da figura ao lado direito
\usepackage{url} % Pacote para citar referências de sites
\usepackage{makeidx} % Pacote para criar índice
\makeindex % Habilitar a criação do índice
\usepackage{longtable} % Pacote para usar tabelas longas

% ----------
% Pacotes para dar quebra de linha dentro da tabela:
%\usepackage{fourier} 
%\usepackage{array}
\usepackage{makecell}

\renewcommand\theadalign{bc}
\renewcommand\theadfont{\bfseries}
\renewcommand\theadgape{\Gape[4pt]}
\renewcommand\cellgape{\Gape[4pt]}
% ----------

\usepackage{wrapfig} % Pacote para alinhar figuras à esquerda ou à direita do texto
\usepackage{ae}
\usepackage{epsfig}
\usepackage{mathtools}
\usepackage[export]{adjustbox}
\usepackage{lineno}
\usepackage{amsfonts}
\usepackage{amsthm}
\usepackage{commath}
\usepackage{color, colortbl}
\usepackage{bbm}

% ----------
% Pacote para inserir cabeçalho e rodapé:
% Atenção! Se remover o cabeçalho e rodapé, o LaTeX coloca como cabeçalho o padrão dele, que é bem interessante também!
% \usepackage{fancyhdr} % Pacote para inserir cabeçalho e rodapé
% \pagestyle{fancy} % Estilo do cabeçalho e rodapé, neste caso apenas uma linha no cabeçalho
%\lhead{{\footnotesize Prandiano\textregistered}} % Inserir texto no canto esquerdo do cabeçalho
%\chead{centro} % Inserir texto no centro do cabeçalho
%\rhead{{\footnotesize Curso de Python}} % Inserir texto no canto direito do cabeçalho
%\lfoot{{\footnotesize Copyright \textcopyright 2018 Prandiano}} % Inserir texto no canto esquerdo do rodapé
%\cfoot{centro} % Inserir texto no centro do rodapé
%\rfoot{{\footnotesize Reprodução proibida: Lei 9.610 de 19/02/1998}} % Inserir texto no canto direito do rodapé
% O comando \renewcommand re-nova o parâmetro especificado
%\renewcommand{\footrulewidth}{0.4pt} % Inserir uma linha acima do rodapé
% ----------

% ----------
% Pacote para inserir múltiplas figuras:
\usepackage{subfigure} % Pacote para inserir múltiplas figuras
% \usepackage{subcaption} % Pacote incompatível com subfigure
%\usepackage{caption}
% ----------

\usepackage{multicol} % Pacote que permite usar multicolunas
\usepackage{float} % Pacote para usar o "H" do figure
\usepackage[dvipsnames]{xcolor} % Pacote de cores para textos e possivelmente outros
\usepackage{amsmath} % Pacote para usar o cases
\usepackage[makeroom]{cancel} % Pacote para usar o cancelar (em e embaixo)
\usepackage{enumitem} % Pacote para enumerar itens
\usepackage{amssymb} % Pacote para inserir os símbolos dos conjuntos numéricos N, Z, Q, R, C
\usepackage[T1]{fontenc}
\usepackage{enumitem}
\usepackage{textcomp}
\usepackage{pdfpages} % Pacote para inserir página em PDF
\usepackage{courier} % Pacote para usar a fonte courier
%\usepackage{matlab-prettifier}
\usepackage{lscape} % Pacote para mudar sentido da página
\newtheorem{theorem}{Teorema} % Pacote para inserir Teorema
\newtheorem{corollary}{Corollary} % Pacote para inserir Corolário
\newtheorem{remark}{Remark}
\newtheorem{mytheorem}{Definition}
\usepackage{hyperref} % Pacote para inserir hyperref
\usepackage{multirow} % Pacote para usar os comandos \multirow
\usepackage{pdfpages} % Pacote para incluir PDF's
\usepackage{listingsutf8} % Pacote para inserir arquivos externos tipo .txt
\usepackage{etoolbox}
\AtBeginEnvironment{verbatim}{\footnotesize}

\usepackage{standalone} % Pacote para inserir sub-arquivos LaTeX
\usepackage{bookmark} % Pacote para NÃO tratar o Índice como Apêndice e colocar letra

% Nome da listagem em português
\renewcommand{\lstlistingname}{Listagem}

% REFERÊNCIAS: ----------
%\usepackage[numbers, compress]{natbib} % Pacote para usar as referências bibliográficas
%\usepackage{usebib}
% \bibliographystyle{references} % Estilo das referências
%\setcitestyle{authoryear,aysep={},open={(},close={)}}
%\newcommand{\myand}{e}

% ----------


% PARÂMETROS DA FUNÇÃO LISTINGS: ----------
\lstset
{
    basicstyle=\small\ttfamily,
    commentstyle=\color{Green},
    keywordstyle=\color{Cerulean},
    frame=single,
    language=Python,
    morekeywords={True, False},
    numbers=left,
    numbersep=10pt,
    numberstyle=\footnotesize\color{Gray},
    showstringspaces=false,
    stringstyle=\color{Mulberry},
    tabsize=3,
    breaklines=true, % Este comando é para quebrar uma linha comprida
    postbreak=\mbox{\textcolor{red}{$\hookrightarrow$}\space}, % Este comando é para indicar com uma seta vermelha que a linha continua embaixo 
    extendedchars=true,
    literate=
        {á}{{\'a}}1 {Á}{{\'A}}1 {à}{{\`a}}1 {À}{{\`A}}1 {ã}{{\~a}}1 {Ã}{{\~A}}1 {â}{{\^a}}1 {Â}{{\^A}}1
        {ç}{{\c{c}}}1 {Ç}{{\c{C}}}1 {é}{{\'e}}1 {É}{{\'E}}1 {ê}{{\^e}}1 {Ê}{{\^E}}1 {í}{{\'\i}}1 {Í}{{\'I}}1
        {ó}{{\'o}}1 {Ó}{{\'O}}1 {ô}{{\^o}}1 {Ô}{{\^O}}1 {õ}{{\~o}}1 {Õ}{{\~O}}1 {ú}{{\'u}}1 {Ú}{{\'U}}1
        {ü}{{\"u}}1 {Ü}{{\"U}}1 {“}{{``}}1 {”}{{''}}1 {’}{{'}}1 {—}{{--}}1 {–}{{--}}1 {−}{{-}}1
        {±}{{$\pm$}}1 {σ}{{$\sigma$}}1 {→}{{$\to$}}1 {…}{{\ldots}}1
}

\lstset{inputencoding=utf8/latin1}

\lstdefinelanguage{Matlab}{
	language=Matlab,
	basicstyle=\ttfamily\small,
	keywordstyle=\color{blue}\bfseries,
	stringstyle=\color{red},
	commentstyle=\color{gray}\itshape,
	morekeywords={function, end, if, else, for, while, switch, case, otherwise},
	sensitive=true
}
% ----------

% ----------
\usepackage{nomencl} % Pacote para criar nomenclatura e lista de nomenclatura
\makenomenclature % Criar a nomenclatura
\renewcommand{\nomname}{Lista de Nomenclaturas} % Renomear o título da nomenclatura
% ----------

\usepackage{amsmath} % Pacote para inserir matrix com chaves, entre outras coisas

\usepackage{booktabs} % Pacote para usar \toprule, \midrule e \bottomrule


% EXECUTAR CÁLCULOS A PARTIR DA LINHA DE COMANDO: ---------
%\usepackage{siunitx, xfp} % Pacote para executar cálculos a partir da linha de comando
%\usepackage{l3regex}%unless you have very new expl3
%%\sisetup{locale=DE}
%%\let\oldnum\num
%%\ExplSyntaxOn
%\def\num#1{\def\tmp{#1}\regex_replace_all:nnN{,([0-9])}{.\1}{\tmp}\oldnum{\tmp}}
%\ExplSyntaxOff
% ----------


%\pagecolor{black} % Alterar a cor da página
%\color{white} % Alterar a cor do texto

% ----------
% Pacotes para inserir algoritmos e pseudo-códigos:
\setlength{\marginparwidth}{2.5cm}
\usepackage[colorinlistoftodos]{todonotes}
\usepackage{algorithm} % Pacote para inserir algoritmos
\usepackage{algpseudocode} % Pacote para inserir algoritmos e pseudo-códigos
\floatname{algorithm}{Algoritmo} % Renomear a palavra Algorithm para Procedimento
\renewcommand{\algorithmicrequire}{\textbf{entrada:}} % Renomear a palavra Input para Entrada
\renewcommand{\algorithmicensure}{\textbf{saída:}} % Renomear a palavra Output para Saída
\renewcommand{\algorithmicwhile}{\textbf{enquanto}} % Renomear a palavra While para Enquanto
\renewcommand{\algorithmicif}{\textbf{se}} % Renomear a palavra If para Se
\renewcommand{\algorithmicend}{\textbf{fim}} % Renomear a palavra End Output para Fim do
\renewcommand{\algorithmicprocedure}{\textbf{procedimento}} % Renomear a palavra Procedure para Procedimento
\renewcommand{\algorithmicfor}{\textbf{para}} % Renomear a palavra For para Para
\renewcommand{\algorithmicelse}{\textbf{senão}} % Renomear a palavra Else para Senão
\renewcommand{\algorithmicfunction}{\textbf{função}} % Renomear a palavra Function para Função 
\renewcommand{\algorithmicthen}{\textbf{então}} % Renomear a palavra Then para Então
\renewcommand{\algorithmicdo}{\textbf{fazer}} % Renomear a palavra Do para Fazer
% ----------

\newcommand{\quotes}[1]{``#1''} % Comando para inserir um texto dentro de aspas-dupla com espaço após o final

% FLUXOGRAMA: ----------
\usepackage{tikz}
\usetikzlibrary{positioning, shapes.geometric}
\tikzstyle{line}=[draw] % here

\tikzstyle{startstop} = [rectangle, rounded corners, minimum width=3cm,
minimum height=1cm,text centered, draw=black, fill=white!30]
\tikzstyle{arrow} = [thick, ->, >=stealth]
\tikzstyle{process} = [rectangle, minimum width=3cm, minimum height=1cm, text centered, draw=black, fill=white!30]
\tikzstyle{io} = [trapezium, trapezium left angle=70, trapezium right angle=110, minimum width=3cm, minimum height=1cm, text centered, draw=black, fill=white!30]
\tikzstyle{connector} = [draw, circle, fill=white!20]
\tikzstyle{decision} = [diamond, minimum width=3cm, minimum height=1cm, text centered, draw=black, fill=white!30, aspect=3]
% ----------


% CAIXA ENTORNO DAS FÓRMULAS: ----------
\usepackage{dashbox}
\newcommand\dboxed[1]{\dbox{\ensuremath{#1}}}

% ----------

%\usepackage{tabularx} % para colunas com largura ajustável

% -------- CONFIGURAÇÃO DE MARGENS --------
% \usepackage[top=2cm,bottom=2cm,left=3cm,right=2cm]{geometry}

% -------- PACOTE fancyhdr --------
\usepackage{fancyhdr}

\setlength{\marginparwidth}{2cm}

% Ajuste do tamanho do cabeçalho para acompanhar exatamente \textwidth
\setlength{\headwidth}{\textwidth}

% Define um headheight mínimo (para evitar avisos)
\setlength{\headheight}{15pt} 

% Limpa configurações padrão de cabeçalho/rodapé
\pagestyle{fancy}
\fancyhf{} 

% Marcações automáticas de seção para o cabeçalho (opcional)
\renewcommand{\sectionmark}[1]{%
	\markboth{\thesection\ #1}{}%
}

% Texto do cabeçalho à esquerda: seção e título
% Usamos \leftmark pois o \sectionmark acima grava no \leftmark
\fancyhead[L]{\small \leftmark}

% Texto do cabeçalho à direita: número de página
\fancyhead[R]{\small \thepage}

% Linha horizontal no cabeçalho (espessura 0.4pt)
\renewcommand{\headrulewidth}{0.4pt}

% Sobrescreve o título padrão “References” para “Referências”
%\renewcommand{\refname}{Referências}
% Tradução de autoref para o português
\def\figureautorefname{Figura}
\def\tableautorefname{Tabela}
\def\equationautorefname{Equa\c{c}\~ao}

% \hypersetup{backref=true}

\ifdefined\backref
        \renewcommand*{\backref}[1]{}
        \renewcommand*{\backrefentrycount}[2]{#1\ifnum#2>1 ~(#2)\fi}
\fi

\DeclarePairedDelimiter{\ceil}{\lceil}{\rceil}

\onehalfspacing

\usepackage{pgfgantt}

%++++++++++++++++++++++++++++++++++++++++++++++++++++++++++++++++++++++++++++++
% Identificacoes (se o trabalho for em inglês, insira os dados em inglês)
% Para entradas abreviadas de Professora (Profa.) em português escreva: Prof$^\textnormal{a}$.
%++++++++++++++++++++++++++++++++++++++++++++++++++++++++++++++++++++++++++++++
\course{Engenharia Aeronáutica e Mecânica} % Programa de PG ou Curso de Graduação
\area{Propulsão Aeroespacial e Energia} % Área de concentração na PG (Não utilizado no caso de TG)

% Autor do trabalho: Nome Sobrenome
\authorgender{masc} %sexo: masc ou fem
\author{\textbf{EDEN DENIS} Ferreira da Silva Lopes}{Santos}
\itaauthoraddress{Rua João Fonseca dos Santos, 20, Apto. 93}{12.230-088}{São José dos Campos--SP}

% Titulo da Tese/Dissertação
\title{Aplicação de Inteligência Artificial (AI) em Banco de Dados de Emissão de Motores}

% Orientador
\advisorgender{masc} % masc ou fem
\advisor{Prof.~Dr.}{Pedro Teixeira \textbf{LACAVA}}{ITA}

% Coorientador (Caso não haja coorientador, colocar ambas as variáveis \coadvisorgender e \coadvisor comentadas, com um % na frente)
\coadvisorgender{masc}
\coadvisor{}{}  % {Prof.~Dr.}{Carlos Alberto Rocha \textbf{PIMENTEL}}{ITA}

% Pró-reitor da Pós-graduação
\bossgender{fem}
\boss{Prof.~Dra.}{Emília \textbf{VILLANI}}

%Coordenador do curso no caso de TG
\bosscoursegender{masc}
\bosscourse{Prof.~Dr.}{ - } % masc ou fem

% Palavras-Chaves informadas pela Biblioteca -> utilizada na CIP
\kwcip{Engine taxonomy}
\kwcip{Emission classification}
\kwcip{Inteligência Artificial}
\kwcip{Artificial Intelligence}
\kwcip{Python}

% membros da banca examinadora
\examiner{Prof. Dr.}{Pedro Teixeira \textbf{LACAVA}}{President}{ITA}
\examiner{Prof. Dr.}{}{Internal Member}{ITA} % Vinícius \textbf{MALATESTA}}{Internal Member}{ITA}
\examiner{Prof. Dr.}{}{Internal Member}{ITA} % {Elisan dos Santos \textbf{MAGALHÃES}}{Internal Member}{ITA}
\examiner{Prof. Dr.}{}{Internal Member}{ITA} % {Carlos Alberto Rocha \textbf{PIMENTEL}}{External Member}{UFABC}
\examiner{Prof. Dr.}{}{Internal Member}{ITA} % {\textbf{KARL} Peter Burr}{External Member}{UFABC}

% Data da defesa (mês em maiúsculo, se trabalho em inglês, e minúsculo se trabalho em português) 
\date{14}{JULHO}{2025}

% Número CDU - (somente para TG)
\cdu{621.38}

% Folha de Registro do Documento
\newcommand{\advisorTitle}{Prof.~Dr.}
\newcommand{\advisorName}{Pedro Teixeira \textbf{LACAVA}}

\newcommand{\coadvisorTitle}{Prof.~Dr.}
\newcommand{\coadvisorName}{Carlos Alberto Rocha \textbf{PIMENTEL}}

\newcommand{\dateDay}{14}
\newcommand{\dateMonth}{JULHO}
\newcommand{\dateYear}{2025}

% Data de publicação e defesa (usada na folha de registro)
\newcommand{\frdDataPublicacao}{04 de agosto de 2025}
\newcommand{\frdNumeroRegistro}{DCTA/ITA/TD-029/2020}
\newcommand{\frdPalavrasAutor}{Big Data; Engine Emission; Emission Databank; Clustering; Artificial Intelligence}
\newcommand{\frdPalavrasResultado}{Taxonomia; Emissão de Motores; Planejamento de Dados; Inteligência Artificial; Processamento de Dados}
\newcommand{\frdApresentacao}{
	ITA, São José dos Campos. Curso de Mestrado. Programa de Pós-Graduação em Engenharia Aeronáutica e Mecânica. Área de Propulsão Aeroespacial e Energia.\\
	Orientador: \advisorTitle~\advisorName.\\
	Coorientador: \coadvisorTitle~\coadvisorName.\\
	Defesa em \dateDay/\dateMonth/\dateYear. Publicada em \frdDataPublicacao.
}

\newcommand{\PorcentagemDaLargura}{0.80}

\begin{document}
	
	\newpage
	\pagestyle{fancy}
	\thispagestyle{empty}
	\renewcommand{\listfigurename}{Lista de Ilustrações}
	\listoffigures % Criar lista de figuras existentes no documento
	\addcontentsline{toc}{section}{Lista de Ilustrações}
	
\end{document}
	\input{list_of_nomenclatures.tex}
	\documentclass[12pt]{article} % Tamanho de letra 12, papel A4, Estilo de artigo

\usepackage[utf8]{inputenc} % Pacote para permitir o uso de acentos
% Ajustes de hifenização
\tolerance=1
\emergencystretch=\maxdimen
\hyphenpenalty=10000
\hbadness=10000
\usepackage{lipsum} % Pacote para gerar texto aleatório
\usepackage{ragged2e} % Pacote para mais opções de alinhamento
\usepackage{graphicx} % Pacote para inserir figuras jpeg, pdf e png
\usepackage{ae}
\usepackage{epsfig}
\usepackage{mathtools}
\usepackage[export]{adjustbox}
\usepackage{lineno}
\usepackage{amsthm}
\usepackage{commath}
\usepackage{color, colortbl}
\usepackage{bbm}
\usepackage{setspace} % Pacote para ajustar espaçamento entre linhas

% Instructions to prevent LaTeX from hyphenating words
\tolerance=1
\emergencystretch=\maxdimen
\hyphenpenalty=10000
\hbadness=10000

\graphicspath{{figures/}} % Indica o diretório no qual as imagens estão sendo salvas, evitando que precisamos digitar parte do caminho da pasta que contém as imagens
\usepackage[rightcaption]{sidecap} % Pacote para mais opções de figura, colocar, título da legenda da figura ao lado direito
\usepackage{url} % Pacote para citar referências de sites
\usepackage{makeidx} % Pacote para criar índice
\makeindex % Habilitar a criação do índice
\usepackage{longtable} % Pacote para usar tabelas longas

% ----------
% Pacotes para dar quebra de linha dentro da tabela:
%\usepackage{fourier} 
%\usepackage{array}
\usepackage{makecell}

\renewcommand\theadalign{bc}
\renewcommand\theadfont{\bfseries}
\renewcommand\theadgape{\Gape[4pt]}
\renewcommand\cellgape{\Gape[4pt]}
% ----------

\usepackage{wrapfig} % Pacote para alinhar figuras à esquerda ou à direita do texto
\usepackage{ae}
\usepackage{epsfig}
\usepackage{mathtools}
\usepackage[export]{adjustbox}
\usepackage{lineno}
\usepackage{amsfonts}
\usepackage{amsthm}
\usepackage{commath}
\usepackage{color, colortbl}
\usepackage{bbm}

% ----------
% Pacote para inserir cabeçalho e rodapé:
% Atenção! Se remover o cabeçalho e rodapé, o LaTeX coloca como cabeçalho o padrão dele, que é bem interessante também!
% \usepackage{fancyhdr} % Pacote para inserir cabeçalho e rodapé
% \pagestyle{fancy} % Estilo do cabeçalho e rodapé, neste caso apenas uma linha no cabeçalho
%\lhead{{\footnotesize Prandiano\textregistered}} % Inserir texto no canto esquerdo do cabeçalho
%\chead{centro} % Inserir texto no centro do cabeçalho
%\rhead{{\footnotesize Curso de Python}} % Inserir texto no canto direito do cabeçalho
%\lfoot{{\footnotesize Copyright \textcopyright 2018 Prandiano}} % Inserir texto no canto esquerdo do rodapé
%\cfoot{centro} % Inserir texto no centro do rodapé
%\rfoot{{\footnotesize Reprodução proibida: Lei 9.610 de 19/02/1998}} % Inserir texto no canto direito do rodapé
% O comando \renewcommand re-nova o parâmetro especificado
%\renewcommand{\footrulewidth}{0.4pt} % Inserir uma linha acima do rodapé
% ----------

% ----------
% Pacote para inserir múltiplas figuras:
\usepackage{subfigure} % Pacote para inserir múltiplas figuras
% \usepackage{subcaption} % Pacote incompatível com subfigure
%\usepackage{caption}
% ----------

\usepackage{multicol} % Pacote que permite usar multicolunas
\usepackage{float} % Pacote para usar o "H" do figure
\usepackage[dvipsnames]{xcolor} % Pacote de cores para textos e possivelmente outros
\usepackage{amsmath} % Pacote para usar o cases
\usepackage[makeroom]{cancel} % Pacote para usar o cancelar (em e embaixo)
\usepackage{enumitem} % Pacote para enumerar itens
\usepackage{amssymb} % Pacote para inserir os símbolos dos conjuntos numéricos N, Z, Q, R, C
\usepackage[T1]{fontenc}
\usepackage{enumitem}
\usepackage{textcomp}
\usepackage{pdfpages} % Pacote para inserir página em PDF
\usepackage{courier} % Pacote para usar a fonte courier
%\usepackage{matlab-prettifier}
\usepackage{lscape} % Pacote para mudar sentido da página
\newtheorem{theorem}{Teorema} % Pacote para inserir Teorema
\newtheorem{corollary}{Corollary} % Pacote para inserir Corolário
\newtheorem{remark}{Remark}
\newtheorem{mytheorem}{Definition}
\usepackage{hyperref} % Pacote para inserir hyperref
\usepackage{multirow} % Pacote para usar os comandos \multirow
\usepackage{pdfpages} % Pacote para incluir PDF's
\usepackage{listingsutf8} % Pacote para inserir arquivos externos tipo .txt
\usepackage{etoolbox}
\AtBeginEnvironment{verbatim}{\footnotesize}

\usepackage{standalone} % Pacote para inserir sub-arquivos LaTeX
\usepackage{bookmark} % Pacote para NÃO tratar o Índice como Apêndice e colocar letra

% Nome da listagem em português
\renewcommand{\lstlistingname}{Listagem}

% REFERÊNCIAS: ----------
%\usepackage[numbers, compress]{natbib} % Pacote para usar as referências bibliográficas
%\usepackage{usebib}
% \bibliographystyle{references} % Estilo das referências
%\setcitestyle{authoryear,aysep={},open={(},close={)}}
%\newcommand{\myand}{e}

% ----------


% PARÂMETROS DA FUNÇÃO LISTINGS: ----------
\lstset
{
    basicstyle=\small\ttfamily,
    commentstyle=\color{Green},
    keywordstyle=\color{Cerulean},
    frame=single,
    language=Python,
    morekeywords={True, False},
    numbers=left,
    numbersep=10pt,
    numberstyle=\footnotesize\color{Gray},
    showstringspaces=false,
    stringstyle=\color{Mulberry},
    tabsize=3,
    breaklines=true, % Este comando é para quebrar uma linha comprida
    postbreak=\mbox{\textcolor{red}{$\hookrightarrow$}\space}, % Este comando é para indicar com uma seta vermelha que a linha continua embaixo 
    extendedchars=true,
    literate=
        {á}{{\'a}}1 {Á}{{\'A}}1 {à}{{\`a}}1 {À}{{\`A}}1 {ã}{{\~a}}1 {Ã}{{\~A}}1 {â}{{\^a}}1 {Â}{{\^A}}1
        {ç}{{\c{c}}}1 {Ç}{{\c{C}}}1 {é}{{\'e}}1 {É}{{\'E}}1 {ê}{{\^e}}1 {Ê}{{\^E}}1 {í}{{\'\i}}1 {Í}{{\'I}}1
        {ó}{{\'o}}1 {Ó}{{\'O}}1 {ô}{{\^o}}1 {Ô}{{\^O}}1 {õ}{{\~o}}1 {Õ}{{\~O}}1 {ú}{{\'u}}1 {Ú}{{\'U}}1
        {ü}{{\"u}}1 {Ü}{{\"U}}1 {“}{{``}}1 {”}{{''}}1 {’}{{'}}1 {—}{{--}}1 {–}{{--}}1 {−}{{-}}1
        {±}{{$\pm$}}1 {σ}{{$\sigma$}}1 {→}{{$\to$}}1 {…}{{\ldots}}1
}

\lstset{inputencoding=utf8/latin1}

\lstdefinelanguage{Matlab}{
	language=Matlab,
	basicstyle=\ttfamily\small,
	keywordstyle=\color{blue}\bfseries,
	stringstyle=\color{red},
	commentstyle=\color{gray}\itshape,
	morekeywords={function, end, if, else, for, while, switch, case, otherwise},
	sensitive=true
}
% ----------

% ----------
\usepackage{nomencl} % Pacote para criar nomenclatura e lista de nomenclatura
\makenomenclature % Criar a nomenclatura
\renewcommand{\nomname}{Lista de Nomenclaturas} % Renomear o título da nomenclatura
% ----------

\usepackage{amsmath} % Pacote para inserir matrix com chaves, entre outras coisas

\usepackage{booktabs} % Pacote para usar \toprule, \midrule e \bottomrule


% EXECUTAR CÁLCULOS A PARTIR DA LINHA DE COMANDO: ---------
%\usepackage{siunitx, xfp} % Pacote para executar cálculos a partir da linha de comando
%\usepackage{l3regex}%unless you have very new expl3
%%\sisetup{locale=DE}
%%\let\oldnum\num
%%\ExplSyntaxOn
%\def\num#1{\def\tmp{#1}\regex_replace_all:nnN{,([0-9])}{.\1}{\tmp}\oldnum{\tmp}}
%\ExplSyntaxOff
% ----------


%\pagecolor{black} % Alterar a cor da página
%\color{white} % Alterar a cor do texto

% ----------
% Pacotes para inserir algoritmos e pseudo-códigos:
\setlength{\marginparwidth}{2.5cm}
\usepackage[colorinlistoftodos]{todonotes}
\usepackage{algorithm} % Pacote para inserir algoritmos
\usepackage{algpseudocode} % Pacote para inserir algoritmos e pseudo-códigos
\floatname{algorithm}{Algoritmo} % Renomear a palavra Algorithm para Procedimento
\renewcommand{\algorithmicrequire}{\textbf{entrada:}} % Renomear a palavra Input para Entrada
\renewcommand{\algorithmicensure}{\textbf{saída:}} % Renomear a palavra Output para Saída
\renewcommand{\algorithmicwhile}{\textbf{enquanto}} % Renomear a palavra While para Enquanto
\renewcommand{\algorithmicif}{\textbf{se}} % Renomear a palavra If para Se
\renewcommand{\algorithmicend}{\textbf{fim}} % Renomear a palavra End Output para Fim do
\renewcommand{\algorithmicprocedure}{\textbf{procedimento}} % Renomear a palavra Procedure para Procedimento
\renewcommand{\algorithmicfor}{\textbf{para}} % Renomear a palavra For para Para
\renewcommand{\algorithmicelse}{\textbf{senão}} % Renomear a palavra Else para Senão
\renewcommand{\algorithmicfunction}{\textbf{função}} % Renomear a palavra Function para Função 
\renewcommand{\algorithmicthen}{\textbf{então}} % Renomear a palavra Then para Então
\renewcommand{\algorithmicdo}{\textbf{fazer}} % Renomear a palavra Do para Fazer
% ----------

\newcommand{\quotes}[1]{``#1''} % Comando para inserir um texto dentro de aspas-dupla com espaço após o final

% FLUXOGRAMA: ----------
\usepackage{tikz}
\usetikzlibrary{positioning, shapes.geometric}
\tikzstyle{line}=[draw] % here

\tikzstyle{startstop} = [rectangle, rounded corners, minimum width=3cm,
minimum height=1cm,text centered, draw=black, fill=white!30]
\tikzstyle{arrow} = [thick, ->, >=stealth]
\tikzstyle{process} = [rectangle, minimum width=3cm, minimum height=1cm, text centered, draw=black, fill=white!30]
\tikzstyle{io} = [trapezium, trapezium left angle=70, trapezium right angle=110, minimum width=3cm, minimum height=1cm, text centered, draw=black, fill=white!30]
\tikzstyle{connector} = [draw, circle, fill=white!20]
\tikzstyle{decision} = [diamond, minimum width=3cm, minimum height=1cm, text centered, draw=black, fill=white!30, aspect=3]
% ----------


% CAIXA ENTORNO DAS FÓRMULAS: ----------
\usepackage{dashbox}
\newcommand\dboxed[1]{\dbox{\ensuremath{#1}}}

% ----------

%\usepackage{tabularx} % para colunas com largura ajustável

% -------- CONFIGURAÇÃO DE MARGENS --------
% \usepackage[top=2cm,bottom=2cm,left=3cm,right=2cm]{geometry}

% -------- PACOTE fancyhdr --------
\usepackage{fancyhdr}

\setlength{\marginparwidth}{2cm}

% Ajuste do tamanho do cabeçalho para acompanhar exatamente \textwidth
\setlength{\headwidth}{\textwidth}

% Define um headheight mínimo (para evitar avisos)
\setlength{\headheight}{15pt} 

% Limpa configurações padrão de cabeçalho/rodapé
\pagestyle{fancy}
\fancyhf{} 

% Marcações automáticas de seção para o cabeçalho (opcional)
\renewcommand{\sectionmark}[1]{%
	\markboth{\thesection\ #1}{}%
}

% Texto do cabeçalho à esquerda: seção e título
% Usamos \leftmark pois o \sectionmark acima grava no \leftmark
\fancyhead[L]{\small \leftmark}

% Texto do cabeçalho à direita: número de página
\fancyhead[R]{\small \thepage}

% Linha horizontal no cabeçalho (espessura 0.4pt)
\renewcommand{\headrulewidth}{0.4pt}

% Sobrescreve o título padrão “References” para “Referências”
%\renewcommand{\refname}{Referências}
% Tradução de autoref para o português
\def\figureautorefname{Figura}
\def\tableautorefname{Tabela}
\def\equationautorefname{Equa\c{c}\~ao}

% \hypersetup{backref=true}

\ifdefined\backref
        \renewcommand*{\backref}[1]{}
        \renewcommand*{\backrefentrycount}[2]{#1\ifnum#2>1 ~(#2)\fi}
\fi

\DeclarePairedDelimiter{\ceil}{\lceil}{\rceil}

\onehalfspacing

\usepackage{pgfgantt}

%++++++++++++++++++++++++++++++++++++++++++++++++++++++++++++++++++++++++++++++
% Identificacoes (se o trabalho for em inglês, insira os dados em inglês)
% Para entradas abreviadas de Professora (Profa.) em português escreva: Prof$^\textnormal{a}$.
%++++++++++++++++++++++++++++++++++++++++++++++++++++++++++++++++++++++++++++++
\course{Engenharia Aeronáutica e Mecânica} % Programa de PG ou Curso de Graduação
\area{Propulsão Aeroespacial e Energia} % Área de concentração na PG (Não utilizado no caso de TG)

% Autor do trabalho: Nome Sobrenome
\authorgender{masc} %sexo: masc ou fem
\author{\textbf{EDEN DENIS} Ferreira da Silva Lopes}{Santos}
\itaauthoraddress{Rua João Fonseca dos Santos, 20, Apto. 93}{12.230-088}{São José dos Campos--SP}

% Titulo da Tese/Dissertação
\title{Aplicação de Inteligência Artificial (AI) em Banco de Dados de Emissão de Motores}

% Orientador
\advisorgender{masc} % masc ou fem
\advisor{Prof.~Dr.}{Pedro Teixeira \textbf{LACAVA}}{ITA}

% Coorientador (Caso não haja coorientador, colocar ambas as variáveis \coadvisorgender e \coadvisor comentadas, com um % na frente)
\coadvisorgender{masc}
\coadvisor{}{}  % {Prof.~Dr.}{Carlos Alberto Rocha \textbf{PIMENTEL}}{ITA}

% Pró-reitor da Pós-graduação
\bossgender{fem}
\boss{Prof.~Dra.}{Emília \textbf{VILLANI}}

%Coordenador do curso no caso de TG
\bosscoursegender{masc}
\bosscourse{Prof.~Dr.}{ - } % masc ou fem

% Palavras-Chaves informadas pela Biblioteca -> utilizada na CIP
\kwcip{Engine taxonomy}
\kwcip{Emission classification}
\kwcip{Inteligência Artificial}
\kwcip{Artificial Intelligence}
\kwcip{Python}

% membros da banca examinadora
\examiner{Prof. Dr.}{Pedro Teixeira \textbf{LACAVA}}{President}{ITA}
\examiner{Prof. Dr.}{}{Internal Member}{ITA} % Vinícius \textbf{MALATESTA}}{Internal Member}{ITA}
\examiner{Prof. Dr.}{}{Internal Member}{ITA} % {Elisan dos Santos \textbf{MAGALHÃES}}{Internal Member}{ITA}
\examiner{Prof. Dr.}{}{Internal Member}{ITA} % {Carlos Alberto Rocha \textbf{PIMENTEL}}{External Member}{UFABC}
\examiner{Prof. Dr.}{}{Internal Member}{ITA} % {\textbf{KARL} Peter Burr}{External Member}{UFABC}

% Data da defesa (mês em maiúsculo, se trabalho em inglês, e minúsculo se trabalho em português) 
\date{14}{JULHO}{2025}

% Número CDU - (somente para TG)
\cdu{621.38}

% Folha de Registro do Documento
\newcommand{\advisorTitle}{Prof.~Dr.}
\newcommand{\advisorName}{Pedro Teixeira \textbf{LACAVA}}

\newcommand{\coadvisorTitle}{Prof.~Dr.}
\newcommand{\coadvisorName}{Carlos Alberto Rocha \textbf{PIMENTEL}}

\newcommand{\dateDay}{14}
\newcommand{\dateMonth}{JULHO}
\newcommand{\dateYear}{2025}

% Data de publicação e defesa (usada na folha de registro)
\newcommand{\frdDataPublicacao}{04 de agosto de 2025}
\newcommand{\frdNumeroRegistro}{DCTA/ITA/TD-029/2020}
\newcommand{\frdPalavrasAutor}{Big Data; Engine Emission; Emission Databank; Clustering; Artificial Intelligence}
\newcommand{\frdPalavrasResultado}{Taxonomia; Emissão de Motores; Planejamento de Dados; Inteligência Artificial; Processamento de Dados}
\newcommand{\frdApresentacao}{
	ITA, São José dos Campos. Curso de Mestrado. Programa de Pós-Graduação em Engenharia Aeronáutica e Mecânica. Área de Propulsão Aeroespacial e Energia.\\
	Orientador: \advisorTitle~\advisorName.\\
	Coorientador: \coadvisorTitle~\coadvisorName.\\
	Defesa em \dateDay/\dateMonth/\dateYear. Publicada em \frdDataPublicacao.
}

\newcommand{\PorcentagemDaLargura}{0.80}

\begin{document}
	
	\newpage
	\pagestyle{fancy}
	\thispagestyle{empty}
	%\shipout\null % Inserir página em branco
	\renewcommand{\listtablename}{Lista de Tabelas}
	\listoftables % Criar lista de tabelas existentes no documento
	\addcontentsline{toc}{section}{Lista de Tabelas}
	
\end{document}
	
	% Sumario
	\tableofcontents
	
	\mainmatter
	% Os capitulos comecam aqui
	
	% Desenvolvimento:
	\documentclass[12pt]{article}

\usepackage[utf8]{inputenc} % Pacote para permitir o uso de acentos
% Ajustes de hifenização
\tolerance=1
\emergencystretch=\maxdimen
\hyphenpenalty=10000
\hbadness=10000
\usepackage{lipsum} % Pacote para gerar texto aleatório
\usepackage{ragged2e} % Pacote para mais opções de alinhamento
\usepackage{graphicx} % Pacote para inserir figuras jpeg, pdf e png
\usepackage{ae}
\usepackage{epsfig}
\usepackage{mathtools}
\usepackage[export]{adjustbox}
\usepackage{lineno}
\usepackage{amsthm}
\usepackage{commath}
\usepackage{color, colortbl}
\usepackage{bbm}
\usepackage{setspace} % Pacote para ajustar espaçamento entre linhas

% Instructions to prevent LaTeX from hyphenating words
\tolerance=1
\emergencystretch=\maxdimen
\hyphenpenalty=10000
\hbadness=10000

\graphicspath{{figures/}} % Indica o diretório no qual as imagens estão sendo salvas, evitando que precisamos digitar parte do caminho da pasta que contém as imagens
\usepackage[rightcaption]{sidecap} % Pacote para mais opções de figura, colocar, título da legenda da figura ao lado direito
\usepackage{url} % Pacote para citar referências de sites
\usepackage{makeidx} % Pacote para criar índice
\makeindex % Habilitar a criação do índice
\usepackage{longtable} % Pacote para usar tabelas longas

% ----------
% Pacotes para dar quebra de linha dentro da tabela:
%\usepackage{fourier} 
%\usepackage{array}
\usepackage{makecell}

\renewcommand\theadalign{bc}
\renewcommand\theadfont{\bfseries}
\renewcommand\theadgape{\Gape[4pt]}
\renewcommand\cellgape{\Gape[4pt]}
% ----------

\usepackage{wrapfig} % Pacote para alinhar figuras à esquerda ou à direita do texto
\usepackage{ae}
\usepackage{epsfig}
\usepackage{mathtools}
\usepackage[export]{adjustbox}
\usepackage{lineno}
\usepackage{amsfonts}
\usepackage{amsthm}
\usepackage{commath}
\usepackage{color, colortbl}
\usepackage{bbm}

% ----------
% Pacote para inserir cabeçalho e rodapé:
% Atenção! Se remover o cabeçalho e rodapé, o LaTeX coloca como cabeçalho o padrão dele, que é bem interessante também!
% \usepackage{fancyhdr} % Pacote para inserir cabeçalho e rodapé
% \pagestyle{fancy} % Estilo do cabeçalho e rodapé, neste caso apenas uma linha no cabeçalho
%\lhead{{\footnotesize Prandiano\textregistered}} % Inserir texto no canto esquerdo do cabeçalho
%\chead{centro} % Inserir texto no centro do cabeçalho
%\rhead{{\footnotesize Curso de Python}} % Inserir texto no canto direito do cabeçalho
%\lfoot{{\footnotesize Copyright \textcopyright 2018 Prandiano}} % Inserir texto no canto esquerdo do rodapé
%\cfoot{centro} % Inserir texto no centro do rodapé
%\rfoot{{\footnotesize Reprodução proibida: Lei 9.610 de 19/02/1998}} % Inserir texto no canto direito do rodapé
% O comando \renewcommand re-nova o parâmetro especificado
%\renewcommand{\footrulewidth}{0.4pt} % Inserir uma linha acima do rodapé
% ----------

% ----------
% Pacote para inserir múltiplas figuras:
\usepackage{subfigure} % Pacote para inserir múltiplas figuras
% \usepackage{subcaption} % Pacote incompatível com subfigure
%\usepackage{caption}
% ----------

\usepackage{multicol} % Pacote que permite usar multicolunas
\usepackage{float} % Pacote para usar o "H" do figure
\usepackage[dvipsnames]{xcolor} % Pacote de cores para textos e possivelmente outros
\usepackage{amsmath} % Pacote para usar o cases
\usepackage[makeroom]{cancel} % Pacote para usar o cancelar (em e embaixo)
\usepackage{enumitem} % Pacote para enumerar itens
\usepackage{amssymb} % Pacote para inserir os símbolos dos conjuntos numéricos N, Z, Q, R, C
\usepackage[T1]{fontenc}
\usepackage{enumitem}
\usepackage{textcomp}
\usepackage{pdfpages} % Pacote para inserir página em PDF
\usepackage{courier} % Pacote para usar a fonte courier
%\usepackage{matlab-prettifier}
\usepackage{lscape} % Pacote para mudar sentido da página
\newtheorem{theorem}{Teorema} % Pacote para inserir Teorema
\newtheorem{corollary}{Corollary} % Pacote para inserir Corolário
\newtheorem{remark}{Remark}
\newtheorem{mytheorem}{Definition}
\usepackage{hyperref} % Pacote para inserir hyperref
\usepackage{multirow} % Pacote para usar os comandos \multirow
\usepackage{pdfpages} % Pacote para incluir PDF's
\usepackage{listingsutf8} % Pacote para inserir arquivos externos tipo .txt
\usepackage{etoolbox}
\AtBeginEnvironment{verbatim}{\footnotesize}

\usepackage{standalone} % Pacote para inserir sub-arquivos LaTeX
\usepackage{bookmark} % Pacote para NÃO tratar o Índice como Apêndice e colocar letra

% Nome da listagem em português
\renewcommand{\lstlistingname}{Listagem}

% REFERÊNCIAS: ----------
%\usepackage[numbers, compress]{natbib} % Pacote para usar as referências bibliográficas
%\usepackage{usebib}
% \bibliographystyle{references} % Estilo das referências
%\setcitestyle{authoryear,aysep={},open={(},close={)}}
%\newcommand{\myand}{e}

% ----------


% PARÂMETROS DA FUNÇÃO LISTINGS: ----------
\lstset
{
    basicstyle=\small\ttfamily,
    commentstyle=\color{Green},
    keywordstyle=\color{Cerulean},
    frame=single,
    language=Python,
    morekeywords={True, False},
    numbers=left,
    numbersep=10pt,
    numberstyle=\footnotesize\color{Gray},
    showstringspaces=false,
    stringstyle=\color{Mulberry},
    tabsize=3,
    breaklines=true, % Este comando é para quebrar uma linha comprida
    postbreak=\mbox{\textcolor{red}{$\hookrightarrow$}\space}, % Este comando é para indicar com uma seta vermelha que a linha continua embaixo 
    extendedchars=true,
    literate=
        {á}{{\'a}}1 {Á}{{\'A}}1 {à}{{\`a}}1 {À}{{\`A}}1 {ã}{{\~a}}1 {Ã}{{\~A}}1 {â}{{\^a}}1 {Â}{{\^A}}1
        {ç}{{\c{c}}}1 {Ç}{{\c{C}}}1 {é}{{\'e}}1 {É}{{\'E}}1 {ê}{{\^e}}1 {Ê}{{\^E}}1 {í}{{\'\i}}1 {Í}{{\'I}}1
        {ó}{{\'o}}1 {Ó}{{\'O}}1 {ô}{{\^o}}1 {Ô}{{\^O}}1 {õ}{{\~o}}1 {Õ}{{\~O}}1 {ú}{{\'u}}1 {Ú}{{\'U}}1
        {ü}{{\"u}}1 {Ü}{{\"U}}1 {“}{{``}}1 {”}{{''}}1 {’}{{'}}1 {—}{{--}}1 {–}{{--}}1 {−}{{-}}1
        {±}{{$\pm$}}1 {σ}{{$\sigma$}}1 {→}{{$\to$}}1 {…}{{\ldots}}1
}

\lstset{inputencoding=utf8/latin1}

\lstdefinelanguage{Matlab}{
	language=Matlab,
	basicstyle=\ttfamily\small,
	keywordstyle=\color{blue}\bfseries,
	stringstyle=\color{red},
	commentstyle=\color{gray}\itshape,
	morekeywords={function, end, if, else, for, while, switch, case, otherwise},
	sensitive=true
}
% ----------

% ----------
\usepackage{nomencl} % Pacote para criar nomenclatura e lista de nomenclatura
\makenomenclature % Criar a nomenclatura
\renewcommand{\nomname}{Lista de Nomenclaturas} % Renomear o título da nomenclatura
% ----------

\usepackage{amsmath} % Pacote para inserir matrix com chaves, entre outras coisas

\usepackage{booktabs} % Pacote para usar \toprule, \midrule e \bottomrule


% EXECUTAR CÁLCULOS A PARTIR DA LINHA DE COMANDO: ---------
%\usepackage{siunitx, xfp} % Pacote para executar cálculos a partir da linha de comando
%\usepackage{l3regex}%unless you have very new expl3
%%\sisetup{locale=DE}
%%\let\oldnum\num
%%\ExplSyntaxOn
%\def\num#1{\def\tmp{#1}\regex_replace_all:nnN{,([0-9])}{.\1}{\tmp}\oldnum{\tmp}}
%\ExplSyntaxOff
% ----------


%\pagecolor{black} % Alterar a cor da página
%\color{white} % Alterar a cor do texto

% ----------
% Pacotes para inserir algoritmos e pseudo-códigos:
\setlength{\marginparwidth}{2.5cm}
\usepackage[colorinlistoftodos]{todonotes}
\usepackage{algorithm} % Pacote para inserir algoritmos
\usepackage{algpseudocode} % Pacote para inserir algoritmos e pseudo-códigos
\floatname{algorithm}{Algoritmo} % Renomear a palavra Algorithm para Procedimento
\renewcommand{\algorithmicrequire}{\textbf{entrada:}} % Renomear a palavra Input para Entrada
\renewcommand{\algorithmicensure}{\textbf{saída:}} % Renomear a palavra Output para Saída
\renewcommand{\algorithmicwhile}{\textbf{enquanto}} % Renomear a palavra While para Enquanto
\renewcommand{\algorithmicif}{\textbf{se}} % Renomear a palavra If para Se
\renewcommand{\algorithmicend}{\textbf{fim}} % Renomear a palavra End Output para Fim do
\renewcommand{\algorithmicprocedure}{\textbf{procedimento}} % Renomear a palavra Procedure para Procedimento
\renewcommand{\algorithmicfor}{\textbf{para}} % Renomear a palavra For para Para
\renewcommand{\algorithmicelse}{\textbf{senão}} % Renomear a palavra Else para Senão
\renewcommand{\algorithmicfunction}{\textbf{função}} % Renomear a palavra Function para Função 
\renewcommand{\algorithmicthen}{\textbf{então}} % Renomear a palavra Then para Então
\renewcommand{\algorithmicdo}{\textbf{fazer}} % Renomear a palavra Do para Fazer
% ----------

\newcommand{\quotes}[1]{``#1''} % Comando para inserir um texto dentro de aspas-dupla com espaço após o final

% FLUXOGRAMA: ----------
\usepackage{tikz}
\usetikzlibrary{positioning, shapes.geometric}
\tikzstyle{line}=[draw] % here

\tikzstyle{startstop} = [rectangle, rounded corners, minimum width=3cm,
minimum height=1cm,text centered, draw=black, fill=white!30]
\tikzstyle{arrow} = [thick, ->, >=stealth]
\tikzstyle{process} = [rectangle, minimum width=3cm, minimum height=1cm, text centered, draw=black, fill=white!30]
\tikzstyle{io} = [trapezium, trapezium left angle=70, trapezium right angle=110, minimum width=3cm, minimum height=1cm, text centered, draw=black, fill=white!30]
\tikzstyle{connector} = [draw, circle, fill=white!20]
\tikzstyle{decision} = [diamond, minimum width=3cm, minimum height=1cm, text centered, draw=black, fill=white!30, aspect=3]
% ----------


% CAIXA ENTORNO DAS FÓRMULAS: ----------
\usepackage{dashbox}
\newcommand\dboxed[1]{\dbox{\ensuremath{#1}}}

% ----------

%\usepackage{tabularx} % para colunas com largura ajustável

% -------- CONFIGURAÇÃO DE MARGENS --------
% \usepackage[top=2cm,bottom=2cm,left=3cm,right=2cm]{geometry}

% -------- PACOTE fancyhdr --------
\usepackage{fancyhdr}

\setlength{\marginparwidth}{2cm}

% Ajuste do tamanho do cabeçalho para acompanhar exatamente \textwidth
\setlength{\headwidth}{\textwidth}

% Define um headheight mínimo (para evitar avisos)
\setlength{\headheight}{15pt} 

% Limpa configurações padrão de cabeçalho/rodapé
\pagestyle{fancy}
\fancyhf{} 

% Marcações automáticas de seção para o cabeçalho (opcional)
\renewcommand{\sectionmark}[1]{%
	\markboth{\thesection\ #1}{}%
}

% Texto do cabeçalho à esquerda: seção e título
% Usamos \leftmark pois o \sectionmark acima grava no \leftmark
\fancyhead[L]{\small \leftmark}

% Texto do cabeçalho à direita: número de página
\fancyhead[R]{\small \thepage}

% Linha horizontal no cabeçalho (espessura 0.4pt)
\renewcommand{\headrulewidth}{0.4pt}

% Sobrescreve o título padrão “References” para “Referências”
%\renewcommand{\refname}{Referências}
% Tradução de autoref para o português
\def\figureautorefname{Figura}
\def\tableautorefname{Tabela}
\def\equationautorefname{Equa\c{c}\~ao}

% \hypersetup{backref=true}

\ifdefined\backref
        \renewcommand*{\backref}[1]{}
        \renewcommand*{\backrefentrycount}[2]{#1\ifnum#2>1 ~(#2)\fi}
\fi

\DeclarePairedDelimiter{\ceil}{\lceil}{\rceil}

\onehalfspacing

\usepackage{pgfgantt}
  % OK
%++++++++++++++++++++++++++++++++++++++++++++++++++++++++++++++++++++++++++++++
% Identificacoes (se o trabalho for em inglês, insira os dados em inglês)
% Para entradas abreviadas de Professora (Profa.) em português escreva: Prof$^\textnormal{a}$.
%++++++++++++++++++++++++++++++++++++++++++++++++++++++++++++++++++++++++++++++
\course{Engenharia Aeronáutica e Mecânica} % Programa de PG ou Curso de Graduação
\area{Propulsão Aeroespacial e Energia} % Área de concentração na PG (Não utilizado no caso de TG)

% Autor do trabalho: Nome Sobrenome
\authorgender{masc} %sexo: masc ou fem
\author{\textbf{EDEN DENIS} Ferreira da Silva Lopes}{Santos}
\itaauthoraddress{Rua João Fonseca dos Santos, 20, Apto. 93}{12.230-088}{São José dos Campos--SP}

% Titulo da Tese/Dissertação
\title{Aplicação de Inteligência Artificial (AI) em Banco de Dados de Emissão de Motores}

% Orientador
\advisorgender{masc} % masc ou fem
\advisor{Prof.~Dr.}{Pedro Teixeira \textbf{LACAVA}}{ITA}

% Coorientador (Caso não haja coorientador, colocar ambas as variáveis \coadvisorgender e \coadvisor comentadas, com um % na frente)
\coadvisorgender{masc}
\coadvisor{}{}  % {Prof.~Dr.}{Carlos Alberto Rocha \textbf{PIMENTEL}}{ITA}

% Pró-reitor da Pós-graduação
\bossgender{fem}
\boss{Prof.~Dra.}{Emília \textbf{VILLANI}}

%Coordenador do curso no caso de TG
\bosscoursegender{masc}
\bosscourse{Prof.~Dr.}{ - } % masc ou fem

% Palavras-Chaves informadas pela Biblioteca -> utilizada na CIP
\kwcip{Engine taxonomy}
\kwcip{Emission classification}
\kwcip{Inteligência Artificial}
\kwcip{Artificial Intelligence}
\kwcip{Python}

% membros da banca examinadora
\examiner{Prof. Dr.}{Pedro Teixeira \textbf{LACAVA}}{President}{ITA}
\examiner{Prof. Dr.}{}{Internal Member}{ITA} % Vinícius \textbf{MALATESTA}}{Internal Member}{ITA}
\examiner{Prof. Dr.}{}{Internal Member}{ITA} % {Elisan dos Santos \textbf{MAGALHÃES}}{Internal Member}{ITA}
\examiner{Prof. Dr.}{}{Internal Member}{ITA} % {Carlos Alberto Rocha \textbf{PIMENTEL}}{External Member}{UFABC}
\examiner{Prof. Dr.}{}{Internal Member}{ITA} % {\textbf{KARL} Peter Burr}{External Member}{UFABC}

% Data da defesa (mês em maiúsculo, se trabalho em inglês, e minúsculo se trabalho em português) 
\date{14}{JULHO}{2025}

% Número CDU - (somente para TG)
\cdu{621.38}

% Folha de Registro do Documento
\newcommand{\advisorTitle}{Prof.~Dr.}
\newcommand{\advisorName}{Pedro Teixeira \textbf{LACAVA}}

\newcommand{\coadvisorTitle}{Prof.~Dr.}
\newcommand{\coadvisorName}{Carlos Alberto Rocha \textbf{PIMENTEL}}

\newcommand{\dateDay}{14}
\newcommand{\dateMonth}{JULHO}
\newcommand{\dateYear}{2025}

% Data de publicação e defesa (usada na folha de registro)
\newcommand{\frdDataPublicacao}{04 de agosto de 2025}
\newcommand{\frdNumeroRegistro}{DCTA/ITA/TD-029/2020}
\newcommand{\frdPalavrasAutor}{Big Data; Engine Emission; Emission Databank; Clustering; Artificial Intelligence}
\newcommand{\frdPalavrasResultado}{Taxonomia; Emissão de Motores; Planejamento de Dados; Inteligência Artificial; Processamento de Dados}
\newcommand{\frdApresentacao}{
	ITA, São José dos Campos. Curso de Mestrado. Programa de Pós-Graduação em Engenharia Aeronáutica e Mecânica. Área de Propulsão Aeroespacial e Energia.\\
	Orientador: \advisorTitle~\advisorName.\\
	Coorientador: \coadvisorTitle~\coadvisorName.\\
	Defesa em \dateDay/\dateMonth/\dateYear. Publicada em \frdDataPublicacao.
}

\newcommand{\PorcentagemDaLargura}{0.80}  % OK

\begin{document}


	\chapter{Introdução}

		Este capítulo apresenta o contexto da pesquisa, o problema de classificação de motores aeronáuticos e a motivação para aplicar métodos de agrupamento aos dados de emissões.
	
		\section{Introdução}
		
			% Eden:
			O problema do particionamento tem grande expressão também na área da Biologia, motivo pelo qual os cientistas criaram um ramo específico chamado de taxonomia, para essa área do conhecimento, a de organismos e plantas. Entretanto, tudo isso, teve início desde a criação das sociedades quando, Aristóteles por volta de $ 300 $ a. C. inicialmente separou, de forma generalista, os seres vivos das plantas e, a \textit{posteriori}, para o caso do primeiro grupo, por exemplo, separou-os pela cor do sangue entre vermelho e não. No século XVIII, Carl von Linné ($ 1707 $ - $ 1778 $), aprimorou o processo taxonômico (utilizado até hoje) que organizou melhor os seres vivos, pois já se conhecia uma quantidade maior de espécies. Atualmente  particionado em cinco reinos: \textit{Monera}, \textit{Protista}, \textit{Fungi}, \textit{Animalia} e \textit{Plantae}; Linné dividiu os seres vivos levando em conta sete critérios (reino, filo, classe, ordem, família, gênero e espécie), o primeiro critério mais generalista e, a medida que se avança o critério, faz-se uma análise mais específica, logo, tem-se uma melhor classificação. Já que o objetivo deste trabalho é classificar as aeronaves, tanto civis como militares, o neologismo Aerotaxonomia torna-se bastante adequado, uma vez que, a palavra aero está relacionada à Engenharia aeroespacial e, por conseguinte, taxonomia significa classificação \cite{Linnaeus1758}, tendo em vista que um termo que seria conveniente para a biologia seria biotaxonomia. Recentemente, foi adicionado um domínio antes do reino. Tendo em vista que a palavra aero (já relacionada) é abrangente, fica conveniente usar, para este caso, Aerotaxonomia de aeronaves (veículos) dado que, poderia ser também de motores (sistema propulsivo). 
			
			% Daqui para baixo quem escreveu foi o prof. Ricieri
			A Inteligência Artificial (AI), em particular suas duas técnicas ou métodos para  tratar dados, $ k $ -means e Dendrogramação, podem servir de modelo para o conhecimento moderno caracterizado pela sua limitada validade temporal. E uma Matemática como a do século XX, que contraria essa particularidade do mundo moderno, século XXI, sociedade liquida de Bauman \cite{Bauman2000}, é inútil para os propósitos de toda e qualquer instituição do saber científico-tecnológico atual e também, muito provável, do futuro.
			
			Para o cientista de dados, a Inteligência Artificial (AI) e suas técnicas é o que deve ser estudado e aplicado na Engenharia moderna, em particular na aeronáutica; não porque ofereça um campo novo para as investigações do conhecimento, como aconteceu com a Matemática clássica até então, mas sim porque seus resultados preditivos encorajam a elevar o \textit{status} da Inteligência Artificial (AI) como a principal disciplina que poderá governar as Engenharias, essas que deverão, obrigatoriamente, preocuparem-se com aspectos nunca antes pensados pelos engenheiros como preferências populares, planos de venda, opções de \textit{marketing}, percentuais de sucesso e subjetividades de formas do produto; isso tudo exigido também pelos produtos aeronáuticos.  
			
			%		Diante dessas exigências do mercado em relação a profissão de Engenharia não existe outro caminho se não o da Inteligência Artificial (AI), que possa ajustar os currículos universitários às demandas atuais. 
			
			Enquanto o engenheiro tradicional busca a todo custo impor aos fenômenos do mundo teorias que possam explicá-los completamente; o engenheiro de dados analisa-os como eles se apresentam a nós, com certo grau de ambiguidade. Essa abordagem nova, pragmática, que tolera possíveis erros ou desvios de previsão, do ponto de vista semiótico–abducional, alegoriza dois movimentos de predições, de mesma direção e sentidos inversos: dos modelos, com seus símbolos sintáxicos aos dados e com suas minerações computacionais aos modelos. Daí a essência da dissonância entre Engenharia clássica, primeiro sentido e de Engenharia de dados, segundo sentido, estar nas explicações teóricas completas e nos modelos computacionais limitados. 
			
			Para um engenheiro de dados, profissional de natureza imediatista, já que atua em uma sociedade cada vez mais liquefeita, a busca de um modelo perfeito, oficial e universalmente abrangente, proíbe-o de qualquer possibilidade de predizer com liberdade, possíveis resultados minimamente assertivos. Assim, a essência do conflito entre esses dois profissionais das exatas deslocados no tempo, resvala na diversidade do conhecimento produzido nos inúmeros setores dessa ou daquela Engenharia, que exigem respostas rápidas e objetivas, mas não extremamente precisas. 
			
			Desta forma, sabendo que a Inteligência Artificial (AI) enquanto Matemática que tem influenciado como nunca antes as Engenharias, permitindo inúmeras respostas, resta àquele que se diz engenheiro ou cientista de dados saber quais são as perguntas. 
			
			% Daqui para baixo quem escreveu foi o Eden
			Portanto, a partir desta Matemática moderna, ajustada a sociedade atual e que tem definido novos padrões de solução somado à taxonomia, o principal objetivo deste trabalho é descrever um processo para particionar uma população $ d $-dimensional de aeronaves em $ K $-sets, sejam estas civis ou militares. A modelagem Matemática utilizada é chamada de $ k $ -means que, trata-se de uma aproximação eficiente quando há variações de determinadas características entre as classes, porém, reduzindo uma possível subjetividade humana. Apesar do interesse ser minimizar possíveis subjetividades na classificação das aeronaves, deve-se corroborar a análise Matemática com a prática da experiência computacional. O método escolhido para particionar, $ k $ -means, é um processo de fácil desenvolvimento e computacionalmente econômico, portanto, factível sobretudo pelo fato de não exigir investimento de altíssimos custos. 
			
			Além de toda a aplicação envolvida, o método $ k $ -means que, originalmente foi concebido como uma tentativa de encontrar um procedimento computacional factível para classificação otimizada de problemas da biologia e propriamente da computação \cite{MacQueen1967}, também é teoricamente pertinente, o que levou a escolhê-lo como ferramenta Matemática deste trabalho. Este método é uma generalização para a média ordinária da amostra, logo, é possível analisar o comportamento assintótico pertinente as características das classes e determinar algum tipo de lei baseada na Inteligência Artificial (AI) das aeronaves quando analisa-se o $ K $ do $ k $ -means para um número grande. O modelo matemático $ k $ -means é desenvolvido na seção (\ref{subsec:kmeans}) e os resultados obtidos sobre o comportamento assintótico são descritos na seção (\ref{sec:ResultadosEDiscussoes}).
			
			Em geral, o método não converge para uma solução ótima global, embora possivelmente existam casos especiais onde isso acontece. Entretanto, até onde se sabe, não existe método factível que sempre se obtêm solução global. Sendo assim, garante-se que pelo menos um resultado falibalista será encontrado, solução ótima local \cite{MacQueen1967}.
	
\end{document}

	\documentclass[12pt]{article} % Tamanho de letra 12, papel A4, Estilo de artigo

\usepackage[utf8]{inputenc} % Pacote para permitir o uso de acentos
% Ajustes de hifenização
\tolerance=1
\emergencystretch=\maxdimen
\hyphenpenalty=10000
\hbadness=10000
\usepackage{lipsum} % Pacote para gerar texto aleatório
\usepackage{ragged2e} % Pacote para mais opções de alinhamento
\usepackage{graphicx} % Pacote para inserir figuras jpeg, pdf e png
\usepackage{ae}
\usepackage{epsfig}
\usepackage{mathtools}
\usepackage[export]{adjustbox}
\usepackage{lineno}
\usepackage{amsthm}
\usepackage{commath}
\usepackage{color, colortbl}
\usepackage{bbm}
\usepackage{setspace} % Pacote para ajustar espaçamento entre linhas

% Instructions to prevent LaTeX from hyphenating words
\tolerance=1
\emergencystretch=\maxdimen
\hyphenpenalty=10000
\hbadness=10000

\graphicspath{{figures/}} % Indica o diretório no qual as imagens estão sendo salvas, evitando que precisamos digitar parte do caminho da pasta que contém as imagens
\usepackage[rightcaption]{sidecap} % Pacote para mais opções de figura, colocar, título da legenda da figura ao lado direito
\usepackage{url} % Pacote para citar referências de sites
\usepackage{makeidx} % Pacote para criar índice
\makeindex % Habilitar a criação do índice
\usepackage{longtable} % Pacote para usar tabelas longas

% ----------
% Pacotes para dar quebra de linha dentro da tabela:
%\usepackage{fourier} 
%\usepackage{array}
\usepackage{makecell}

\renewcommand\theadalign{bc}
\renewcommand\theadfont{\bfseries}
\renewcommand\theadgape{\Gape[4pt]}
\renewcommand\cellgape{\Gape[4pt]}
% ----------

\usepackage{wrapfig} % Pacote para alinhar figuras à esquerda ou à direita do texto
\usepackage{ae}
\usepackage{epsfig}
\usepackage{mathtools}
\usepackage[export]{adjustbox}
\usepackage{lineno}
\usepackage{amsfonts}
\usepackage{amsthm}
\usepackage{commath}
\usepackage{color, colortbl}
\usepackage{bbm}

% ----------
% Pacote para inserir cabeçalho e rodapé:
% Atenção! Se remover o cabeçalho e rodapé, o LaTeX coloca como cabeçalho o padrão dele, que é bem interessante também!
% \usepackage{fancyhdr} % Pacote para inserir cabeçalho e rodapé
% \pagestyle{fancy} % Estilo do cabeçalho e rodapé, neste caso apenas uma linha no cabeçalho
%\lhead{{\footnotesize Prandiano\textregistered}} % Inserir texto no canto esquerdo do cabeçalho
%\chead{centro} % Inserir texto no centro do cabeçalho
%\rhead{{\footnotesize Curso de Python}} % Inserir texto no canto direito do cabeçalho
%\lfoot{{\footnotesize Copyright \textcopyright 2018 Prandiano}} % Inserir texto no canto esquerdo do rodapé
%\cfoot{centro} % Inserir texto no centro do rodapé
%\rfoot{{\footnotesize Reprodução proibida: Lei 9.610 de 19/02/1998}} % Inserir texto no canto direito do rodapé
% O comando \renewcommand re-nova o parâmetro especificado
%\renewcommand{\footrulewidth}{0.4pt} % Inserir uma linha acima do rodapé
% ----------

% ----------
% Pacote para inserir múltiplas figuras:
\usepackage{subfigure} % Pacote para inserir múltiplas figuras
% \usepackage{subcaption} % Pacote incompatível com subfigure
%\usepackage{caption}
% ----------

\usepackage{multicol} % Pacote que permite usar multicolunas
\usepackage{float} % Pacote para usar o "H" do figure
\usepackage[dvipsnames]{xcolor} % Pacote de cores para textos e possivelmente outros
\usepackage{amsmath} % Pacote para usar o cases
\usepackage[makeroom]{cancel} % Pacote para usar o cancelar (em e embaixo)
\usepackage{enumitem} % Pacote para enumerar itens
\usepackage{amssymb} % Pacote para inserir os símbolos dos conjuntos numéricos N, Z, Q, R, C
\usepackage[T1]{fontenc}
\usepackage{enumitem}
\usepackage{textcomp}
\usepackage{pdfpages} % Pacote para inserir página em PDF
\usepackage{courier} % Pacote para usar a fonte courier
%\usepackage{matlab-prettifier}
\usepackage{lscape} % Pacote para mudar sentido da página
\newtheorem{theorem}{Teorema} % Pacote para inserir Teorema
\newtheorem{corollary}{Corollary} % Pacote para inserir Corolário
\newtheorem{remark}{Remark}
\newtheorem{mytheorem}{Definition}
\usepackage{hyperref} % Pacote para inserir hyperref
\usepackage{multirow} % Pacote para usar os comandos \multirow
\usepackage{pdfpages} % Pacote para incluir PDF's
\usepackage{listingsutf8} % Pacote para inserir arquivos externos tipo .txt
\usepackage{etoolbox}
\AtBeginEnvironment{verbatim}{\footnotesize}

\usepackage{standalone} % Pacote para inserir sub-arquivos LaTeX
\usepackage{bookmark} % Pacote para NÃO tratar o Índice como Apêndice e colocar letra

% Nome da listagem em português
\renewcommand{\lstlistingname}{Listagem}

% REFERÊNCIAS: ----------
%\usepackage[numbers, compress]{natbib} % Pacote para usar as referências bibliográficas
%\usepackage{usebib}
% \bibliographystyle{references} % Estilo das referências
%\setcitestyle{authoryear,aysep={},open={(},close={)}}
%\newcommand{\myand}{e}

% ----------


% PARÂMETROS DA FUNÇÃO LISTINGS: ----------
\lstset
{
    basicstyle=\small\ttfamily,
    commentstyle=\color{Green},
    keywordstyle=\color{Cerulean},
    frame=single,
    language=Python,
    morekeywords={True, False},
    numbers=left,
    numbersep=10pt,
    numberstyle=\footnotesize\color{Gray},
    showstringspaces=false,
    stringstyle=\color{Mulberry},
    tabsize=3,
    breaklines=true, % Este comando é para quebrar uma linha comprida
    postbreak=\mbox{\textcolor{red}{$\hookrightarrow$}\space}, % Este comando é para indicar com uma seta vermelha que a linha continua embaixo 
    extendedchars=true,
    literate=
        {á}{{\'a}}1 {Á}{{\'A}}1 {à}{{\`a}}1 {À}{{\`A}}1 {ã}{{\~a}}1 {Ã}{{\~A}}1 {â}{{\^a}}1 {Â}{{\^A}}1
        {ç}{{\c{c}}}1 {Ç}{{\c{C}}}1 {é}{{\'e}}1 {É}{{\'E}}1 {ê}{{\^e}}1 {Ê}{{\^E}}1 {í}{{\'\i}}1 {Í}{{\'I}}1
        {ó}{{\'o}}1 {Ó}{{\'O}}1 {ô}{{\^o}}1 {Ô}{{\^O}}1 {õ}{{\~o}}1 {Õ}{{\~O}}1 {ú}{{\'u}}1 {Ú}{{\'U}}1
        {ü}{{\"u}}1 {Ü}{{\"U}}1 {“}{{``}}1 {”}{{''}}1 {’}{{'}}1 {—}{{--}}1 {–}{{--}}1 {−}{{-}}1
        {±}{{$\pm$}}1 {σ}{{$\sigma$}}1 {→}{{$\to$}}1 {…}{{\ldots}}1
}

\lstset{inputencoding=utf8/latin1}

\lstdefinelanguage{Matlab}{
	language=Matlab,
	basicstyle=\ttfamily\small,
	keywordstyle=\color{blue}\bfseries,
	stringstyle=\color{red},
	commentstyle=\color{gray}\itshape,
	morekeywords={function, end, if, else, for, while, switch, case, otherwise},
	sensitive=true
}
% ----------

% ----------
\usepackage{nomencl} % Pacote para criar nomenclatura e lista de nomenclatura
\makenomenclature % Criar a nomenclatura
\renewcommand{\nomname}{Lista de Nomenclaturas} % Renomear o título da nomenclatura
% ----------

\usepackage{amsmath} % Pacote para inserir matrix com chaves, entre outras coisas

\usepackage{booktabs} % Pacote para usar \toprule, \midrule e \bottomrule


% EXECUTAR CÁLCULOS A PARTIR DA LINHA DE COMANDO: ---------
%\usepackage{siunitx, xfp} % Pacote para executar cálculos a partir da linha de comando
%\usepackage{l3regex}%unless you have very new expl3
%%\sisetup{locale=DE}
%%\let\oldnum\num
%%\ExplSyntaxOn
%\def\num#1{\def\tmp{#1}\regex_replace_all:nnN{,([0-9])}{.\1}{\tmp}\oldnum{\tmp}}
%\ExplSyntaxOff
% ----------


%\pagecolor{black} % Alterar a cor da página
%\color{white} % Alterar a cor do texto

% ----------
% Pacotes para inserir algoritmos e pseudo-códigos:
\setlength{\marginparwidth}{2.5cm}
\usepackage[colorinlistoftodos]{todonotes}
\usepackage{algorithm} % Pacote para inserir algoritmos
\usepackage{algpseudocode} % Pacote para inserir algoritmos e pseudo-códigos
\floatname{algorithm}{Algoritmo} % Renomear a palavra Algorithm para Procedimento
\renewcommand{\algorithmicrequire}{\textbf{entrada:}} % Renomear a palavra Input para Entrada
\renewcommand{\algorithmicensure}{\textbf{saída:}} % Renomear a palavra Output para Saída
\renewcommand{\algorithmicwhile}{\textbf{enquanto}} % Renomear a palavra While para Enquanto
\renewcommand{\algorithmicif}{\textbf{se}} % Renomear a palavra If para Se
\renewcommand{\algorithmicend}{\textbf{fim}} % Renomear a palavra End Output para Fim do
\renewcommand{\algorithmicprocedure}{\textbf{procedimento}} % Renomear a palavra Procedure para Procedimento
\renewcommand{\algorithmicfor}{\textbf{para}} % Renomear a palavra For para Para
\renewcommand{\algorithmicelse}{\textbf{senão}} % Renomear a palavra Else para Senão
\renewcommand{\algorithmicfunction}{\textbf{função}} % Renomear a palavra Function para Função 
\renewcommand{\algorithmicthen}{\textbf{então}} % Renomear a palavra Then para Então
\renewcommand{\algorithmicdo}{\textbf{fazer}} % Renomear a palavra Do para Fazer
% ----------

\newcommand{\quotes}[1]{``#1''} % Comando para inserir um texto dentro de aspas-dupla com espaço após o final

% FLUXOGRAMA: ----------
\usepackage{tikz}
\usetikzlibrary{positioning, shapes.geometric}
\tikzstyle{line}=[draw] % here

\tikzstyle{startstop} = [rectangle, rounded corners, minimum width=3cm,
minimum height=1cm,text centered, draw=black, fill=white!30]
\tikzstyle{arrow} = [thick, ->, >=stealth]
\tikzstyle{process} = [rectangle, minimum width=3cm, minimum height=1cm, text centered, draw=black, fill=white!30]
\tikzstyle{io} = [trapezium, trapezium left angle=70, trapezium right angle=110, minimum width=3cm, minimum height=1cm, text centered, draw=black, fill=white!30]
\tikzstyle{connector} = [draw, circle, fill=white!20]
\tikzstyle{decision} = [diamond, minimum width=3cm, minimum height=1cm, text centered, draw=black, fill=white!30, aspect=3]
% ----------


% CAIXA ENTORNO DAS FÓRMULAS: ----------
\usepackage{dashbox}
\newcommand\dboxed[1]{\dbox{\ensuremath{#1}}}

% ----------

%\usepackage{tabularx} % para colunas com largura ajustável

% -------- CONFIGURAÇÃO DE MARGENS --------
% \usepackage[top=2cm,bottom=2cm,left=3cm,right=2cm]{geometry}

% -------- PACOTE fancyhdr --------
\usepackage{fancyhdr}

\setlength{\marginparwidth}{2cm}

% Ajuste do tamanho do cabeçalho para acompanhar exatamente \textwidth
\setlength{\headwidth}{\textwidth}

% Define um headheight mínimo (para evitar avisos)
\setlength{\headheight}{15pt} 

% Limpa configurações padrão de cabeçalho/rodapé
\pagestyle{fancy}
\fancyhf{} 

% Marcações automáticas de seção para o cabeçalho (opcional)
\renewcommand{\sectionmark}[1]{%
	\markboth{\thesection\ #1}{}%
}

% Texto do cabeçalho à esquerda: seção e título
% Usamos \leftmark pois o \sectionmark acima grava no \leftmark
\fancyhead[L]{\small \leftmark}

% Texto do cabeçalho à direita: número de página
\fancyhead[R]{\small \thepage}

% Linha horizontal no cabeçalho (espessura 0.4pt)
\renewcommand{\headrulewidth}{0.4pt}

% Sobrescreve o título padrão “References” para “Referências”
%\renewcommand{\refname}{Referências}
% Tradução de autoref para o português
\def\figureautorefname{Figura}
\def\tableautorefname{Tabela}
\def\equationautorefname{Equa\c{c}\~ao}

% \hypersetup{backref=true}

\ifdefined\backref
        \renewcommand*{\backref}[1]{}
        \renewcommand*{\backrefentrycount}[2]{#1\ifnum#2>1 ~(#2)\fi}
\fi

\DeclarePairedDelimiter{\ceil}{\lceil}{\rceil}

\onehalfspacing

\usepackage{pgfgantt}

%++++++++++++++++++++++++++++++++++++++++++++++++++++++++++++++++++++++++++++++
% Identificacoes (se o trabalho for em inglês, insira os dados em inglês)
% Para entradas abreviadas de Professora (Profa.) em português escreva: Prof$^\textnormal{a}$.
%++++++++++++++++++++++++++++++++++++++++++++++++++++++++++++++++++++++++++++++
\course{Engenharia Aeronáutica e Mecânica} % Programa de PG ou Curso de Graduação
\area{Propulsão Aeroespacial e Energia} % Área de concentração na PG (Não utilizado no caso de TG)

% Autor do trabalho: Nome Sobrenome
\authorgender{masc} %sexo: masc ou fem
\author{\textbf{EDEN DENIS} Ferreira da Silva Lopes}{Santos}
\itaauthoraddress{Rua João Fonseca dos Santos, 20, Apto. 93}{12.230-088}{São José dos Campos--SP}

% Titulo da Tese/Dissertação
\title{Aplicação de Inteligência Artificial (AI) em Banco de Dados de Emissão de Motores}

% Orientador
\advisorgender{masc} % masc ou fem
\advisor{Prof.~Dr.}{Pedro Teixeira \textbf{LACAVA}}{ITA}

% Coorientador (Caso não haja coorientador, colocar ambas as variáveis \coadvisorgender e \coadvisor comentadas, com um % na frente)
\coadvisorgender{masc}
\coadvisor{}{}  % {Prof.~Dr.}{Carlos Alberto Rocha \textbf{PIMENTEL}}{ITA}

% Pró-reitor da Pós-graduação
\bossgender{fem}
\boss{Prof.~Dra.}{Emília \textbf{VILLANI}}

%Coordenador do curso no caso de TG
\bosscoursegender{masc}
\bosscourse{Prof.~Dr.}{ - } % masc ou fem

% Palavras-Chaves informadas pela Biblioteca -> utilizada na CIP
\kwcip{Engine taxonomy}
\kwcip{Emission classification}
\kwcip{Inteligência Artificial}
\kwcip{Artificial Intelligence}
\kwcip{Python}

% membros da banca examinadora
\examiner{Prof. Dr.}{Pedro Teixeira \textbf{LACAVA}}{President}{ITA}
\examiner{Prof. Dr.}{}{Internal Member}{ITA} % Vinícius \textbf{MALATESTA}}{Internal Member}{ITA}
\examiner{Prof. Dr.}{}{Internal Member}{ITA} % {Elisan dos Santos \textbf{MAGALHÃES}}{Internal Member}{ITA}
\examiner{Prof. Dr.}{}{Internal Member}{ITA} % {Carlos Alberto Rocha \textbf{PIMENTEL}}{External Member}{UFABC}
\examiner{Prof. Dr.}{}{Internal Member}{ITA} % {\textbf{KARL} Peter Burr}{External Member}{UFABC}

% Data da defesa (mês em maiúsculo, se trabalho em inglês, e minúsculo se trabalho em português) 
\date{14}{JULHO}{2025}

% Número CDU - (somente para TG)
\cdu{621.38}

% Folha de Registro do Documento
\newcommand{\advisorTitle}{Prof.~Dr.}
\newcommand{\advisorName}{Pedro Teixeira \textbf{LACAVA}}

\newcommand{\coadvisorTitle}{Prof.~Dr.}
\newcommand{\coadvisorName}{Carlos Alberto Rocha \textbf{PIMENTEL}}

\newcommand{\dateDay}{14}
\newcommand{\dateMonth}{JULHO}
\newcommand{\dateYear}{2025}

% Data de publicação e defesa (usada na folha de registro)
\newcommand{\frdDataPublicacao}{04 de agosto de 2025}
\newcommand{\frdNumeroRegistro}{DCTA/ITA/TD-029/2020}
\newcommand{\frdPalavrasAutor}{Big Data; Engine Emission; Emission Databank; Clustering; Artificial Intelligence}
\newcommand{\frdPalavrasResultado}{Taxonomia; Emissão de Motores; Planejamento de Dados; Inteligência Artificial; Processamento de Dados}
\newcommand{\frdApresentacao}{
	ITA, São José dos Campos. Curso de Mestrado. Programa de Pós-Graduação em Engenharia Aeronáutica e Mecânica. Área de Propulsão Aeroespacial e Energia.\\
	Orientador: \advisorTitle~\advisorName.\\
	Coorientador: \coadvisorTitle~\coadvisorName.\\
	Defesa em \dateDay/\dateMonth/\dateYear. Publicada em \frdDataPublicacao.
}

\newcommand{\PorcentagemDaLargura}{0.80}

\begin{document}
	
	\section{Contextualização do Problema}
	
		A proposta deste trabalho é classificar a partir de Inteligência Artificial (AI) um Big Data de motores de aeronaves, sobretudo pela emissão dos gases de saída em possíveis grupos, sendo que cada grupo deve conter motores de aeronaves com alguma semelhança categórica segundo suas características de projeto: condições de pressão, temperatura, dimensões etc..
		
		Sendo assim, com uma Taxonomia Numérica \cite{Manktelow2010} ou Análise Computacional de Clustering (Data Mining) \cite{Anderberg1973}, escolhe-se a \textit{priori}, o número de grupos ($ k $) que deverão dendrogramar os motores de aeronaves de pouca, média ou muita emissão; o que servirá como insumo informativo para pré-projetos de novas motores aeronaves.
	
	\section{Identificação do Problema de Pesquisa}	
	
		Os problemas da pesquisa podem ser listados, conforme abaixo:
		
		\begin{itemize}
			
			\item \textbf{Classificar os motores de aeronaves por grupos de emissões pela suas similaridades o que explicita e ranqueia os projetos}: classificar por ordem do melhor para o \quotes{priori}, neste caso, pelas emissões dos motores de aeronaves;
			
			\item \textbf{Determinar qual o número de grupos que os motores de aeronaves devem ser classificadas para obter a melhor ranque}: destacar o número de segregações (2, 3, 4, ..., $ N $) os motores das aeronaves devem ser separados;
			
			\item \textbf{Determinar quais as melhores combinações das características de projeto de motores, sobretudo pelas suas emissões, para um desenvolvimento (ou ajuste) de um projeto, segue as características}:
			
			\input{tables/heatmap_variables_contextualization.tex}
			
			As variáveis foram numeradas, pois quando são gerados os mapas de calor, tanto para as correlações quanto para os $ p $-values, foi trocado o nome da variável pelo índice, \autoref{sec:mapadecalorfiltrado} e \autoref{sec:pvaluesfiltrado}.
			
			\item \textbf{Determinar qual são os projeto ímpares/únicos}: destacar qual o(s) melhor(es) motores (alvo) bem como o(s) piore(s) impróprios;
			
			\item \textbf{Determinar qual a melhor técnica ou combinação a ser aplicada}: aplicar várias técnicas matemáticas -- possivelmente usando o processamento paralelo -- para indicar qual a melhor entre elas, modelos competidores.
			
%			\item \textbf{Desenvolver uma aplicação que possa ser executada em qualquer plataforma}: a(s) aplicação(ões) computacionais devem ser pensadas para ser(em) executadas nas diversas plataformas (\texttt{Windows}, \texttt{Linux}, \texttt{MacOs} etc.);
			
%			\item \textbf{Direcionar o estudo dos subsistemas dos motores de aeronaves (propulsivo, estrutural etc.)}: indicar o sistema de estudo propulsivo, foco da pesquisa em uma segunda etapa.
			
		\end{itemize}
		
\end{document}

	
	% Descrição numérica:
	\documentclass[12pt]{article} % Tamanho de letra 12, papel A4, Estilo de artigo

\usepackage[utf8]{inputenc} % Pacote para permitir o uso de acentos
% Ajustes de hifenização
\tolerance=1
\emergencystretch=\maxdimen
\hyphenpenalty=10000
\hbadness=10000
\usepackage{lipsum} % Pacote para gerar texto aleatório
\usepackage{ragged2e} % Pacote para mais opções de alinhamento
\usepackage{graphicx} % Pacote para inserir figuras jpeg, pdf e png
\usepackage{ae}
\usepackage{epsfig}
\usepackage{mathtools}
\usepackage[export]{adjustbox}
\usepackage{lineno}
\usepackage{amsthm}
\usepackage{commath}
\usepackage{color, colortbl}
\usepackage{bbm}
\usepackage{setspace} % Pacote para ajustar espaçamento entre linhas

% Instructions to prevent LaTeX from hyphenating words
\tolerance=1
\emergencystretch=\maxdimen
\hyphenpenalty=10000
\hbadness=10000

\graphicspath{{figures/}} % Indica o diretório no qual as imagens estão sendo salvas, evitando que precisamos digitar parte do caminho da pasta que contém as imagens
\usepackage[rightcaption]{sidecap} % Pacote para mais opções de figura, colocar, título da legenda da figura ao lado direito
\usepackage{url} % Pacote para citar referências de sites
\usepackage{makeidx} % Pacote para criar índice
\makeindex % Habilitar a criação do índice
\usepackage{longtable} % Pacote para usar tabelas longas

% ----------
% Pacotes para dar quebra de linha dentro da tabela:
%\usepackage{fourier} 
%\usepackage{array}
\usepackage{makecell}

\renewcommand\theadalign{bc}
\renewcommand\theadfont{\bfseries}
\renewcommand\theadgape{\Gape[4pt]}
\renewcommand\cellgape{\Gape[4pt]}
% ----------

\usepackage{wrapfig} % Pacote para alinhar figuras à esquerda ou à direita do texto
\usepackage{ae}
\usepackage{epsfig}
\usepackage{mathtools}
\usepackage[export]{adjustbox}
\usepackage{lineno}
\usepackage{amsfonts}
\usepackage{amsthm}
\usepackage{commath}
\usepackage{color, colortbl}
\usepackage{bbm}

% ----------
% Pacote para inserir cabeçalho e rodapé:
% Atenção! Se remover o cabeçalho e rodapé, o LaTeX coloca como cabeçalho o padrão dele, que é bem interessante também!
% \usepackage{fancyhdr} % Pacote para inserir cabeçalho e rodapé
% \pagestyle{fancy} % Estilo do cabeçalho e rodapé, neste caso apenas uma linha no cabeçalho
%\lhead{{\footnotesize Prandiano\textregistered}} % Inserir texto no canto esquerdo do cabeçalho
%\chead{centro} % Inserir texto no centro do cabeçalho
%\rhead{{\footnotesize Curso de Python}} % Inserir texto no canto direito do cabeçalho
%\lfoot{{\footnotesize Copyright \textcopyright 2018 Prandiano}} % Inserir texto no canto esquerdo do rodapé
%\cfoot{centro} % Inserir texto no centro do rodapé
%\rfoot{{\footnotesize Reprodução proibida: Lei 9.610 de 19/02/1998}} % Inserir texto no canto direito do rodapé
% O comando \renewcommand re-nova o parâmetro especificado
%\renewcommand{\footrulewidth}{0.4pt} % Inserir uma linha acima do rodapé
% ----------

% ----------
% Pacote para inserir múltiplas figuras:
\usepackage{subfigure} % Pacote para inserir múltiplas figuras
% \usepackage{subcaption} % Pacote incompatível com subfigure
%\usepackage{caption}
% ----------

\usepackage{multicol} % Pacote que permite usar multicolunas
\usepackage{float} % Pacote para usar o "H" do figure
\usepackage[dvipsnames]{xcolor} % Pacote de cores para textos e possivelmente outros
\usepackage{amsmath} % Pacote para usar o cases
\usepackage[makeroom]{cancel} % Pacote para usar o cancelar (em e embaixo)
\usepackage{enumitem} % Pacote para enumerar itens
\usepackage{amssymb} % Pacote para inserir os símbolos dos conjuntos numéricos N, Z, Q, R, C
\usepackage[T1]{fontenc}
\usepackage{enumitem}
\usepackage{textcomp}
\usepackage{pdfpages} % Pacote para inserir página em PDF
\usepackage{courier} % Pacote para usar a fonte courier
%\usepackage{matlab-prettifier}
\usepackage{lscape} % Pacote para mudar sentido da página
\newtheorem{theorem}{Teorema} % Pacote para inserir Teorema
\newtheorem{corollary}{Corollary} % Pacote para inserir Corolário
\newtheorem{remark}{Remark}
\newtheorem{mytheorem}{Definition}
\usepackage{hyperref} % Pacote para inserir hyperref
\usepackage{multirow} % Pacote para usar os comandos \multirow
\usepackage{pdfpages} % Pacote para incluir PDF's
\usepackage{listingsutf8} % Pacote para inserir arquivos externos tipo .txt
\usepackage{etoolbox}
\AtBeginEnvironment{verbatim}{\footnotesize}

\usepackage{standalone} % Pacote para inserir sub-arquivos LaTeX
\usepackage{bookmark} % Pacote para NÃO tratar o Índice como Apêndice e colocar letra

% Nome da listagem em português
\renewcommand{\lstlistingname}{Listagem}

% REFERÊNCIAS: ----------
%\usepackage[numbers, compress]{natbib} % Pacote para usar as referências bibliográficas
%\usepackage{usebib}
% \bibliographystyle{references} % Estilo das referências
%\setcitestyle{authoryear,aysep={},open={(},close={)}}
%\newcommand{\myand}{e}

% ----------


% PARÂMETROS DA FUNÇÃO LISTINGS: ----------
\lstset
{
    basicstyle=\small\ttfamily,
    commentstyle=\color{Green},
    keywordstyle=\color{Cerulean},
    frame=single,
    language=Python,
    morekeywords={True, False},
    numbers=left,
    numbersep=10pt,
    numberstyle=\footnotesize\color{Gray},
    showstringspaces=false,
    stringstyle=\color{Mulberry},
    tabsize=3,
    breaklines=true, % Este comando é para quebrar uma linha comprida
    postbreak=\mbox{\textcolor{red}{$\hookrightarrow$}\space}, % Este comando é para indicar com uma seta vermelha que a linha continua embaixo 
    extendedchars=true,
    literate=
        {á}{{\'a}}1 {Á}{{\'A}}1 {à}{{\`a}}1 {À}{{\`A}}1 {ã}{{\~a}}1 {Ã}{{\~A}}1 {â}{{\^a}}1 {Â}{{\^A}}1
        {ç}{{\c{c}}}1 {Ç}{{\c{C}}}1 {é}{{\'e}}1 {É}{{\'E}}1 {ê}{{\^e}}1 {Ê}{{\^E}}1 {í}{{\'\i}}1 {Í}{{\'I}}1
        {ó}{{\'o}}1 {Ó}{{\'O}}1 {ô}{{\^o}}1 {Ô}{{\^O}}1 {õ}{{\~o}}1 {Õ}{{\~O}}1 {ú}{{\'u}}1 {Ú}{{\'U}}1
        {ü}{{\"u}}1 {Ü}{{\"U}}1 {“}{{``}}1 {”}{{''}}1 {’}{{'}}1 {—}{{--}}1 {–}{{--}}1 {−}{{-}}1
        {±}{{$\pm$}}1 {σ}{{$\sigma$}}1 {→}{{$\to$}}1 {…}{{\ldots}}1
}

\lstset{inputencoding=utf8/latin1}

\lstdefinelanguage{Matlab}{
	language=Matlab,
	basicstyle=\ttfamily\small,
	keywordstyle=\color{blue}\bfseries,
	stringstyle=\color{red},
	commentstyle=\color{gray}\itshape,
	morekeywords={function, end, if, else, for, while, switch, case, otherwise},
	sensitive=true
}
% ----------

% ----------
\usepackage{nomencl} % Pacote para criar nomenclatura e lista de nomenclatura
\makenomenclature % Criar a nomenclatura
\renewcommand{\nomname}{Lista de Nomenclaturas} % Renomear o título da nomenclatura
% ----------

\usepackage{amsmath} % Pacote para inserir matrix com chaves, entre outras coisas

\usepackage{booktabs} % Pacote para usar \toprule, \midrule e \bottomrule


% EXECUTAR CÁLCULOS A PARTIR DA LINHA DE COMANDO: ---------
%\usepackage{siunitx, xfp} % Pacote para executar cálculos a partir da linha de comando
%\usepackage{l3regex}%unless you have very new expl3
%%\sisetup{locale=DE}
%%\let\oldnum\num
%%\ExplSyntaxOn
%\def\num#1{\def\tmp{#1}\regex_replace_all:nnN{,([0-9])}{.\1}{\tmp}\oldnum{\tmp}}
%\ExplSyntaxOff
% ----------


%\pagecolor{black} % Alterar a cor da página
%\color{white} % Alterar a cor do texto

% ----------
% Pacotes para inserir algoritmos e pseudo-códigos:
\setlength{\marginparwidth}{2.5cm}
\usepackage[colorinlistoftodos]{todonotes}
\usepackage{algorithm} % Pacote para inserir algoritmos
\usepackage{algpseudocode} % Pacote para inserir algoritmos e pseudo-códigos
\floatname{algorithm}{Algoritmo} % Renomear a palavra Algorithm para Procedimento
\renewcommand{\algorithmicrequire}{\textbf{entrada:}} % Renomear a palavra Input para Entrada
\renewcommand{\algorithmicensure}{\textbf{saída:}} % Renomear a palavra Output para Saída
\renewcommand{\algorithmicwhile}{\textbf{enquanto}} % Renomear a palavra While para Enquanto
\renewcommand{\algorithmicif}{\textbf{se}} % Renomear a palavra If para Se
\renewcommand{\algorithmicend}{\textbf{fim}} % Renomear a palavra End Output para Fim do
\renewcommand{\algorithmicprocedure}{\textbf{procedimento}} % Renomear a palavra Procedure para Procedimento
\renewcommand{\algorithmicfor}{\textbf{para}} % Renomear a palavra For para Para
\renewcommand{\algorithmicelse}{\textbf{senão}} % Renomear a palavra Else para Senão
\renewcommand{\algorithmicfunction}{\textbf{função}} % Renomear a palavra Function para Função 
\renewcommand{\algorithmicthen}{\textbf{então}} % Renomear a palavra Then para Então
\renewcommand{\algorithmicdo}{\textbf{fazer}} % Renomear a palavra Do para Fazer
% ----------

\newcommand{\quotes}[1]{``#1''} % Comando para inserir um texto dentro de aspas-dupla com espaço após o final

% FLUXOGRAMA: ----------
\usepackage{tikz}
\usetikzlibrary{positioning, shapes.geometric}
\tikzstyle{line}=[draw] % here

\tikzstyle{startstop} = [rectangle, rounded corners, minimum width=3cm,
minimum height=1cm,text centered, draw=black, fill=white!30]
\tikzstyle{arrow} = [thick, ->, >=stealth]
\tikzstyle{process} = [rectangle, minimum width=3cm, minimum height=1cm, text centered, draw=black, fill=white!30]
\tikzstyle{io} = [trapezium, trapezium left angle=70, trapezium right angle=110, minimum width=3cm, minimum height=1cm, text centered, draw=black, fill=white!30]
\tikzstyle{connector} = [draw, circle, fill=white!20]
\tikzstyle{decision} = [diamond, minimum width=3cm, minimum height=1cm, text centered, draw=black, fill=white!30, aspect=3]
% ----------


% CAIXA ENTORNO DAS FÓRMULAS: ----------
\usepackage{dashbox}
\newcommand\dboxed[1]{\dbox{\ensuremath{#1}}}

% ----------

%\usepackage{tabularx} % para colunas com largura ajustável

% -------- CONFIGURAÇÃO DE MARGENS --------
% \usepackage[top=2cm,bottom=2cm,left=3cm,right=2cm]{geometry}

% -------- PACOTE fancyhdr --------
\usepackage{fancyhdr}

\setlength{\marginparwidth}{2cm}

% Ajuste do tamanho do cabeçalho para acompanhar exatamente \textwidth
\setlength{\headwidth}{\textwidth}

% Define um headheight mínimo (para evitar avisos)
\setlength{\headheight}{15pt} 

% Limpa configurações padrão de cabeçalho/rodapé
\pagestyle{fancy}
\fancyhf{} 

% Marcações automáticas de seção para o cabeçalho (opcional)
\renewcommand{\sectionmark}[1]{%
	\markboth{\thesection\ #1}{}%
}

% Texto do cabeçalho à esquerda: seção e título
% Usamos \leftmark pois o \sectionmark acima grava no \leftmark
\fancyhead[L]{\small \leftmark}

% Texto do cabeçalho à direita: número de página
\fancyhead[R]{\small \thepage}

% Linha horizontal no cabeçalho (espessura 0.4pt)
\renewcommand{\headrulewidth}{0.4pt}

% Sobrescreve o título padrão “References” para “Referências”
%\renewcommand{\refname}{Referências}
% Tradução de autoref para o português
\def\figureautorefname{Figura}
\def\tableautorefname{Tabela}
\def\equationautorefname{Equa\c{c}\~ao}

% \hypersetup{backref=true}

\ifdefined\backref
        \renewcommand*{\backref}[1]{}
        \renewcommand*{\backrefentrycount}[2]{#1\ifnum#2>1 ~(#2)\fi}
\fi

\DeclarePairedDelimiter{\ceil}{\lceil}{\rceil}

\onehalfspacing

\usepackage{pgfgantt}

%++++++++++++++++++++++++++++++++++++++++++++++++++++++++++++++++++++++++++++++
% Identificacoes (se o trabalho for em inglês, insira os dados em inglês)
% Para entradas abreviadas de Professora (Profa.) em português escreva: Prof$^\textnormal{a}$.
%++++++++++++++++++++++++++++++++++++++++++++++++++++++++++++++++++++++++++++++
\course{Engenharia Aeronáutica e Mecânica} % Programa de PG ou Curso de Graduação
\area{Propulsão Aeroespacial e Energia} % Área de concentração na PG (Não utilizado no caso de TG)

% Autor do trabalho: Nome Sobrenome
\authorgender{masc} %sexo: masc ou fem
\author{\textbf{EDEN DENIS} Ferreira da Silva Lopes}{Santos}
\itaauthoraddress{Rua João Fonseca dos Santos, 20, Apto. 93}{12.230-088}{São José dos Campos--SP}

% Titulo da Tese/Dissertação
\title{Aplicação de Inteligência Artificial (AI) em Banco de Dados de Emissão de Motores}

% Orientador
\advisorgender{masc} % masc ou fem
\advisor{Prof.~Dr.}{Pedro Teixeira \textbf{LACAVA}}{ITA}

% Coorientador (Caso não haja coorientador, colocar ambas as variáveis \coadvisorgender e \coadvisor comentadas, com um % na frente)
\coadvisorgender{masc}
\coadvisor{}{}  % {Prof.~Dr.}{Carlos Alberto Rocha \textbf{PIMENTEL}}{ITA}

% Pró-reitor da Pós-graduação
\bossgender{fem}
\boss{Prof.~Dra.}{Emília \textbf{VILLANI}}

%Coordenador do curso no caso de TG
\bosscoursegender{masc}
\bosscourse{Prof.~Dr.}{ - } % masc ou fem

% Palavras-Chaves informadas pela Biblioteca -> utilizada na CIP
\kwcip{Engine taxonomy}
\kwcip{Emission classification}
\kwcip{Inteligência Artificial}
\kwcip{Artificial Intelligence}
\kwcip{Python}

% membros da banca examinadora
\examiner{Prof. Dr.}{Pedro Teixeira \textbf{LACAVA}}{President}{ITA}
\examiner{Prof. Dr.}{}{Internal Member}{ITA} % Vinícius \textbf{MALATESTA}}{Internal Member}{ITA}
\examiner{Prof. Dr.}{}{Internal Member}{ITA} % {Elisan dos Santos \textbf{MAGALHÃES}}{Internal Member}{ITA}
\examiner{Prof. Dr.}{}{Internal Member}{ITA} % {Carlos Alberto Rocha \textbf{PIMENTEL}}{External Member}{UFABC}
\examiner{Prof. Dr.}{}{Internal Member}{ITA} % {\textbf{KARL} Peter Burr}{External Member}{UFABC}

% Data da defesa (mês em maiúsculo, se trabalho em inglês, e minúsculo se trabalho em português) 
\date{14}{JULHO}{2025}

% Número CDU - (somente para TG)
\cdu{621.38}

% Folha de Registro do Documento
\newcommand{\advisorTitle}{Prof.~Dr.}
\newcommand{\advisorName}{Pedro Teixeira \textbf{LACAVA}}

\newcommand{\coadvisorTitle}{Prof.~Dr.}
\newcommand{\coadvisorName}{Carlos Alberto Rocha \textbf{PIMENTEL}}

\newcommand{\dateDay}{14}
\newcommand{\dateMonth}{JULHO}
\newcommand{\dateYear}{2025}

% Data de publicação e defesa (usada na folha de registro)
\newcommand{\frdDataPublicacao}{04 de agosto de 2025}
\newcommand{\frdNumeroRegistro}{DCTA/ITA/TD-029/2020}
\newcommand{\frdPalavrasAutor}{Big Data; Engine Emission; Emission Databank; Clustering; Artificial Intelligence}
\newcommand{\frdPalavrasResultado}{Taxonomia; Emissão de Motores; Planejamento de Dados; Inteligência Artificial; Processamento de Dados}
\newcommand{\frdApresentacao}{
	ITA, São José dos Campos. Curso de Mestrado. Programa de Pós-Graduação em Engenharia Aeronáutica e Mecânica. Área de Propulsão Aeroespacial e Energia.\\
	Orientador: \advisorTitle~\advisorName.\\
	Coorientador: \coadvisorTitle~\coadvisorName.\\
	Defesa em \dateDay/\dateMonth/\dateYear. Publicada em \frdDataPublicacao.
}

\newcommand{\PorcentagemDaLargura}{0.80}

\begin{document}
	
	\newpage
	%\shipout\null % Inserir página em branco
	\chapter{Descrição Numérica}\label{chap:descricao_numerica}

		Este capítulo reúne a descrição do conjunto de dados, as etapas de preparação, as métricas de qualidade e os métodos de análise numérica aplicados.
	
		
		\input{audithas.tex}
		\input{praeparatio.tex}
		\input{correlsp_correlations.tex}
%		\input{dendrogramacao.tex}
		\input{k_means_method.tex}
		
		\input{python.tex}
	
				
\end{document}

	\input{resultados_e_discussoes.tex}
	
	% Conclusão e trabalhos futuros:
%	\documentclass[12pt]{article} % Tamanho de letra 12, papel A4, Estilo de artigo

\usepackage[utf8]{inputenc} % Pacote para permitir o uso de acentos
% Ajustes de hifenização
\tolerance=1
\emergencystretch=\maxdimen
\hyphenpenalty=10000
\hbadness=10000
\usepackage{lipsum} % Pacote para gerar texto aleatório
\usepackage{ragged2e} % Pacote para mais opções de alinhamento
\usepackage{graphicx} % Pacote para inserir figuras jpeg, pdf e png
\usepackage{ae}
\usepackage{epsfig}
\usepackage{mathtools}
\usepackage[export]{adjustbox}
\usepackage{lineno}
\usepackage{amsthm}
\usepackage{commath}
\usepackage{color, colortbl}
\usepackage{bbm}
\usepackage{setspace} % Pacote para ajustar espaçamento entre linhas

% Instructions to prevent LaTeX from hyphenating words
\tolerance=1
\emergencystretch=\maxdimen
\hyphenpenalty=10000
\hbadness=10000

\graphicspath{{figures/}} % Indica o diretório no qual as imagens estão sendo salvas, evitando que precisamos digitar parte do caminho da pasta que contém as imagens
\usepackage[rightcaption]{sidecap} % Pacote para mais opções de figura, colocar, título da legenda da figura ao lado direito
\usepackage{url} % Pacote para citar referências de sites
\usepackage{makeidx} % Pacote para criar índice
\makeindex % Habilitar a criação do índice
\usepackage{longtable} % Pacote para usar tabelas longas

% ----------
% Pacotes para dar quebra de linha dentro da tabela:
%\usepackage{fourier} 
%\usepackage{array}
\usepackage{makecell}

\renewcommand\theadalign{bc}
\renewcommand\theadfont{\bfseries}
\renewcommand\theadgape{\Gape[4pt]}
\renewcommand\cellgape{\Gape[4pt]}
% ----------

\usepackage{wrapfig} % Pacote para alinhar figuras à esquerda ou à direita do texto
\usepackage{ae}
\usepackage{epsfig}
\usepackage{mathtools}
\usepackage[export]{adjustbox}
\usepackage{lineno}
\usepackage{amsfonts}
\usepackage{amsthm}
\usepackage{commath}
\usepackage{color, colortbl}
\usepackage{bbm}

% ----------
% Pacote para inserir cabeçalho e rodapé:
% Atenção! Se remover o cabeçalho e rodapé, o LaTeX coloca como cabeçalho o padrão dele, que é bem interessante também!
% \usepackage{fancyhdr} % Pacote para inserir cabeçalho e rodapé
% \pagestyle{fancy} % Estilo do cabeçalho e rodapé, neste caso apenas uma linha no cabeçalho
%\lhead{{\footnotesize Prandiano\textregistered}} % Inserir texto no canto esquerdo do cabeçalho
%\chead{centro} % Inserir texto no centro do cabeçalho
%\rhead{{\footnotesize Curso de Python}} % Inserir texto no canto direito do cabeçalho
%\lfoot{{\footnotesize Copyright \textcopyright 2018 Prandiano}} % Inserir texto no canto esquerdo do rodapé
%\cfoot{centro} % Inserir texto no centro do rodapé
%\rfoot{{\footnotesize Reprodução proibida: Lei 9.610 de 19/02/1998}} % Inserir texto no canto direito do rodapé
% O comando \renewcommand re-nova o parâmetro especificado
%\renewcommand{\footrulewidth}{0.4pt} % Inserir uma linha acima do rodapé
% ----------

% ----------
% Pacote para inserir múltiplas figuras:
\usepackage{subfigure} % Pacote para inserir múltiplas figuras
% \usepackage{subcaption} % Pacote incompatível com subfigure
%\usepackage{caption}
% ----------

\usepackage{multicol} % Pacote que permite usar multicolunas
\usepackage{float} % Pacote para usar o "H" do figure
\usepackage[dvipsnames]{xcolor} % Pacote de cores para textos e possivelmente outros
\usepackage{amsmath} % Pacote para usar o cases
\usepackage[makeroom]{cancel} % Pacote para usar o cancelar (em e embaixo)
\usepackage{enumitem} % Pacote para enumerar itens
\usepackage{amssymb} % Pacote para inserir os símbolos dos conjuntos numéricos N, Z, Q, R, C
\usepackage[T1]{fontenc}
\usepackage{enumitem}
\usepackage{textcomp}
\usepackage{pdfpages} % Pacote para inserir página em PDF
\usepackage{courier} % Pacote para usar a fonte courier
%\usepackage{matlab-prettifier}
\usepackage{lscape} % Pacote para mudar sentido da página
\newtheorem{theorem}{Teorema} % Pacote para inserir Teorema
\newtheorem{corollary}{Corollary} % Pacote para inserir Corolário
\newtheorem{remark}{Remark}
\newtheorem{mytheorem}{Definition}
\usepackage{hyperref} % Pacote para inserir hyperref
\usepackage{multirow} % Pacote para usar os comandos \multirow
\usepackage{pdfpages} % Pacote para incluir PDF's
\usepackage{listingsutf8} % Pacote para inserir arquivos externos tipo .txt
\usepackage{etoolbox}
\AtBeginEnvironment{verbatim}{\footnotesize}

\usepackage{standalone} % Pacote para inserir sub-arquivos LaTeX
\usepackage{bookmark} % Pacote para NÃO tratar o Índice como Apêndice e colocar letra

% Nome da listagem em português
\renewcommand{\lstlistingname}{Listagem}

% REFERÊNCIAS: ----------
%\usepackage[numbers, compress]{natbib} % Pacote para usar as referências bibliográficas
%\usepackage{usebib}
% \bibliographystyle{references} % Estilo das referências
%\setcitestyle{authoryear,aysep={},open={(},close={)}}
%\newcommand{\myand}{e}

% ----------


% PARÂMETROS DA FUNÇÃO LISTINGS: ----------
\lstset
{
    basicstyle=\small\ttfamily,
    commentstyle=\color{Green},
    keywordstyle=\color{Cerulean},
    frame=single,
    language=Python,
    morekeywords={True, False},
    numbers=left,
    numbersep=10pt,
    numberstyle=\footnotesize\color{Gray},
    showstringspaces=false,
    stringstyle=\color{Mulberry},
    tabsize=3,
    breaklines=true, % Este comando é para quebrar uma linha comprida
    postbreak=\mbox{\textcolor{red}{$\hookrightarrow$}\space}, % Este comando é para indicar com uma seta vermelha que a linha continua embaixo 
    extendedchars=true,
    literate=
        {á}{{\'a}}1 {Á}{{\'A}}1 {à}{{\`a}}1 {À}{{\`A}}1 {ã}{{\~a}}1 {Ã}{{\~A}}1 {â}{{\^a}}1 {Â}{{\^A}}1
        {ç}{{\c{c}}}1 {Ç}{{\c{C}}}1 {é}{{\'e}}1 {É}{{\'E}}1 {ê}{{\^e}}1 {Ê}{{\^E}}1 {í}{{\'\i}}1 {Í}{{\'I}}1
        {ó}{{\'o}}1 {Ó}{{\'O}}1 {ô}{{\^o}}1 {Ô}{{\^O}}1 {õ}{{\~o}}1 {Õ}{{\~O}}1 {ú}{{\'u}}1 {Ú}{{\'U}}1
        {ü}{{\"u}}1 {Ü}{{\"U}}1 {“}{{``}}1 {”}{{''}}1 {’}{{'}}1 {—}{{--}}1 {–}{{--}}1 {−}{{-}}1
        {±}{{$\pm$}}1 {σ}{{$\sigma$}}1 {→}{{$\to$}}1 {…}{{\ldots}}1
}

\lstset{inputencoding=utf8/latin1}

\lstdefinelanguage{Matlab}{
	language=Matlab,
	basicstyle=\ttfamily\small,
	keywordstyle=\color{blue}\bfseries,
	stringstyle=\color{red},
	commentstyle=\color{gray}\itshape,
	morekeywords={function, end, if, else, for, while, switch, case, otherwise},
	sensitive=true
}
% ----------

% ----------
\usepackage{nomencl} % Pacote para criar nomenclatura e lista de nomenclatura
\makenomenclature % Criar a nomenclatura
\renewcommand{\nomname}{Lista de Nomenclaturas} % Renomear o título da nomenclatura
% ----------

\usepackage{amsmath} % Pacote para inserir matrix com chaves, entre outras coisas

\usepackage{booktabs} % Pacote para usar \toprule, \midrule e \bottomrule


% EXECUTAR CÁLCULOS A PARTIR DA LINHA DE COMANDO: ---------
%\usepackage{siunitx, xfp} % Pacote para executar cálculos a partir da linha de comando
%\usepackage{l3regex}%unless you have very new expl3
%%\sisetup{locale=DE}
%%\let\oldnum\num
%%\ExplSyntaxOn
%\def\num#1{\def\tmp{#1}\regex_replace_all:nnN{,([0-9])}{.\1}{\tmp}\oldnum{\tmp}}
%\ExplSyntaxOff
% ----------


%\pagecolor{black} % Alterar a cor da página
%\color{white} % Alterar a cor do texto

% ----------
% Pacotes para inserir algoritmos e pseudo-códigos:
\setlength{\marginparwidth}{2.5cm}
\usepackage[colorinlistoftodos]{todonotes}
\usepackage{algorithm} % Pacote para inserir algoritmos
\usepackage{algpseudocode} % Pacote para inserir algoritmos e pseudo-códigos
\floatname{algorithm}{Algoritmo} % Renomear a palavra Algorithm para Procedimento
\renewcommand{\algorithmicrequire}{\textbf{entrada:}} % Renomear a palavra Input para Entrada
\renewcommand{\algorithmicensure}{\textbf{saída:}} % Renomear a palavra Output para Saída
\renewcommand{\algorithmicwhile}{\textbf{enquanto}} % Renomear a palavra While para Enquanto
\renewcommand{\algorithmicif}{\textbf{se}} % Renomear a palavra If para Se
\renewcommand{\algorithmicend}{\textbf{fim}} % Renomear a palavra End Output para Fim do
\renewcommand{\algorithmicprocedure}{\textbf{procedimento}} % Renomear a palavra Procedure para Procedimento
\renewcommand{\algorithmicfor}{\textbf{para}} % Renomear a palavra For para Para
\renewcommand{\algorithmicelse}{\textbf{senão}} % Renomear a palavra Else para Senão
\renewcommand{\algorithmicfunction}{\textbf{função}} % Renomear a palavra Function para Função 
\renewcommand{\algorithmicthen}{\textbf{então}} % Renomear a palavra Then para Então
\renewcommand{\algorithmicdo}{\textbf{fazer}} % Renomear a palavra Do para Fazer
% ----------

\newcommand{\quotes}[1]{``#1''} % Comando para inserir um texto dentro de aspas-dupla com espaço após o final

% FLUXOGRAMA: ----------
\usepackage{tikz}
\usetikzlibrary{positioning, shapes.geometric}
\tikzstyle{line}=[draw] % here

\tikzstyle{startstop} = [rectangle, rounded corners, minimum width=3cm,
minimum height=1cm,text centered, draw=black, fill=white!30]
\tikzstyle{arrow} = [thick, ->, >=stealth]
\tikzstyle{process} = [rectangle, minimum width=3cm, minimum height=1cm, text centered, draw=black, fill=white!30]
\tikzstyle{io} = [trapezium, trapezium left angle=70, trapezium right angle=110, minimum width=3cm, minimum height=1cm, text centered, draw=black, fill=white!30]
\tikzstyle{connector} = [draw, circle, fill=white!20]
\tikzstyle{decision} = [diamond, minimum width=3cm, minimum height=1cm, text centered, draw=black, fill=white!30, aspect=3]
% ----------


% CAIXA ENTORNO DAS FÓRMULAS: ----------
\usepackage{dashbox}
\newcommand\dboxed[1]{\dbox{\ensuremath{#1}}}

% ----------

%\usepackage{tabularx} % para colunas com largura ajustável

% -------- CONFIGURAÇÃO DE MARGENS --------
% \usepackage[top=2cm,bottom=2cm,left=3cm,right=2cm]{geometry}

% -------- PACOTE fancyhdr --------
\usepackage{fancyhdr}

\setlength{\marginparwidth}{2cm}

% Ajuste do tamanho do cabeçalho para acompanhar exatamente \textwidth
\setlength{\headwidth}{\textwidth}

% Define um headheight mínimo (para evitar avisos)
\setlength{\headheight}{15pt} 

% Limpa configurações padrão de cabeçalho/rodapé
\pagestyle{fancy}
\fancyhf{} 

% Marcações automáticas de seção para o cabeçalho (opcional)
\renewcommand{\sectionmark}[1]{%
	\markboth{\thesection\ #1}{}%
}

% Texto do cabeçalho à esquerda: seção e título
% Usamos \leftmark pois o \sectionmark acima grava no \leftmark
\fancyhead[L]{\small \leftmark}

% Texto do cabeçalho à direita: número de página
\fancyhead[R]{\small \thepage}

% Linha horizontal no cabeçalho (espessura 0.4pt)
\renewcommand{\headrulewidth}{0.4pt}

% Sobrescreve o título padrão “References” para “Referências”
%\renewcommand{\refname}{Referências}
% Tradução de autoref para o português
\def\figureautorefname{Figura}
\def\tableautorefname{Tabela}
\def\equationautorefname{Equa\c{c}\~ao}

% \hypersetup{backref=true}

\ifdefined\backref
        \renewcommand*{\backref}[1]{}
        \renewcommand*{\backrefentrycount}[2]{#1\ifnum#2>1 ~(#2)\fi}
\fi

\DeclarePairedDelimiter{\ceil}{\lceil}{\rceil}

\onehalfspacing

\usepackage{pgfgantt}

%++++++++++++++++++++++++++++++++++++++++++++++++++++++++++++++++++++++++++++++
% Identificacoes (se o trabalho for em inglês, insira os dados em inglês)
% Para entradas abreviadas de Professora (Profa.) em português escreva: Prof$^\textnormal{a}$.
%++++++++++++++++++++++++++++++++++++++++++++++++++++++++++++++++++++++++++++++
\course{Engenharia Aeronáutica e Mecânica} % Programa de PG ou Curso de Graduação
\area{Propulsão Aeroespacial e Energia} % Área de concentração na PG (Não utilizado no caso de TG)

% Autor do trabalho: Nome Sobrenome
\authorgender{masc} %sexo: masc ou fem
\author{\textbf{EDEN DENIS} Ferreira da Silva Lopes}{Santos}
\itaauthoraddress{Rua João Fonseca dos Santos, 20, Apto. 93}{12.230-088}{São José dos Campos--SP}

% Titulo da Tese/Dissertação
\title{Aplicação de Inteligência Artificial (AI) em Banco de Dados de Emissão de Motores}

% Orientador
\advisorgender{masc} % masc ou fem
\advisor{Prof.~Dr.}{Pedro Teixeira \textbf{LACAVA}}{ITA}

% Coorientador (Caso não haja coorientador, colocar ambas as variáveis \coadvisorgender e \coadvisor comentadas, com um % na frente)
\coadvisorgender{masc}
\coadvisor{}{}  % {Prof.~Dr.}{Carlos Alberto Rocha \textbf{PIMENTEL}}{ITA}

% Pró-reitor da Pós-graduação
\bossgender{fem}
\boss{Prof.~Dra.}{Emília \textbf{VILLANI}}

%Coordenador do curso no caso de TG
\bosscoursegender{masc}
\bosscourse{Prof.~Dr.}{ - } % masc ou fem

% Palavras-Chaves informadas pela Biblioteca -> utilizada na CIP
\kwcip{Engine taxonomy}
\kwcip{Emission classification}
\kwcip{Inteligência Artificial}
\kwcip{Artificial Intelligence}
\kwcip{Python}

% membros da banca examinadora
\examiner{Prof. Dr.}{Pedro Teixeira \textbf{LACAVA}}{President}{ITA}
\examiner{Prof. Dr.}{}{Internal Member}{ITA} % Vinícius \textbf{MALATESTA}}{Internal Member}{ITA}
\examiner{Prof. Dr.}{}{Internal Member}{ITA} % {Elisan dos Santos \textbf{MAGALHÃES}}{Internal Member}{ITA}
\examiner{Prof. Dr.}{}{Internal Member}{ITA} % {Carlos Alberto Rocha \textbf{PIMENTEL}}{External Member}{UFABC}
\examiner{Prof. Dr.}{}{Internal Member}{ITA} % {\textbf{KARL} Peter Burr}{External Member}{UFABC}

% Data da defesa (mês em maiúsculo, se trabalho em inglês, e minúsculo se trabalho em português) 
\date{14}{JULHO}{2025}

% Número CDU - (somente para TG)
\cdu{621.38}

% Folha de Registro do Documento
\newcommand{\advisorTitle}{Prof.~Dr.}
\newcommand{\advisorName}{Pedro Teixeira \textbf{LACAVA}}

\newcommand{\coadvisorTitle}{Prof.~Dr.}
\newcommand{\coadvisorName}{Carlos Alberto Rocha \textbf{PIMENTEL}}

\newcommand{\dateDay}{14}
\newcommand{\dateMonth}{JULHO}
\newcommand{\dateYear}{2025}

% Data de publicação e defesa (usada na folha de registro)
\newcommand{\frdDataPublicacao}{04 de agosto de 2025}
\newcommand{\frdNumeroRegistro}{DCTA/ITA/TD-029/2020}
\newcommand{\frdPalavrasAutor}{Big Data; Engine Emission; Emission Databank; Clustering; Artificial Intelligence}
\newcommand{\frdPalavrasResultado}{Taxonomia; Emissão de Motores; Planejamento de Dados; Inteligência Artificial; Processamento de Dados}
\newcommand{\frdApresentacao}{
	ITA, São José dos Campos. Curso de Mestrado. Programa de Pós-Graduação em Engenharia Aeronáutica e Mecânica. Área de Propulsão Aeroespacial e Energia.\\
	Orientador: \advisorTitle~\advisorName.\\
	Coorientador: \coadvisorTitle~\coadvisorName.\\
	Defesa em \dateDay/\dateMonth/\dateYear. Publicada em \frdDataPublicacao.
}

\newcommand{\PorcentagemDaLargura}{0.80}

\begin{document}

	\newpage
	\chapter{Conclusão}
		
		\section{Conclusão}
		
			\textcolor{red}{Ao desenvolver o estudo e as atividades, foi possível observar que, o projeto caminhou na direção correta, baseada em soluções históricas, a Matemática escolhida é solução para o problema de particionar as aeronaves. A aplicação do método apresentado foi desenvolvida genericamente a qual serve de ferramenta Matemática-computacional para classificar outros bancos de dados.}
			
%			Um dos grandes problemas da indústria aeroespacial, em particular a de construção de aeronaves civis e militares, é o gasto no ante-projeto desses veículos. É comum o aporte de gastos vultosos e de um grande tempo dos engenheiros projetistas para descobrir que o projeto focado é simplesmente inadequado a missão ou ao mercado. Em alguns casos, tudo isso é descoberto depois de toda a verba e tempo despendidos, o que pode ocasionar a falência da companhia. Portanto, o $ k $-means, técnica de Big Data aerotaxonométrica permite, como ficou confirmado neste trabalho, uma melhor orientação inicial para um possível projeto bem sucedido, em termos do seu início, meio e fim; evolução ajustada segundo as preferências vigentes do mercado aeronaútico. Outro sim, é quanto a evolução da Engenharia atual na direção do futuro apontado pelo Big Data que, além de prescrever a minimização do tempo de trabalho aponta também para uma redução em escala do esforço econômico-financeiro das empresas de um modo em geral.
%			
%			O $ k $-means desdobrado na dendrogramação, permite através dos seus dendros, a alocação de um futuro projeto, por semelhanças de dados ser associado pictograficamente com todos os outros projetos, já realizados, com sucesso ou fracasso seja do ponto de vista comercial, do funcional ou de ambos. Essa classificação aerotaxonômica, possibilita ao engenheiro aeroespacial, diante mão e instantaneamente, ter uma direção fundamentada para a sua futura aeronave cuja possível modificação, nessa ou naquela característica de projeto, aloca-o nesse ou naquele dendro, nesse ou no conjunto de aeronaves projetados anteriormente. Assim, durante o desenvolvimento do projeto, pode-se verificar em qualquer etapa, se o mesmo encontra-se dentro de uma classificação pré-estipulada e, caso esse tenda à um outro grupo, indesejado, pode-se fazer as modificações técnicas necessárias para que o projeto não tome outro rumo e sua verba, na maioria das vezes em \textit{déficit}, seja consumida equivocadamente. Outro ponto que toda essa Matemática destaca, são possíveis semelhanças entre aeronaves jamais comparadas devido ao grande número de projetos e de variáveis, entre outros detalhes. Não apenas isso, mas também traz a possibilidade de projetar um produto de maior custo-benefício, portanto, cada vez mais competitivo.
%			
%			Constatou-se também que o $ k $-means, numericamente funciona como uma espécie de sistema planetário onde, os medioides são as estrelas e os elementos, planetas que são gravitacionados por uma força de atração produzida por esses corpos de maior massa, medioides. Consequentemente, o modelo, ao variar a posição dos medioides (sóis) a cada iteração, consegue convergir para o ponto ótimo (de mínima ação) e, a partir daí, a posição do medioide não se altera, pois encontrou-se o ponto de equilíbrio do sistema, tudo isso sem a necessidade de se executar \quotes{infinitas} iterações já que o método converge rapidamente. Portanto, o método iterativo matemático-computacional adotado elimina os equívocos naturais das iterações anteriores, uma vez que, o critério de convergência adotado e, sobretudo, a Matemática desenvolvida, faz com que o ponto de equilíbrio seja atingido. Isto é, um elemento alocado equivocadamente em um medioide que não seja o seu ideal, será remanejado a cada passo em direção ao seu medioide ideal, até estabilizar neste.   
%			
%			No que diz respeito ao domínio comercial e, sobretudo militar, essa Matemática, atual e baseada em grandes quantidades de dados (Big Data), é ferramente fundamental na análise, comparação e desenvolvimento de projetos. Não é à toa que as empresas deste ramo da Engenharia, buscam cada vez mais profissionais com este tipo de conhecimento, que conhecem soluções não apenas específicas, mas que também permitem enxergar o mercado de maneira ampla, acima de tudo, circunstânciais. 
%			
%			Ao analisar os resultados (ver seção \ref{sec:ResultadosEDiscussoes}) chega-se a certas conclusões que somente alguém muito experiente saberia. No entanto, o $ k $-means, traz a possibilidade de uma pessoa que não seja perita neste tipo de trabalho, possa extrair conhecimento através da informação (banco de dados), mas também para quem for provecto, que este perceba certas nuances que apenas uma análise mais profunda é capaz de destacar muito além do que  já tenha feito. Vale lembrar que as discussões realizadas neste trabalho, não se limitam aqui, existe mais a ser discutido sobre tudo isso, graças aos novos métodos matemático que vem surgindo nas últimas décadas os quais trazem inúmeras possibilidades. Outro detalhe é que informações conhecidas podem ser confirmadas como é citado e comparado ao longo do texto, vale lembrar o caso do A380 e do F/A-22 Raptor discutidos na seção (\ref{sec:ResultadosEDiscussoes}).
%			
%			A Matemática é adequada e se mostrou, apesar de falibalista, com um alto grau de acerto, um excelente método para particionar bancos de dados; de fato, os resultados obtidos além de outras análises que possam ser feitas, servem de base para novos projetos e, não só isso, mas todo o estudo pode ser utilizado em outras áreas do conhecimento ou também pode ser feita diversas análises dentro do mesmo trabalho, neste caso, como se fosse uma espécie de árvores de decisão. Ao executar o particionamento do banco de dados, escolhe-se, a partir dos dendros, um grupo no qual é feita uma nova classificação que resultará em outro dendrograma e, assim, sucessivamente, pode-se extrair resultados muito mais profundos e específicos quanto às aeronaves. Da mesma forma, ao escolher um certo grupo como base de um novo projeto, esta mesma análise, pode sugerir a separação dos subsistemas  das aeronaves (superfícies aerodinâmicas, sistemas de propulsão, controle, aviônica, comunicação e componentes estruturais, dentre outros) do grupo escolhido, logo, podem ser escolhidos os subsistemas do novo projeto com o que há de melhor em termos de custo-benefício das aeronaves daquele grupo adotado inicialmente.
%			
%			O trabalho desenvolvido permite acompanhar as tendências do mercado tanto civil, quanto militar, pois é baseado, como visto, no banco de dados do que já foi produzido, ou melhor, deve-se fazer com certa frequência a análise para acompanhar tais modulações e, isto conduzirá, de fato, ao que o mercado exige em termos financeiros e ao mesmo tempo no que diz respeito à sua missão. Exemplos dessas tendências são os voos comerciais de curta distâncias, mais frequentes no mercado atual, daí a fusão entre empresas de maior porte com as de menor e as atualizações que são feitas na aviônica, em específico (mas, não apenas isso) o sistema radar dos caças que ocorrem para não se defasarem em relação à outro caça de mesma grandeza. Isso possibilita a companhia, estar sempre renovada quanto às tendências e exigências do mercado, requisito fundamental para que as empresas se mantenham na competição empresarial bem como para que uma nação mantenha sua soberania aérea. 
%			
%			Enfim, percebeu-se no $ k $-means uma possibilidade de entrever, não apenas estas, mas também outras soluções extremamente abrangentes e eminentemente funcionais para a Engenharia, porém, com um outro estilo de Matemática que contrapõe às maneiras clássicas. Fundamentado no DNA (conjunto de variáveis) das aeronaves, ao distanciar-se cada vez mais da subjetividade, a Matemática traçou as verossimilhanças entre estas, ou seja, a não absoluta, no entanto verdadeira semelhança. Este método, intercambiável, é versátil, líquido o suficiente para acompanhar as perspectivas da sociedade, logo, traz consigo, uma probabilidade real de termos em mãos a Matemática da Engenharia do futuro; no caso do presente...    

\end{document}
%	\documentclass[12pt]{article} % Tamanho de letra 12, papel A4, Estilo de artigo

\usepackage[utf8]{inputenc} % Pacote para permitir o uso de acentos
% Ajustes de hifenização
\tolerance=1
\emergencystretch=\maxdimen
\hyphenpenalty=10000
\hbadness=10000
\usepackage{lipsum} % Pacote para gerar texto aleatório
\usepackage{ragged2e} % Pacote para mais opções de alinhamento
\usepackage{graphicx} % Pacote para inserir figuras jpeg, pdf e png
\usepackage{ae}
\usepackage{epsfig}
\usepackage{mathtools}
\usepackage[export]{adjustbox}
\usepackage{lineno}
\usepackage{amsthm}
\usepackage{commath}
\usepackage{color, colortbl}
\usepackage{bbm}
\usepackage{setspace} % Pacote para ajustar espaçamento entre linhas

% Instructions to prevent LaTeX from hyphenating words
\tolerance=1
\emergencystretch=\maxdimen
\hyphenpenalty=10000
\hbadness=10000

\graphicspath{{figures/}} % Indica o diretório no qual as imagens estão sendo salvas, evitando que precisamos digitar parte do caminho da pasta que contém as imagens
\usepackage[rightcaption]{sidecap} % Pacote para mais opções de figura, colocar, título da legenda da figura ao lado direito
\usepackage{url} % Pacote para citar referências de sites
\usepackage{makeidx} % Pacote para criar índice
\makeindex % Habilitar a criação do índice
\usepackage{longtable} % Pacote para usar tabelas longas

% ----------
% Pacotes para dar quebra de linha dentro da tabela:
%\usepackage{fourier} 
%\usepackage{array}
\usepackage{makecell}

\renewcommand\theadalign{bc}
\renewcommand\theadfont{\bfseries}
\renewcommand\theadgape{\Gape[4pt]}
\renewcommand\cellgape{\Gape[4pt]}
% ----------

\usepackage{wrapfig} % Pacote para alinhar figuras à esquerda ou à direita do texto
\usepackage{ae}
\usepackage{epsfig}
\usepackage{mathtools}
\usepackage[export]{adjustbox}
\usepackage{lineno}
\usepackage{amsfonts}
\usepackage{amsthm}
\usepackage{commath}
\usepackage{color, colortbl}
\usepackage{bbm}

% ----------
% Pacote para inserir cabeçalho e rodapé:
% Atenção! Se remover o cabeçalho e rodapé, o LaTeX coloca como cabeçalho o padrão dele, que é bem interessante também!
% \usepackage{fancyhdr} % Pacote para inserir cabeçalho e rodapé
% \pagestyle{fancy} % Estilo do cabeçalho e rodapé, neste caso apenas uma linha no cabeçalho
%\lhead{{\footnotesize Prandiano\textregistered}} % Inserir texto no canto esquerdo do cabeçalho
%\chead{centro} % Inserir texto no centro do cabeçalho
%\rhead{{\footnotesize Curso de Python}} % Inserir texto no canto direito do cabeçalho
%\lfoot{{\footnotesize Copyright \textcopyright 2018 Prandiano}} % Inserir texto no canto esquerdo do rodapé
%\cfoot{centro} % Inserir texto no centro do rodapé
%\rfoot{{\footnotesize Reprodução proibida: Lei 9.610 de 19/02/1998}} % Inserir texto no canto direito do rodapé
% O comando \renewcommand re-nova o parâmetro especificado
%\renewcommand{\footrulewidth}{0.4pt} % Inserir uma linha acima do rodapé
% ----------

% ----------
% Pacote para inserir múltiplas figuras:
\usepackage{subfigure} % Pacote para inserir múltiplas figuras
% \usepackage{subcaption} % Pacote incompatível com subfigure
%\usepackage{caption}
% ----------

\usepackage{multicol} % Pacote que permite usar multicolunas
\usepackage{float} % Pacote para usar o "H" do figure
\usepackage[dvipsnames]{xcolor} % Pacote de cores para textos e possivelmente outros
\usepackage{amsmath} % Pacote para usar o cases
\usepackage[makeroom]{cancel} % Pacote para usar o cancelar (em e embaixo)
\usepackage{enumitem} % Pacote para enumerar itens
\usepackage{amssymb} % Pacote para inserir os símbolos dos conjuntos numéricos N, Z, Q, R, C
\usepackage[T1]{fontenc}
\usepackage{enumitem}
\usepackage{textcomp}
\usepackage{pdfpages} % Pacote para inserir página em PDF
\usepackage{courier} % Pacote para usar a fonte courier
%\usepackage{matlab-prettifier}
\usepackage{lscape} % Pacote para mudar sentido da página
\newtheorem{theorem}{Teorema} % Pacote para inserir Teorema
\newtheorem{corollary}{Corollary} % Pacote para inserir Corolário
\newtheorem{remark}{Remark}
\newtheorem{mytheorem}{Definition}
\usepackage{hyperref} % Pacote para inserir hyperref
\usepackage{multirow} % Pacote para usar os comandos \multirow
\usepackage{pdfpages} % Pacote para incluir PDF's
\usepackage{listingsutf8} % Pacote para inserir arquivos externos tipo .txt
\usepackage{etoolbox}
\AtBeginEnvironment{verbatim}{\footnotesize}

\usepackage{standalone} % Pacote para inserir sub-arquivos LaTeX
\usepackage{bookmark} % Pacote para NÃO tratar o Índice como Apêndice e colocar letra

% Nome da listagem em português
\renewcommand{\lstlistingname}{Listagem}

% REFERÊNCIAS: ----------
%\usepackage[numbers, compress]{natbib} % Pacote para usar as referências bibliográficas
%\usepackage{usebib}
% \bibliographystyle{references} % Estilo das referências
%\setcitestyle{authoryear,aysep={},open={(},close={)}}
%\newcommand{\myand}{e}

% ----------


% PARÂMETROS DA FUNÇÃO LISTINGS: ----------
\lstset
{
    basicstyle=\small\ttfamily,
    commentstyle=\color{Green},
    keywordstyle=\color{Cerulean},
    frame=single,
    language=Python,
    morekeywords={True, False},
    numbers=left,
    numbersep=10pt,
    numberstyle=\footnotesize\color{Gray},
    showstringspaces=false,
    stringstyle=\color{Mulberry},
    tabsize=3,
    breaklines=true, % Este comando é para quebrar uma linha comprida
    postbreak=\mbox{\textcolor{red}{$\hookrightarrow$}\space}, % Este comando é para indicar com uma seta vermelha que a linha continua embaixo 
    extendedchars=true,
    literate=
        {á}{{\'a}}1 {Á}{{\'A}}1 {à}{{\`a}}1 {À}{{\`A}}1 {ã}{{\~a}}1 {Ã}{{\~A}}1 {â}{{\^a}}1 {Â}{{\^A}}1
        {ç}{{\c{c}}}1 {Ç}{{\c{C}}}1 {é}{{\'e}}1 {É}{{\'E}}1 {ê}{{\^e}}1 {Ê}{{\^E}}1 {í}{{\'\i}}1 {Í}{{\'I}}1
        {ó}{{\'o}}1 {Ó}{{\'O}}1 {ô}{{\^o}}1 {Ô}{{\^O}}1 {õ}{{\~o}}1 {Õ}{{\~O}}1 {ú}{{\'u}}1 {Ú}{{\'U}}1
        {ü}{{\"u}}1 {Ü}{{\"U}}1 {“}{{``}}1 {”}{{''}}1 {’}{{'}}1 {—}{{--}}1 {–}{{--}}1 {−}{{-}}1
        {±}{{$\pm$}}1 {σ}{{$\sigma$}}1 {→}{{$\to$}}1 {…}{{\ldots}}1
}

\lstset{inputencoding=utf8/latin1}

\lstdefinelanguage{Matlab}{
	language=Matlab,
	basicstyle=\ttfamily\small,
	keywordstyle=\color{blue}\bfseries,
	stringstyle=\color{red},
	commentstyle=\color{gray}\itshape,
	morekeywords={function, end, if, else, for, while, switch, case, otherwise},
	sensitive=true
}
% ----------

% ----------
\usepackage{nomencl} % Pacote para criar nomenclatura e lista de nomenclatura
\makenomenclature % Criar a nomenclatura
\renewcommand{\nomname}{Lista de Nomenclaturas} % Renomear o título da nomenclatura
% ----------

\usepackage{amsmath} % Pacote para inserir matrix com chaves, entre outras coisas

\usepackage{booktabs} % Pacote para usar \toprule, \midrule e \bottomrule


% EXECUTAR CÁLCULOS A PARTIR DA LINHA DE COMANDO: ---------
%\usepackage{siunitx, xfp} % Pacote para executar cálculos a partir da linha de comando
%\usepackage{l3regex}%unless you have very new expl3
%%\sisetup{locale=DE}
%%\let\oldnum\num
%%\ExplSyntaxOn
%\def\num#1{\def\tmp{#1}\regex_replace_all:nnN{,([0-9])}{.\1}{\tmp}\oldnum{\tmp}}
%\ExplSyntaxOff
% ----------


%\pagecolor{black} % Alterar a cor da página
%\color{white} % Alterar a cor do texto

% ----------
% Pacotes para inserir algoritmos e pseudo-códigos:
\setlength{\marginparwidth}{2.5cm}
\usepackage[colorinlistoftodos]{todonotes}
\usepackage{algorithm} % Pacote para inserir algoritmos
\usepackage{algpseudocode} % Pacote para inserir algoritmos e pseudo-códigos
\floatname{algorithm}{Algoritmo} % Renomear a palavra Algorithm para Procedimento
\renewcommand{\algorithmicrequire}{\textbf{entrada:}} % Renomear a palavra Input para Entrada
\renewcommand{\algorithmicensure}{\textbf{saída:}} % Renomear a palavra Output para Saída
\renewcommand{\algorithmicwhile}{\textbf{enquanto}} % Renomear a palavra While para Enquanto
\renewcommand{\algorithmicif}{\textbf{se}} % Renomear a palavra If para Se
\renewcommand{\algorithmicend}{\textbf{fim}} % Renomear a palavra End Output para Fim do
\renewcommand{\algorithmicprocedure}{\textbf{procedimento}} % Renomear a palavra Procedure para Procedimento
\renewcommand{\algorithmicfor}{\textbf{para}} % Renomear a palavra For para Para
\renewcommand{\algorithmicelse}{\textbf{senão}} % Renomear a palavra Else para Senão
\renewcommand{\algorithmicfunction}{\textbf{função}} % Renomear a palavra Function para Função 
\renewcommand{\algorithmicthen}{\textbf{então}} % Renomear a palavra Then para Então
\renewcommand{\algorithmicdo}{\textbf{fazer}} % Renomear a palavra Do para Fazer
% ----------

\newcommand{\quotes}[1]{``#1''} % Comando para inserir um texto dentro de aspas-dupla com espaço após o final

% FLUXOGRAMA: ----------
\usepackage{tikz}
\usetikzlibrary{positioning, shapes.geometric}
\tikzstyle{line}=[draw] % here

\tikzstyle{startstop} = [rectangle, rounded corners, minimum width=3cm,
minimum height=1cm,text centered, draw=black, fill=white!30]
\tikzstyle{arrow} = [thick, ->, >=stealth]
\tikzstyle{process} = [rectangle, minimum width=3cm, minimum height=1cm, text centered, draw=black, fill=white!30]
\tikzstyle{io} = [trapezium, trapezium left angle=70, trapezium right angle=110, minimum width=3cm, minimum height=1cm, text centered, draw=black, fill=white!30]
\tikzstyle{connector} = [draw, circle, fill=white!20]
\tikzstyle{decision} = [diamond, minimum width=3cm, minimum height=1cm, text centered, draw=black, fill=white!30, aspect=3]
% ----------


% CAIXA ENTORNO DAS FÓRMULAS: ----------
\usepackage{dashbox}
\newcommand\dboxed[1]{\dbox{\ensuremath{#1}}}

% ----------

%\usepackage{tabularx} % para colunas com largura ajustável

% -------- CONFIGURAÇÃO DE MARGENS --------
% \usepackage[top=2cm,bottom=2cm,left=3cm,right=2cm]{geometry}

% -------- PACOTE fancyhdr --------
\usepackage{fancyhdr}

\setlength{\marginparwidth}{2cm}

% Ajuste do tamanho do cabeçalho para acompanhar exatamente \textwidth
\setlength{\headwidth}{\textwidth}

% Define um headheight mínimo (para evitar avisos)
\setlength{\headheight}{15pt} 

% Limpa configurações padrão de cabeçalho/rodapé
\pagestyle{fancy}
\fancyhf{} 

% Marcações automáticas de seção para o cabeçalho (opcional)
\renewcommand{\sectionmark}[1]{%
	\markboth{\thesection\ #1}{}%
}

% Texto do cabeçalho à esquerda: seção e título
% Usamos \leftmark pois o \sectionmark acima grava no \leftmark
\fancyhead[L]{\small \leftmark}

% Texto do cabeçalho à direita: número de página
\fancyhead[R]{\small \thepage}

% Linha horizontal no cabeçalho (espessura 0.4pt)
\renewcommand{\headrulewidth}{0.4pt}

% Sobrescreve o título padrão “References” para “Referências”
%\renewcommand{\refname}{Referências}
% Tradução de autoref para o português
\def\figureautorefname{Figura}
\def\tableautorefname{Tabela}
\def\equationautorefname{Equa\c{c}\~ao}

% \hypersetup{backref=true}

\ifdefined\backref
        \renewcommand*{\backref}[1]{}
        \renewcommand*{\backrefentrycount}[2]{#1\ifnum#2>1 ~(#2)\fi}
\fi

\DeclarePairedDelimiter{\ceil}{\lceil}{\rceil}

\onehalfspacing

\usepackage{pgfgantt}

%++++++++++++++++++++++++++++++++++++++++++++++++++++++++++++++++++++++++++++++
% Identificacoes (se o trabalho for em inglês, insira os dados em inglês)
% Para entradas abreviadas de Professora (Profa.) em português escreva: Prof$^\textnormal{a}$.
%++++++++++++++++++++++++++++++++++++++++++++++++++++++++++++++++++++++++++++++
\course{Engenharia Aeronáutica e Mecânica} % Programa de PG ou Curso de Graduação
\area{Propulsão Aeroespacial e Energia} % Área de concentração na PG (Não utilizado no caso de TG)

% Autor do trabalho: Nome Sobrenome
\authorgender{masc} %sexo: masc ou fem
\author{\textbf{EDEN DENIS} Ferreira da Silva Lopes}{Santos}
\itaauthoraddress{Rua João Fonseca dos Santos, 20, Apto. 93}{12.230-088}{São José dos Campos--SP}

% Titulo da Tese/Dissertação
\title{Aplicação de Inteligência Artificial (AI) em Banco de Dados de Emissão de Motores}

% Orientador
\advisorgender{masc} % masc ou fem
\advisor{Prof.~Dr.}{Pedro Teixeira \textbf{LACAVA}}{ITA}

% Coorientador (Caso não haja coorientador, colocar ambas as variáveis \coadvisorgender e \coadvisor comentadas, com um % na frente)
\coadvisorgender{masc}
\coadvisor{}{}  % {Prof.~Dr.}{Carlos Alberto Rocha \textbf{PIMENTEL}}{ITA}

% Pró-reitor da Pós-graduação
\bossgender{fem}
\boss{Prof.~Dra.}{Emília \textbf{VILLANI}}

%Coordenador do curso no caso de TG
\bosscoursegender{masc}
\bosscourse{Prof.~Dr.}{ - } % masc ou fem

% Palavras-Chaves informadas pela Biblioteca -> utilizada na CIP
\kwcip{Engine taxonomy}
\kwcip{Emission classification}
\kwcip{Inteligência Artificial}
\kwcip{Artificial Intelligence}
\kwcip{Python}

% membros da banca examinadora
\examiner{Prof. Dr.}{Pedro Teixeira \textbf{LACAVA}}{President}{ITA}
\examiner{Prof. Dr.}{}{Internal Member}{ITA} % Vinícius \textbf{MALATESTA}}{Internal Member}{ITA}
\examiner{Prof. Dr.}{}{Internal Member}{ITA} % {Elisan dos Santos \textbf{MAGALHÃES}}{Internal Member}{ITA}
\examiner{Prof. Dr.}{}{Internal Member}{ITA} % {Carlos Alberto Rocha \textbf{PIMENTEL}}{External Member}{UFABC}
\examiner{Prof. Dr.}{}{Internal Member}{ITA} % {\textbf{KARL} Peter Burr}{External Member}{UFABC}

% Data da defesa (mês em maiúsculo, se trabalho em inglês, e minúsculo se trabalho em português) 
\date{14}{JULHO}{2025}

% Número CDU - (somente para TG)
\cdu{621.38}

% Folha de Registro do Documento
\newcommand{\advisorTitle}{Prof.~Dr.}
\newcommand{\advisorName}{Pedro Teixeira \textbf{LACAVA}}

\newcommand{\coadvisorTitle}{Prof.~Dr.}
\newcommand{\coadvisorName}{Carlos Alberto Rocha \textbf{PIMENTEL}}

\newcommand{\dateDay}{14}
\newcommand{\dateMonth}{JULHO}
\newcommand{\dateYear}{2025}

% Data de publicação e defesa (usada na folha de registro)
\newcommand{\frdDataPublicacao}{04 de agosto de 2025}
\newcommand{\frdNumeroRegistro}{DCTA/ITA/TD-029/2020}
\newcommand{\frdPalavrasAutor}{Big Data; Engine Emission; Emission Databank; Clustering; Artificial Intelligence}
\newcommand{\frdPalavrasResultado}{Taxonomia; Emissão de Motores; Planejamento de Dados; Inteligência Artificial; Processamento de Dados}
\newcommand{\frdApresentacao}{
	ITA, São José dos Campos. Curso de Mestrado. Programa de Pós-Graduação em Engenharia Aeronáutica e Mecânica. Área de Propulsão Aeroespacial e Energia.\\
	Orientador: \advisorTitle~\advisorName.\\
	Coorientador: \coadvisorTitle~\coadvisorName.\\
	Defesa em \dateDay/\dateMonth/\dateYear. Publicada em \frdDataPublicacao.
}

\newcommand{\PorcentagemDaLargura}{0.80}

\begin{document}
	
	\newpage
	\chapter{Trabalhos futuros}
	

\end{document}
%	\documentclass[12pt]{article} % Tamanho de letra 12, papel A4, Estilo de artigo

\usepackage[utf8]{inputenc} % Pacote para permitir o uso de acentos
% Ajustes de hifenização
\tolerance=1
\emergencystretch=\maxdimen
\hyphenpenalty=10000
\hbadness=10000
\usepackage{lipsum} % Pacote para gerar texto aleatório
\usepackage{ragged2e} % Pacote para mais opções de alinhamento
\usepackage{graphicx} % Pacote para inserir figuras jpeg, pdf e png
\usepackage{ae}
\usepackage{epsfig}
\usepackage{mathtools}
\usepackage[export]{adjustbox}
\usepackage{lineno}
\usepackage{amsthm}
\usepackage{commath}
\usepackage{color, colortbl}
\usepackage{bbm}
\usepackage{setspace} % Pacote para ajustar espaçamento entre linhas

% Instructions to prevent LaTeX from hyphenating words
\tolerance=1
\emergencystretch=\maxdimen
\hyphenpenalty=10000
\hbadness=10000

\graphicspath{{figures/}} % Indica o diretório no qual as imagens estão sendo salvas, evitando que precisamos digitar parte do caminho da pasta que contém as imagens
\usepackage[rightcaption]{sidecap} % Pacote para mais opções de figura, colocar, título da legenda da figura ao lado direito
\usepackage{url} % Pacote para citar referências de sites
\usepackage{makeidx} % Pacote para criar índice
\makeindex % Habilitar a criação do índice
\usepackage{longtable} % Pacote para usar tabelas longas

% ----------
% Pacotes para dar quebra de linha dentro da tabela:
%\usepackage{fourier} 
%\usepackage{array}
\usepackage{makecell}

\renewcommand\theadalign{bc}
\renewcommand\theadfont{\bfseries}
\renewcommand\theadgape{\Gape[4pt]}
\renewcommand\cellgape{\Gape[4pt]}
% ----------

\usepackage{wrapfig} % Pacote para alinhar figuras à esquerda ou à direita do texto
\usepackage{ae}
\usepackage{epsfig}
\usepackage{mathtools}
\usepackage[export]{adjustbox}
\usepackage{lineno}
\usepackage{amsfonts}
\usepackage{amsthm}
\usepackage{commath}
\usepackage{color, colortbl}
\usepackage{bbm}

% ----------
% Pacote para inserir cabeçalho e rodapé:
% Atenção! Se remover o cabeçalho e rodapé, o LaTeX coloca como cabeçalho o padrão dele, que é bem interessante também!
% \usepackage{fancyhdr} % Pacote para inserir cabeçalho e rodapé
% \pagestyle{fancy} % Estilo do cabeçalho e rodapé, neste caso apenas uma linha no cabeçalho
%\lhead{{\footnotesize Prandiano\textregistered}} % Inserir texto no canto esquerdo do cabeçalho
%\chead{centro} % Inserir texto no centro do cabeçalho
%\rhead{{\footnotesize Curso de Python}} % Inserir texto no canto direito do cabeçalho
%\lfoot{{\footnotesize Copyright \textcopyright 2018 Prandiano}} % Inserir texto no canto esquerdo do rodapé
%\cfoot{centro} % Inserir texto no centro do rodapé
%\rfoot{{\footnotesize Reprodução proibida: Lei 9.610 de 19/02/1998}} % Inserir texto no canto direito do rodapé
% O comando \renewcommand re-nova o parâmetro especificado
%\renewcommand{\footrulewidth}{0.4pt} % Inserir uma linha acima do rodapé
% ----------

% ----------
% Pacote para inserir múltiplas figuras:
\usepackage{subfigure} % Pacote para inserir múltiplas figuras
% \usepackage{subcaption} % Pacote incompatível com subfigure
%\usepackage{caption}
% ----------

\usepackage{multicol} % Pacote que permite usar multicolunas
\usepackage{float} % Pacote para usar o "H" do figure
\usepackage[dvipsnames]{xcolor} % Pacote de cores para textos e possivelmente outros
\usepackage{amsmath} % Pacote para usar o cases
\usepackage[makeroom]{cancel} % Pacote para usar o cancelar (em e embaixo)
\usepackage{enumitem} % Pacote para enumerar itens
\usepackage{amssymb} % Pacote para inserir os símbolos dos conjuntos numéricos N, Z, Q, R, C
\usepackage[T1]{fontenc}
\usepackage{enumitem}
\usepackage{textcomp}
\usepackage{pdfpages} % Pacote para inserir página em PDF
\usepackage{courier} % Pacote para usar a fonte courier
%\usepackage{matlab-prettifier}
\usepackage{lscape} % Pacote para mudar sentido da página
\newtheorem{theorem}{Teorema} % Pacote para inserir Teorema
\newtheorem{corollary}{Corollary} % Pacote para inserir Corolário
\newtheorem{remark}{Remark}
\newtheorem{mytheorem}{Definition}
\usepackage{hyperref} % Pacote para inserir hyperref
\usepackage{multirow} % Pacote para usar os comandos \multirow
\usepackage{pdfpages} % Pacote para incluir PDF's
\usepackage{listingsutf8} % Pacote para inserir arquivos externos tipo .txt
\usepackage{etoolbox}
\AtBeginEnvironment{verbatim}{\footnotesize}

\usepackage{standalone} % Pacote para inserir sub-arquivos LaTeX
\usepackage{bookmark} % Pacote para NÃO tratar o Índice como Apêndice e colocar letra

% Nome da listagem em português
\renewcommand{\lstlistingname}{Listagem}

% REFERÊNCIAS: ----------
%\usepackage[numbers, compress]{natbib} % Pacote para usar as referências bibliográficas
%\usepackage{usebib}
% \bibliographystyle{references} % Estilo das referências
%\setcitestyle{authoryear,aysep={},open={(},close={)}}
%\newcommand{\myand}{e}

% ----------


% PARÂMETROS DA FUNÇÃO LISTINGS: ----------
\lstset
{
    basicstyle=\small\ttfamily,
    commentstyle=\color{Green},
    keywordstyle=\color{Cerulean},
    frame=single,
    language=Python,
    morekeywords={True, False},
    numbers=left,
    numbersep=10pt,
    numberstyle=\footnotesize\color{Gray},
    showstringspaces=false,
    stringstyle=\color{Mulberry},
    tabsize=3,
    breaklines=true, % Este comando é para quebrar uma linha comprida
    postbreak=\mbox{\textcolor{red}{$\hookrightarrow$}\space}, % Este comando é para indicar com uma seta vermelha que a linha continua embaixo 
    extendedchars=true,
    literate=
        {á}{{\'a}}1 {Á}{{\'A}}1 {à}{{\`a}}1 {À}{{\`A}}1 {ã}{{\~a}}1 {Ã}{{\~A}}1 {â}{{\^a}}1 {Â}{{\^A}}1
        {ç}{{\c{c}}}1 {Ç}{{\c{C}}}1 {é}{{\'e}}1 {É}{{\'E}}1 {ê}{{\^e}}1 {Ê}{{\^E}}1 {í}{{\'\i}}1 {Í}{{\'I}}1
        {ó}{{\'o}}1 {Ó}{{\'O}}1 {ô}{{\^o}}1 {Ô}{{\^O}}1 {õ}{{\~o}}1 {Õ}{{\~O}}1 {ú}{{\'u}}1 {Ú}{{\'U}}1
        {ü}{{\"u}}1 {Ü}{{\"U}}1 {“}{{``}}1 {”}{{''}}1 {’}{{'}}1 {—}{{--}}1 {–}{{--}}1 {−}{{-}}1
        {±}{{$\pm$}}1 {σ}{{$\sigma$}}1 {→}{{$\to$}}1 {…}{{\ldots}}1
}

\lstset{inputencoding=utf8/latin1}

\lstdefinelanguage{Matlab}{
	language=Matlab,
	basicstyle=\ttfamily\small,
	keywordstyle=\color{blue}\bfseries,
	stringstyle=\color{red},
	commentstyle=\color{gray}\itshape,
	morekeywords={function, end, if, else, for, while, switch, case, otherwise},
	sensitive=true
}
% ----------

% ----------
\usepackage{nomencl} % Pacote para criar nomenclatura e lista de nomenclatura
\makenomenclature % Criar a nomenclatura
\renewcommand{\nomname}{Lista de Nomenclaturas} % Renomear o título da nomenclatura
% ----------

\usepackage{amsmath} % Pacote para inserir matrix com chaves, entre outras coisas

\usepackage{booktabs} % Pacote para usar \toprule, \midrule e \bottomrule


% EXECUTAR CÁLCULOS A PARTIR DA LINHA DE COMANDO: ---------
%\usepackage{siunitx, xfp} % Pacote para executar cálculos a partir da linha de comando
%\usepackage{l3regex}%unless you have very new expl3
%%\sisetup{locale=DE}
%%\let\oldnum\num
%%\ExplSyntaxOn
%\def\num#1{\def\tmp{#1}\regex_replace_all:nnN{,([0-9])}{.\1}{\tmp}\oldnum{\tmp}}
%\ExplSyntaxOff
% ----------


%\pagecolor{black} % Alterar a cor da página
%\color{white} % Alterar a cor do texto

% ----------
% Pacotes para inserir algoritmos e pseudo-códigos:
\setlength{\marginparwidth}{2.5cm}
\usepackage[colorinlistoftodos]{todonotes}
\usepackage{algorithm} % Pacote para inserir algoritmos
\usepackage{algpseudocode} % Pacote para inserir algoritmos e pseudo-códigos
\floatname{algorithm}{Algoritmo} % Renomear a palavra Algorithm para Procedimento
\renewcommand{\algorithmicrequire}{\textbf{entrada:}} % Renomear a palavra Input para Entrada
\renewcommand{\algorithmicensure}{\textbf{saída:}} % Renomear a palavra Output para Saída
\renewcommand{\algorithmicwhile}{\textbf{enquanto}} % Renomear a palavra While para Enquanto
\renewcommand{\algorithmicif}{\textbf{se}} % Renomear a palavra If para Se
\renewcommand{\algorithmicend}{\textbf{fim}} % Renomear a palavra End Output para Fim do
\renewcommand{\algorithmicprocedure}{\textbf{procedimento}} % Renomear a palavra Procedure para Procedimento
\renewcommand{\algorithmicfor}{\textbf{para}} % Renomear a palavra For para Para
\renewcommand{\algorithmicelse}{\textbf{senão}} % Renomear a palavra Else para Senão
\renewcommand{\algorithmicfunction}{\textbf{função}} % Renomear a palavra Function para Função 
\renewcommand{\algorithmicthen}{\textbf{então}} % Renomear a palavra Then para Então
\renewcommand{\algorithmicdo}{\textbf{fazer}} % Renomear a palavra Do para Fazer
% ----------

\newcommand{\quotes}[1]{``#1''} % Comando para inserir um texto dentro de aspas-dupla com espaço após o final

% FLUXOGRAMA: ----------
\usepackage{tikz}
\usetikzlibrary{positioning, shapes.geometric}
\tikzstyle{line}=[draw] % here

\tikzstyle{startstop} = [rectangle, rounded corners, minimum width=3cm,
minimum height=1cm,text centered, draw=black, fill=white!30]
\tikzstyle{arrow} = [thick, ->, >=stealth]
\tikzstyle{process} = [rectangle, minimum width=3cm, minimum height=1cm, text centered, draw=black, fill=white!30]
\tikzstyle{io} = [trapezium, trapezium left angle=70, trapezium right angle=110, minimum width=3cm, minimum height=1cm, text centered, draw=black, fill=white!30]
\tikzstyle{connector} = [draw, circle, fill=white!20]
\tikzstyle{decision} = [diamond, minimum width=3cm, minimum height=1cm, text centered, draw=black, fill=white!30, aspect=3]
% ----------


% CAIXA ENTORNO DAS FÓRMULAS: ----------
\usepackage{dashbox}
\newcommand\dboxed[1]{\dbox{\ensuremath{#1}}}

% ----------

%\usepackage{tabularx} % para colunas com largura ajustável

% -------- CONFIGURAÇÃO DE MARGENS --------
% \usepackage[top=2cm,bottom=2cm,left=3cm,right=2cm]{geometry}

% -------- PACOTE fancyhdr --------
\usepackage{fancyhdr}

\setlength{\marginparwidth}{2cm}

% Ajuste do tamanho do cabeçalho para acompanhar exatamente \textwidth
\setlength{\headwidth}{\textwidth}

% Define um headheight mínimo (para evitar avisos)
\setlength{\headheight}{15pt} 

% Limpa configurações padrão de cabeçalho/rodapé
\pagestyle{fancy}
\fancyhf{} 

% Marcações automáticas de seção para o cabeçalho (opcional)
\renewcommand{\sectionmark}[1]{%
	\markboth{\thesection\ #1}{}%
}

% Texto do cabeçalho à esquerda: seção e título
% Usamos \leftmark pois o \sectionmark acima grava no \leftmark
\fancyhead[L]{\small \leftmark}

% Texto do cabeçalho à direita: número de página
\fancyhead[R]{\small \thepage}

% Linha horizontal no cabeçalho (espessura 0.4pt)
\renewcommand{\headrulewidth}{0.4pt}

% Sobrescreve o título padrão “References” para “Referências”
%\renewcommand{\refname}{Referências}
% Tradução de autoref para o português
\def\figureautorefname{Figura}
\def\tableautorefname{Tabela}
\def\equationautorefname{Equa\c{c}\~ao}

% \hypersetup{backref=true}

\ifdefined\backref
        \renewcommand*{\backref}[1]{}
        \renewcommand*{\backrefentrycount}[2]{#1\ifnum#2>1 ~(#2)\fi}
\fi

\DeclarePairedDelimiter{\ceil}{\lceil}{\rceil}

\onehalfspacing

\usepackage{pgfgantt}

%++++++++++++++++++++++++++++++++++++++++++++++++++++++++++++++++++++++++++++++
% Identificacoes (se o trabalho for em inglês, insira os dados em inglês)
% Para entradas abreviadas de Professora (Profa.) em português escreva: Prof$^\textnormal{a}$.
%++++++++++++++++++++++++++++++++++++++++++++++++++++++++++++++++++++++++++++++
\course{Engenharia Aeronáutica e Mecânica} % Programa de PG ou Curso de Graduação
\area{Propulsão Aeroespacial e Energia} % Área de concentração na PG (Não utilizado no caso de TG)

% Autor do trabalho: Nome Sobrenome
\authorgender{masc} %sexo: masc ou fem
\author{\textbf{EDEN DENIS} Ferreira da Silva Lopes}{Santos}
\itaauthoraddress{Rua João Fonseca dos Santos, 20, Apto. 93}{12.230-088}{São José dos Campos--SP}

% Titulo da Tese/Dissertação
\title{Aplicação de Inteligência Artificial (AI) em Banco de Dados de Emissão de Motores}

% Orientador
\advisorgender{masc} % masc ou fem
\advisor{Prof.~Dr.}{Pedro Teixeira \textbf{LACAVA}}{ITA}

% Coorientador (Caso não haja coorientador, colocar ambas as variáveis \coadvisorgender e \coadvisor comentadas, com um % na frente)
\coadvisorgender{masc}
\coadvisor{}{}  % {Prof.~Dr.}{Carlos Alberto Rocha \textbf{PIMENTEL}}{ITA}

% Pró-reitor da Pós-graduação
\bossgender{fem}
\boss{Prof.~Dra.}{Emília \textbf{VILLANI}}

%Coordenador do curso no caso de TG
\bosscoursegender{masc}
\bosscourse{Prof.~Dr.}{ - } % masc ou fem

% Palavras-Chaves informadas pela Biblioteca -> utilizada na CIP
\kwcip{Engine taxonomy}
\kwcip{Emission classification}
\kwcip{Inteligência Artificial}
\kwcip{Artificial Intelligence}
\kwcip{Python}

% membros da banca examinadora
\examiner{Prof. Dr.}{Pedro Teixeira \textbf{LACAVA}}{President}{ITA}
\examiner{Prof. Dr.}{}{Internal Member}{ITA} % Vinícius \textbf{MALATESTA}}{Internal Member}{ITA}
\examiner{Prof. Dr.}{}{Internal Member}{ITA} % {Elisan dos Santos \textbf{MAGALHÃES}}{Internal Member}{ITA}
\examiner{Prof. Dr.}{}{Internal Member}{ITA} % {Carlos Alberto Rocha \textbf{PIMENTEL}}{External Member}{UFABC}
\examiner{Prof. Dr.}{}{Internal Member}{ITA} % {\textbf{KARL} Peter Burr}{External Member}{UFABC}

% Data da defesa (mês em maiúsculo, se trabalho em inglês, e minúsculo se trabalho em português) 
\date{14}{JULHO}{2025}

% Número CDU - (somente para TG)
\cdu{621.38}

% Folha de Registro do Documento
\newcommand{\advisorTitle}{Prof.~Dr.}
\newcommand{\advisorName}{Pedro Teixeira \textbf{LACAVA}}

\newcommand{\coadvisorTitle}{Prof.~Dr.}
\newcommand{\coadvisorName}{Carlos Alberto Rocha \textbf{PIMENTEL}}

\newcommand{\dateDay}{14}
\newcommand{\dateMonth}{JULHO}
\newcommand{\dateYear}{2025}

% Data de publicação e defesa (usada na folha de registro)
\newcommand{\frdDataPublicacao}{04 de agosto de 2025}
\newcommand{\frdNumeroRegistro}{DCTA/ITA/TD-029/2020}
\newcommand{\frdPalavrasAutor}{Big Data; Engine Emission; Emission Databank; Clustering; Artificial Intelligence}
\newcommand{\frdPalavrasResultado}{Taxonomia; Emissão de Motores; Planejamento de Dados; Inteligência Artificial; Processamento de Dados}
\newcommand{\frdApresentacao}{
	ITA, São José dos Campos. Curso de Mestrado. Programa de Pós-Graduação em Engenharia Aeronáutica e Mecânica. Área de Propulsão Aeroespacial e Energia.\\
	Orientador: \advisorTitle~\advisorName.\\
	Coorientador: \coadvisorTitle~\coadvisorName.\\
	Defesa em \dateDay/\dateMonth/\dateYear. Publicada em \frdDataPublicacao.
}

\newcommand{\PorcentagemDaLargura}{0.80}

\begin{document}
	
	\newpage
	%\shipout\null % Inserir página em branco
	\chapter{Outras Aplicações em Engenharia}
	
		\section{Outras Aplicações em Engenharia}\label{sec:OutrasAplicacoesEmEngenharia}
		
		Dentre as possibilidades de aplicações, tem-se:
		
		\begin{enumerate}
	                \item Métodos para similaridade de grupos;
	
	                \item Previsão não-linear;
	
	                \item Aproximação de distribuição multivariada;
	
	                \item Testes não-paramétricos de independência entre várias variáveis;
	
	                \item Otimização de rotas logísticas;
	
	                \item Monitoramento de integridade estrutural;
	
	                \item Detecção de anomalias em sensores;
	
	                \item Planejamento de manutenção preditiva;
	
	                \item Classificação de motores de aeronaves a partir das emissões;
	
	                \item Análise de BIG Data para desempenho de sistemas propulsivos;
	
	                \item Integração de técnicas de Inteligência Artificial em projetos de turbinas;
	
	                \item Processamento paralelo para otimizar algoritmos de agrupamento;
	
	                \item Aplicações multiplataforma para avaliação de motores de aeronaves.
	        \end{enumerate}
	
\end{document}

	
	% Apêndices:
    \appendix
\chapter{Apêndice}

Para garantir transparência e reprodutibilidade, listamos a seguir os principais arquivos dos códigos-fonte e seus pré e pós-processamentos. Os trechos são incluídos com numeração de linhas para referência.

\subsection{\texorpdfstring{\texttt{normal\_sigma\_bands\_outliers.py}}{normal_sigma_bands_outliers.py}}\label{app:normal_sigma_bands_outliers}

    \lstinputlisting[
        caption={Distribuição normal com bandas sigma e outliers (script de gráfico).},
        label={lst:normal_sigma_bands_outliers},
        language=Python,
        inputencoding=utf8,
        basicstyle=\scriptsize\ttfamily,
        numberstyle=\scriptsize\color{Gray}
        ]{\detokenize{scripts/python/normal_sigma_bands_outliers.py}}

\subsection{\texorpdfstring{\texttt{main\_audithas.py}}{main_audithas.py}}\label{app:main_audithas}

    \lstinputlisting[
        caption={Fluxo principal do Audithas, contemplando integrações de pré e pós-processamento.},
        label={lst:main_audithas},
        language=Python,
        inputencoding=utf8,
        basicstyle=\scriptsize\ttfamily,
        numberstyle=\scriptsize\color{Gray}
        ]{\detokenize{subs/submodules/audithas/main_audithas.py}}

\subsection{\texorpdfstring{\texttt{main\_k\_means.py}}{main_k_means.py}}\label{app:main_k_means}

    \lstinputlisting[
        caption={Fluxo principal do módulo \texttt{k\_means\_thesis} e sua orquestração de rotinas.},
        label={lst:main_k_means},
        language=Python,
        inputencoding=utf8,
        basicstyle=\scriptsize\ttfamily,
        numberstyle=\scriptsize\color{Gray}
        ]{\detokenize{subs/submodules/k_means_master_thesis/main_k_means.py}}

	
	% Anexos:
	\documentclass[12pt]{article} % Tamanho de letra 12, papel A4, Estilo de artigo

\usepackage[utf8]{inputenc} % Pacote para permitir o uso de acentos
% Ajustes de hifenização
\tolerance=1
\emergencystretch=\maxdimen
\hyphenpenalty=10000
\hbadness=10000
\usepackage{lipsum} % Pacote para gerar texto aleatório
\usepackage{ragged2e} % Pacote para mais opções de alinhamento
\usepackage{graphicx} % Pacote para inserir figuras jpeg, pdf e png
\usepackage{ae}
\usepackage{epsfig}
\usepackage{mathtools}
\usepackage[export]{adjustbox}
\usepackage{lineno}
\usepackage{amsthm}
\usepackage{commath}
\usepackage{color, colortbl}
\usepackage{bbm}
\usepackage{setspace} % Pacote para ajustar espaçamento entre linhas

% Instructions to prevent LaTeX from hyphenating words
\tolerance=1
\emergencystretch=\maxdimen
\hyphenpenalty=10000
\hbadness=10000

\graphicspath{{figures/}} % Indica o diretório no qual as imagens estão sendo salvas, evitando que precisamos digitar parte do caminho da pasta que contém as imagens
\usepackage[rightcaption]{sidecap} % Pacote para mais opções de figura, colocar, título da legenda da figura ao lado direito
\usepackage{url} % Pacote para citar referências de sites
\usepackage{makeidx} % Pacote para criar índice
\makeindex % Habilitar a criação do índice
\usepackage{longtable} % Pacote para usar tabelas longas

% ----------
% Pacotes para dar quebra de linha dentro da tabela:
%\usepackage{fourier} 
%\usepackage{array}
\usepackage{makecell}

\renewcommand\theadalign{bc}
\renewcommand\theadfont{\bfseries}
\renewcommand\theadgape{\Gape[4pt]}
\renewcommand\cellgape{\Gape[4pt]}
% ----------

\usepackage{wrapfig} % Pacote para alinhar figuras à esquerda ou à direita do texto
\usepackage{ae}
\usepackage{epsfig}
\usepackage{mathtools}
\usepackage[export]{adjustbox}
\usepackage{lineno}
\usepackage{amsfonts}
\usepackage{amsthm}
\usepackage{commath}
\usepackage{color, colortbl}
\usepackage{bbm}

% ----------
% Pacote para inserir cabeçalho e rodapé:
% Atenção! Se remover o cabeçalho e rodapé, o LaTeX coloca como cabeçalho o padrão dele, que é bem interessante também!
% \usepackage{fancyhdr} % Pacote para inserir cabeçalho e rodapé
% \pagestyle{fancy} % Estilo do cabeçalho e rodapé, neste caso apenas uma linha no cabeçalho
%\lhead{{\footnotesize Prandiano\textregistered}} % Inserir texto no canto esquerdo do cabeçalho
%\chead{centro} % Inserir texto no centro do cabeçalho
%\rhead{{\footnotesize Curso de Python}} % Inserir texto no canto direito do cabeçalho
%\lfoot{{\footnotesize Copyright \textcopyright 2018 Prandiano}} % Inserir texto no canto esquerdo do rodapé
%\cfoot{centro} % Inserir texto no centro do rodapé
%\rfoot{{\footnotesize Reprodução proibida: Lei 9.610 de 19/02/1998}} % Inserir texto no canto direito do rodapé
% O comando \renewcommand re-nova o parâmetro especificado
%\renewcommand{\footrulewidth}{0.4pt} % Inserir uma linha acima do rodapé
% ----------

% ----------
% Pacote para inserir múltiplas figuras:
\usepackage{subfigure} % Pacote para inserir múltiplas figuras
% \usepackage{subcaption} % Pacote incompatível com subfigure
%\usepackage{caption}
% ----------

\usepackage{multicol} % Pacote que permite usar multicolunas
\usepackage{float} % Pacote para usar o "H" do figure
\usepackage[dvipsnames]{xcolor} % Pacote de cores para textos e possivelmente outros
\usepackage{amsmath} % Pacote para usar o cases
\usepackage[makeroom]{cancel} % Pacote para usar o cancelar (em e embaixo)
\usepackage{enumitem} % Pacote para enumerar itens
\usepackage{amssymb} % Pacote para inserir os símbolos dos conjuntos numéricos N, Z, Q, R, C
\usepackage[T1]{fontenc}
\usepackage{enumitem}
\usepackage{textcomp}
\usepackage{pdfpages} % Pacote para inserir página em PDF
\usepackage{courier} % Pacote para usar a fonte courier
%\usepackage{matlab-prettifier}
\usepackage{lscape} % Pacote para mudar sentido da página
\newtheorem{theorem}{Teorema} % Pacote para inserir Teorema
\newtheorem{corollary}{Corollary} % Pacote para inserir Corolário
\newtheorem{remark}{Remark}
\newtheorem{mytheorem}{Definition}
\usepackage{hyperref} % Pacote para inserir hyperref
\usepackage{multirow} % Pacote para usar os comandos \multirow
\usepackage{pdfpages} % Pacote para incluir PDF's
\usepackage{listingsutf8} % Pacote para inserir arquivos externos tipo .txt
\usepackage{etoolbox}
\AtBeginEnvironment{verbatim}{\footnotesize}

\usepackage{standalone} % Pacote para inserir sub-arquivos LaTeX
\usepackage{bookmark} % Pacote para NÃO tratar o Índice como Apêndice e colocar letra

% Nome da listagem em português
\renewcommand{\lstlistingname}{Listagem}

% REFERÊNCIAS: ----------
%\usepackage[numbers, compress]{natbib} % Pacote para usar as referências bibliográficas
%\usepackage{usebib}
% \bibliographystyle{references} % Estilo das referências
%\setcitestyle{authoryear,aysep={},open={(},close={)}}
%\newcommand{\myand}{e}

% ----------


% PARÂMETROS DA FUNÇÃO LISTINGS: ----------
\lstset
{
    basicstyle=\small\ttfamily,
    commentstyle=\color{Green},
    keywordstyle=\color{Cerulean},
    frame=single,
    language=Python,
    morekeywords={True, False},
    numbers=left,
    numbersep=10pt,
    numberstyle=\footnotesize\color{Gray},
    showstringspaces=false,
    stringstyle=\color{Mulberry},
    tabsize=3,
    breaklines=true, % Este comando é para quebrar uma linha comprida
    postbreak=\mbox{\textcolor{red}{$\hookrightarrow$}\space}, % Este comando é para indicar com uma seta vermelha que a linha continua embaixo 
    extendedchars=true,
    literate=
        {á}{{\'a}}1 {Á}{{\'A}}1 {à}{{\`a}}1 {À}{{\`A}}1 {ã}{{\~a}}1 {Ã}{{\~A}}1 {â}{{\^a}}1 {Â}{{\^A}}1
        {ç}{{\c{c}}}1 {Ç}{{\c{C}}}1 {é}{{\'e}}1 {É}{{\'E}}1 {ê}{{\^e}}1 {Ê}{{\^E}}1 {í}{{\'\i}}1 {Í}{{\'I}}1
        {ó}{{\'o}}1 {Ó}{{\'O}}1 {ô}{{\^o}}1 {Ô}{{\^O}}1 {õ}{{\~o}}1 {Õ}{{\~O}}1 {ú}{{\'u}}1 {Ú}{{\'U}}1
        {ü}{{\"u}}1 {Ü}{{\"U}}1 {“}{{``}}1 {”}{{''}}1 {’}{{'}}1 {—}{{--}}1 {–}{{--}}1 {−}{{-}}1
        {±}{{$\pm$}}1 {σ}{{$\sigma$}}1 {→}{{$\to$}}1 {…}{{\ldots}}1
}

\lstset{inputencoding=utf8/latin1}

\lstdefinelanguage{Matlab}{
	language=Matlab,
	basicstyle=\ttfamily\small,
	keywordstyle=\color{blue}\bfseries,
	stringstyle=\color{red},
	commentstyle=\color{gray}\itshape,
	morekeywords={function, end, if, else, for, while, switch, case, otherwise},
	sensitive=true
}
% ----------

% ----------
\usepackage{nomencl} % Pacote para criar nomenclatura e lista de nomenclatura
\makenomenclature % Criar a nomenclatura
\renewcommand{\nomname}{Lista de Nomenclaturas} % Renomear o título da nomenclatura
% ----------

\usepackage{amsmath} % Pacote para inserir matrix com chaves, entre outras coisas

\usepackage{booktabs} % Pacote para usar \toprule, \midrule e \bottomrule


% EXECUTAR CÁLCULOS A PARTIR DA LINHA DE COMANDO: ---------
%\usepackage{siunitx, xfp} % Pacote para executar cálculos a partir da linha de comando
%\usepackage{l3regex}%unless you have very new expl3
%%\sisetup{locale=DE}
%%\let\oldnum\num
%%\ExplSyntaxOn
%\def\num#1{\def\tmp{#1}\regex_replace_all:nnN{,([0-9])}{.\1}{\tmp}\oldnum{\tmp}}
%\ExplSyntaxOff
% ----------


%\pagecolor{black} % Alterar a cor da página
%\color{white} % Alterar a cor do texto

% ----------
% Pacotes para inserir algoritmos e pseudo-códigos:
\setlength{\marginparwidth}{2.5cm}
\usepackage[colorinlistoftodos]{todonotes}
\usepackage{algorithm} % Pacote para inserir algoritmos
\usepackage{algpseudocode} % Pacote para inserir algoritmos e pseudo-códigos
\floatname{algorithm}{Algoritmo} % Renomear a palavra Algorithm para Procedimento
\renewcommand{\algorithmicrequire}{\textbf{entrada:}} % Renomear a palavra Input para Entrada
\renewcommand{\algorithmicensure}{\textbf{saída:}} % Renomear a palavra Output para Saída
\renewcommand{\algorithmicwhile}{\textbf{enquanto}} % Renomear a palavra While para Enquanto
\renewcommand{\algorithmicif}{\textbf{se}} % Renomear a palavra If para Se
\renewcommand{\algorithmicend}{\textbf{fim}} % Renomear a palavra End Output para Fim do
\renewcommand{\algorithmicprocedure}{\textbf{procedimento}} % Renomear a palavra Procedure para Procedimento
\renewcommand{\algorithmicfor}{\textbf{para}} % Renomear a palavra For para Para
\renewcommand{\algorithmicelse}{\textbf{senão}} % Renomear a palavra Else para Senão
\renewcommand{\algorithmicfunction}{\textbf{função}} % Renomear a palavra Function para Função 
\renewcommand{\algorithmicthen}{\textbf{então}} % Renomear a palavra Then para Então
\renewcommand{\algorithmicdo}{\textbf{fazer}} % Renomear a palavra Do para Fazer
% ----------

\newcommand{\quotes}[1]{``#1''} % Comando para inserir um texto dentro de aspas-dupla com espaço após o final

% FLUXOGRAMA: ----------
\usepackage{tikz}
\usetikzlibrary{positioning, shapes.geometric}
\tikzstyle{line}=[draw] % here

\tikzstyle{startstop} = [rectangle, rounded corners, minimum width=3cm,
minimum height=1cm,text centered, draw=black, fill=white!30]
\tikzstyle{arrow} = [thick, ->, >=stealth]
\tikzstyle{process} = [rectangle, minimum width=3cm, minimum height=1cm, text centered, draw=black, fill=white!30]
\tikzstyle{io} = [trapezium, trapezium left angle=70, trapezium right angle=110, minimum width=3cm, minimum height=1cm, text centered, draw=black, fill=white!30]
\tikzstyle{connector} = [draw, circle, fill=white!20]
\tikzstyle{decision} = [diamond, minimum width=3cm, minimum height=1cm, text centered, draw=black, fill=white!30, aspect=3]
% ----------


% CAIXA ENTORNO DAS FÓRMULAS: ----------
\usepackage{dashbox}
\newcommand\dboxed[1]{\dbox{\ensuremath{#1}}}

% ----------

%\usepackage{tabularx} % para colunas com largura ajustável

% -------- CONFIGURAÇÃO DE MARGENS --------
% \usepackage[top=2cm,bottom=2cm,left=3cm,right=2cm]{geometry}

% -------- PACOTE fancyhdr --------
\usepackage{fancyhdr}

\setlength{\marginparwidth}{2cm}

% Ajuste do tamanho do cabeçalho para acompanhar exatamente \textwidth
\setlength{\headwidth}{\textwidth}

% Define um headheight mínimo (para evitar avisos)
\setlength{\headheight}{15pt} 

% Limpa configurações padrão de cabeçalho/rodapé
\pagestyle{fancy}
\fancyhf{} 

% Marcações automáticas de seção para o cabeçalho (opcional)
\renewcommand{\sectionmark}[1]{%
	\markboth{\thesection\ #1}{}%
}

% Texto do cabeçalho à esquerda: seção e título
% Usamos \leftmark pois o \sectionmark acima grava no \leftmark
\fancyhead[L]{\small \leftmark}

% Texto do cabeçalho à direita: número de página
\fancyhead[R]{\small \thepage}

% Linha horizontal no cabeçalho (espessura 0.4pt)
\renewcommand{\headrulewidth}{0.4pt}

% Sobrescreve o título padrão “References” para “Referências”
%\renewcommand{\refname}{Referências}
% Tradução de autoref para o português
\def\figureautorefname{Figura}
\def\tableautorefname{Tabela}
\def\equationautorefname{Equa\c{c}\~ao}

% \hypersetup{backref=true}

\ifdefined\backref
        \renewcommand*{\backref}[1]{}
        \renewcommand*{\backrefentrycount}[2]{#1\ifnum#2>1 ~(#2)\fi}
\fi

\DeclarePairedDelimiter{\ceil}{\lceil}{\rceil}

\onehalfspacing

\usepackage{pgfgantt}

%++++++++++++++++++++++++++++++++++++++++++++++++++++++++++++++++++++++++++++++
% Identificacoes (se o trabalho for em inglês, insira os dados em inglês)
% Para entradas abreviadas de Professora (Profa.) em português escreva: Prof$^\textnormal{a}$.
%++++++++++++++++++++++++++++++++++++++++++++++++++++++++++++++++++++++++++++++
\course{Engenharia Aeronáutica e Mecânica} % Programa de PG ou Curso de Graduação
\area{Propulsão Aeroespacial e Energia} % Área de concentração na PG (Não utilizado no caso de TG)

% Autor do trabalho: Nome Sobrenome
\authorgender{masc} %sexo: masc ou fem
\author{\textbf{EDEN DENIS} Ferreira da Silva Lopes}{Santos}
\itaauthoraddress{Rua João Fonseca dos Santos, 20, Apto. 93}{12.230-088}{São José dos Campos--SP}

% Titulo da Tese/Dissertação
\title{Aplicação de Inteligência Artificial (AI) em Banco de Dados de Emissão de Motores}

% Orientador
\advisorgender{masc} % masc ou fem
\advisor{Prof.~Dr.}{Pedro Teixeira \textbf{LACAVA}}{ITA}

% Coorientador (Caso não haja coorientador, colocar ambas as variáveis \coadvisorgender e \coadvisor comentadas, com um % na frente)
\coadvisorgender{masc}
\coadvisor{}{}  % {Prof.~Dr.}{Carlos Alberto Rocha \textbf{PIMENTEL}}{ITA}

% Pró-reitor da Pós-graduação
\bossgender{fem}
\boss{Prof.~Dra.}{Emília \textbf{VILLANI}}

%Coordenador do curso no caso de TG
\bosscoursegender{masc}
\bosscourse{Prof.~Dr.}{ - } % masc ou fem

% Palavras-Chaves informadas pela Biblioteca -> utilizada na CIP
\kwcip{Engine taxonomy}
\kwcip{Emission classification}
\kwcip{Inteligência Artificial}
\kwcip{Artificial Intelligence}
\kwcip{Python}

% membros da banca examinadora
\examiner{Prof. Dr.}{Pedro Teixeira \textbf{LACAVA}}{President}{ITA}
\examiner{Prof. Dr.}{}{Internal Member}{ITA} % Vinícius \textbf{MALATESTA}}{Internal Member}{ITA}
\examiner{Prof. Dr.}{}{Internal Member}{ITA} % {Elisan dos Santos \textbf{MAGALHÃES}}{Internal Member}{ITA}
\examiner{Prof. Dr.}{}{Internal Member}{ITA} % {Carlos Alberto Rocha \textbf{PIMENTEL}}{External Member}{UFABC}
\examiner{Prof. Dr.}{}{Internal Member}{ITA} % {\textbf{KARL} Peter Burr}{External Member}{UFABC}

% Data da defesa (mês em maiúsculo, se trabalho em inglês, e minúsculo se trabalho em português) 
\date{14}{JULHO}{2025}

% Número CDU - (somente para TG)
\cdu{621.38}

% Folha de Registro do Documento
\newcommand{\advisorTitle}{Prof.~Dr.}
\newcommand{\advisorName}{Pedro Teixeira \textbf{LACAVA}}

\newcommand{\coadvisorTitle}{Prof.~Dr.}
\newcommand{\coadvisorName}{Carlos Alberto Rocha \textbf{PIMENTEL}}

\newcommand{\dateDay}{14}
\newcommand{\dateMonth}{JULHO}
\newcommand{\dateYear}{2025}

% Data de publicação e defesa (usada na folha de registro)
\newcommand{\frdDataPublicacao}{04 de agosto de 2025}
\newcommand{\frdNumeroRegistro}{DCTA/ITA/TD-029/2020}
\newcommand{\frdPalavrasAutor}{Big Data; Engine Emission; Emission Databank; Clustering; Artificial Intelligence}
\newcommand{\frdPalavrasResultado}{Taxonomia; Emissão de Motores; Planejamento de Dados; Inteligência Artificial; Processamento de Dados}
\newcommand{\frdApresentacao}{
	ITA, São José dos Campos. Curso de Mestrado. Programa de Pós-Graduação em Engenharia Aeronáutica e Mecânica. Área de Propulsão Aeroespacial e Energia.\\
	Orientador: \advisorTitle~\advisorName.\\
	Coorientador: \coadvisorTitle~\coadvisorName.\\
	Defesa em \dateDay/\dateMonth/\dateYear. Publicada em \frdDataPublicacao.
}

\newcommand{\PorcentagemDaLargura}{0.80}

\begin{document}
	
	% \annex
	\chapter{Anexos} %opcional

		Esta seção reúne materiais complementares que documentam os dados e recursos associados ao trabalho.
	
	\section*{Anexos}

	% \begin{landscape}
		
%		\includepdf[pages=1, scale=1, angle=0, pagecommand=\section{Dados das Aeronaves Civis},offset=0cm 0cm]{data/aeronaves_civis.pdf}
%		
%		\includepdf[pages=2-12, scale=1, angle=0, pagecommand={},offset=0cm 0cm]{data/aeronaves_civis.pdf}
%		
%		\includepdf[pages=1, scale=1, angle=0, pagecommand={\section{Dados das Aeronaves Militares}},offset=0cm 0cm]{data/aeronaves_militares.pdf}
%		
%		\includepdf[pages=2-28, scale=1, angle=0,offset=0cm 0cm]{data/aeronaves_militares.pdf}
		
		
	% \end{landscape}
	
\end{document}

	
	% REFERENCIAS BIBLIOGRAFICAS
%	\input{references.tex}
	\renewcommand\bibname{\itareferencesnamebabel}
	\bibliography{references}	
		
	\documentclass[12pt]{article} % Tamanho de letra 12, papel A4, Estilo de artigo

\usepackage[utf8]{inputenc} % Pacote para permitir o uso de acentos
% Ajustes de hifenização
\tolerance=1
\emergencystretch=\maxdimen
\hyphenpenalty=10000
\hbadness=10000
\usepackage{lipsum} % Pacote para gerar texto aleatório
\usepackage{ragged2e} % Pacote para mais opções de alinhamento
\usepackage{graphicx} % Pacote para inserir figuras jpeg, pdf e png
\usepackage{ae}
\usepackage{epsfig}
\usepackage{mathtools}
\usepackage[export]{adjustbox}
\usepackage{lineno}
\usepackage{amsthm}
\usepackage{commath}
\usepackage{color, colortbl}
\usepackage{bbm}
\usepackage{setspace} % Pacote para ajustar espaçamento entre linhas

% Instructions to prevent LaTeX from hyphenating words
\tolerance=1
\emergencystretch=\maxdimen
\hyphenpenalty=10000
\hbadness=10000

\graphicspath{{figures/}} % Indica o diretório no qual as imagens estão sendo salvas, evitando que precisamos digitar parte do caminho da pasta que contém as imagens
\usepackage[rightcaption]{sidecap} % Pacote para mais opções de figura, colocar, título da legenda da figura ao lado direito
\usepackage{url} % Pacote para citar referências de sites
\usepackage{makeidx} % Pacote para criar índice
\makeindex % Habilitar a criação do índice
\usepackage{longtable} % Pacote para usar tabelas longas

% ----------
% Pacotes para dar quebra de linha dentro da tabela:
%\usepackage{fourier} 
%\usepackage{array}
\usepackage{makecell}

\renewcommand\theadalign{bc}
\renewcommand\theadfont{\bfseries}
\renewcommand\theadgape{\Gape[4pt]}
\renewcommand\cellgape{\Gape[4pt]}
% ----------

\usepackage{wrapfig} % Pacote para alinhar figuras à esquerda ou à direita do texto
\usepackage{ae}
\usepackage{epsfig}
\usepackage{mathtools}
\usepackage[export]{adjustbox}
\usepackage{lineno}
\usepackage{amsfonts}
\usepackage{amsthm}
\usepackage{commath}
\usepackage{color, colortbl}
\usepackage{bbm}

% ----------
% Pacote para inserir cabeçalho e rodapé:
% Atenção! Se remover o cabeçalho e rodapé, o LaTeX coloca como cabeçalho o padrão dele, que é bem interessante também!
% \usepackage{fancyhdr} % Pacote para inserir cabeçalho e rodapé
% \pagestyle{fancy} % Estilo do cabeçalho e rodapé, neste caso apenas uma linha no cabeçalho
%\lhead{{\footnotesize Prandiano\textregistered}} % Inserir texto no canto esquerdo do cabeçalho
%\chead{centro} % Inserir texto no centro do cabeçalho
%\rhead{{\footnotesize Curso de Python}} % Inserir texto no canto direito do cabeçalho
%\lfoot{{\footnotesize Copyright \textcopyright 2018 Prandiano}} % Inserir texto no canto esquerdo do rodapé
%\cfoot{centro} % Inserir texto no centro do rodapé
%\rfoot{{\footnotesize Reprodução proibida: Lei 9.610 de 19/02/1998}} % Inserir texto no canto direito do rodapé
% O comando \renewcommand re-nova o parâmetro especificado
%\renewcommand{\footrulewidth}{0.4pt} % Inserir uma linha acima do rodapé
% ----------

% ----------
% Pacote para inserir múltiplas figuras:
\usepackage{subfigure} % Pacote para inserir múltiplas figuras
% \usepackage{subcaption} % Pacote incompatível com subfigure
%\usepackage{caption}
% ----------

\usepackage{multicol} % Pacote que permite usar multicolunas
\usepackage{float} % Pacote para usar o "H" do figure
\usepackage[dvipsnames]{xcolor} % Pacote de cores para textos e possivelmente outros
\usepackage{amsmath} % Pacote para usar o cases
\usepackage[makeroom]{cancel} % Pacote para usar o cancelar (em e embaixo)
\usepackage{enumitem} % Pacote para enumerar itens
\usepackage{amssymb} % Pacote para inserir os símbolos dos conjuntos numéricos N, Z, Q, R, C
\usepackage[T1]{fontenc}
\usepackage{enumitem}
\usepackage{textcomp}
\usepackage{pdfpages} % Pacote para inserir página em PDF
\usepackage{courier} % Pacote para usar a fonte courier
%\usepackage{matlab-prettifier}
\usepackage{lscape} % Pacote para mudar sentido da página
\newtheorem{theorem}{Teorema} % Pacote para inserir Teorema
\newtheorem{corollary}{Corollary} % Pacote para inserir Corolário
\newtheorem{remark}{Remark}
\newtheorem{mytheorem}{Definition}
\usepackage{hyperref} % Pacote para inserir hyperref
\usepackage{multirow} % Pacote para usar os comandos \multirow
\usepackage{pdfpages} % Pacote para incluir PDF's
\usepackage{listingsutf8} % Pacote para inserir arquivos externos tipo .txt
\usepackage{etoolbox}
\AtBeginEnvironment{verbatim}{\footnotesize}

\usepackage{standalone} % Pacote para inserir sub-arquivos LaTeX
\usepackage{bookmark} % Pacote para NÃO tratar o Índice como Apêndice e colocar letra

% Nome da listagem em português
\renewcommand{\lstlistingname}{Listagem}

% REFERÊNCIAS: ----------
%\usepackage[numbers, compress]{natbib} % Pacote para usar as referências bibliográficas
%\usepackage{usebib}
% \bibliographystyle{references} % Estilo das referências
%\setcitestyle{authoryear,aysep={},open={(},close={)}}
%\newcommand{\myand}{e}

% ----------


% PARÂMETROS DA FUNÇÃO LISTINGS: ----------
\lstset
{
    basicstyle=\small\ttfamily,
    commentstyle=\color{Green},
    keywordstyle=\color{Cerulean},
    frame=single,
    language=Python,
    morekeywords={True, False},
    numbers=left,
    numbersep=10pt,
    numberstyle=\footnotesize\color{Gray},
    showstringspaces=false,
    stringstyle=\color{Mulberry},
    tabsize=3,
    breaklines=true, % Este comando é para quebrar uma linha comprida
    postbreak=\mbox{\textcolor{red}{$\hookrightarrow$}\space}, % Este comando é para indicar com uma seta vermelha que a linha continua embaixo 
    extendedchars=true,
    literate=
        {á}{{\'a}}1 {Á}{{\'A}}1 {à}{{\`a}}1 {À}{{\`A}}1 {ã}{{\~a}}1 {Ã}{{\~A}}1 {â}{{\^a}}1 {Â}{{\^A}}1
        {ç}{{\c{c}}}1 {Ç}{{\c{C}}}1 {é}{{\'e}}1 {É}{{\'E}}1 {ê}{{\^e}}1 {Ê}{{\^E}}1 {í}{{\'\i}}1 {Í}{{\'I}}1
        {ó}{{\'o}}1 {Ó}{{\'O}}1 {ô}{{\^o}}1 {Ô}{{\^O}}1 {õ}{{\~o}}1 {Õ}{{\~O}}1 {ú}{{\'u}}1 {Ú}{{\'U}}1
        {ü}{{\"u}}1 {Ü}{{\"U}}1 {“}{{``}}1 {”}{{''}}1 {’}{{'}}1 {—}{{--}}1 {–}{{--}}1 {−}{{-}}1
        {±}{{$\pm$}}1 {σ}{{$\sigma$}}1 {→}{{$\to$}}1 {…}{{\ldots}}1
}

\lstset{inputencoding=utf8/latin1}

\lstdefinelanguage{Matlab}{
	language=Matlab,
	basicstyle=\ttfamily\small,
	keywordstyle=\color{blue}\bfseries,
	stringstyle=\color{red},
	commentstyle=\color{gray}\itshape,
	morekeywords={function, end, if, else, for, while, switch, case, otherwise},
	sensitive=true
}
% ----------

% ----------
\usepackage{nomencl} % Pacote para criar nomenclatura e lista de nomenclatura
\makenomenclature % Criar a nomenclatura
\renewcommand{\nomname}{Lista de Nomenclaturas} % Renomear o título da nomenclatura
% ----------

\usepackage{amsmath} % Pacote para inserir matrix com chaves, entre outras coisas

\usepackage{booktabs} % Pacote para usar \toprule, \midrule e \bottomrule


% EXECUTAR CÁLCULOS A PARTIR DA LINHA DE COMANDO: ---------
%\usepackage{siunitx, xfp} % Pacote para executar cálculos a partir da linha de comando
%\usepackage{l3regex}%unless you have very new expl3
%%\sisetup{locale=DE}
%%\let\oldnum\num
%%\ExplSyntaxOn
%\def\num#1{\def\tmp{#1}\regex_replace_all:nnN{,([0-9])}{.\1}{\tmp}\oldnum{\tmp}}
%\ExplSyntaxOff
% ----------


%\pagecolor{black} % Alterar a cor da página
%\color{white} % Alterar a cor do texto

% ----------
% Pacotes para inserir algoritmos e pseudo-códigos:
\setlength{\marginparwidth}{2.5cm}
\usepackage[colorinlistoftodos]{todonotes}
\usepackage{algorithm} % Pacote para inserir algoritmos
\usepackage{algpseudocode} % Pacote para inserir algoritmos e pseudo-códigos
\floatname{algorithm}{Algoritmo} % Renomear a palavra Algorithm para Procedimento
\renewcommand{\algorithmicrequire}{\textbf{entrada:}} % Renomear a palavra Input para Entrada
\renewcommand{\algorithmicensure}{\textbf{saída:}} % Renomear a palavra Output para Saída
\renewcommand{\algorithmicwhile}{\textbf{enquanto}} % Renomear a palavra While para Enquanto
\renewcommand{\algorithmicif}{\textbf{se}} % Renomear a palavra If para Se
\renewcommand{\algorithmicend}{\textbf{fim}} % Renomear a palavra End Output para Fim do
\renewcommand{\algorithmicprocedure}{\textbf{procedimento}} % Renomear a palavra Procedure para Procedimento
\renewcommand{\algorithmicfor}{\textbf{para}} % Renomear a palavra For para Para
\renewcommand{\algorithmicelse}{\textbf{senão}} % Renomear a palavra Else para Senão
\renewcommand{\algorithmicfunction}{\textbf{função}} % Renomear a palavra Function para Função 
\renewcommand{\algorithmicthen}{\textbf{então}} % Renomear a palavra Then para Então
\renewcommand{\algorithmicdo}{\textbf{fazer}} % Renomear a palavra Do para Fazer
% ----------

\newcommand{\quotes}[1]{``#1''} % Comando para inserir um texto dentro de aspas-dupla com espaço após o final

% FLUXOGRAMA: ----------
\usepackage{tikz}
\usetikzlibrary{positioning, shapes.geometric}
\tikzstyle{line}=[draw] % here

\tikzstyle{startstop} = [rectangle, rounded corners, minimum width=3cm,
minimum height=1cm,text centered, draw=black, fill=white!30]
\tikzstyle{arrow} = [thick, ->, >=stealth]
\tikzstyle{process} = [rectangle, minimum width=3cm, minimum height=1cm, text centered, draw=black, fill=white!30]
\tikzstyle{io} = [trapezium, trapezium left angle=70, trapezium right angle=110, minimum width=3cm, minimum height=1cm, text centered, draw=black, fill=white!30]
\tikzstyle{connector} = [draw, circle, fill=white!20]
\tikzstyle{decision} = [diamond, minimum width=3cm, minimum height=1cm, text centered, draw=black, fill=white!30, aspect=3]
% ----------


% CAIXA ENTORNO DAS FÓRMULAS: ----------
\usepackage{dashbox}
\newcommand\dboxed[1]{\dbox{\ensuremath{#1}}}

% ----------

%\usepackage{tabularx} % para colunas com largura ajustável

% -------- CONFIGURAÇÃO DE MARGENS --------
% \usepackage[top=2cm,bottom=2cm,left=3cm,right=2cm]{geometry}

% -------- PACOTE fancyhdr --------
\usepackage{fancyhdr}

\setlength{\marginparwidth}{2cm}

% Ajuste do tamanho do cabeçalho para acompanhar exatamente \textwidth
\setlength{\headwidth}{\textwidth}

% Define um headheight mínimo (para evitar avisos)
\setlength{\headheight}{15pt} 

% Limpa configurações padrão de cabeçalho/rodapé
\pagestyle{fancy}
\fancyhf{} 

% Marcações automáticas de seção para o cabeçalho (opcional)
\renewcommand{\sectionmark}[1]{%
	\markboth{\thesection\ #1}{}%
}

% Texto do cabeçalho à esquerda: seção e título
% Usamos \leftmark pois o \sectionmark acima grava no \leftmark
\fancyhead[L]{\small \leftmark}

% Texto do cabeçalho à direita: número de página
\fancyhead[R]{\small \thepage}

% Linha horizontal no cabeçalho (espessura 0.4pt)
\renewcommand{\headrulewidth}{0.4pt}

% Sobrescreve o título padrão “References” para “Referências”
%\renewcommand{\refname}{Referências}
% Tradução de autoref para o português
\def\figureautorefname{Figura}
\def\tableautorefname{Tabela}
\def\equationautorefname{Equa\c{c}\~ao}

% \hypersetup{backref=true}

\ifdefined\backref
        \renewcommand*{\backref}[1]{}
        \renewcommand*{\backrefentrycount}[2]{#1\ifnum#2>1 ~(#2)\fi}
\fi

\DeclarePairedDelimiter{\ceil}{\lceil}{\rceil}

\onehalfspacing

\usepackage{pgfgantt}

%++++++++++++++++++++++++++++++++++++++++++++++++++++++++++++++++++++++++++++++
% Identificacoes (se o trabalho for em inglês, insira os dados em inglês)
% Para entradas abreviadas de Professora (Profa.) em português escreva: Prof$^\textnormal{a}$.
%++++++++++++++++++++++++++++++++++++++++++++++++++++++++++++++++++++++++++++++
\course{Engenharia Aeronáutica e Mecânica} % Programa de PG ou Curso de Graduação
\area{Propulsão Aeroespacial e Energia} % Área de concentração na PG (Não utilizado no caso de TG)

% Autor do trabalho: Nome Sobrenome
\authorgender{masc} %sexo: masc ou fem
\author{\textbf{EDEN DENIS} Ferreira da Silva Lopes}{Santos}
\itaauthoraddress{Rua João Fonseca dos Santos, 20, Apto. 93}{12.230-088}{São José dos Campos--SP}

% Titulo da Tese/Dissertação
\title{Aplicação de Inteligência Artificial (AI) em Banco de Dados de Emissão de Motores}

% Orientador
\advisorgender{masc} % masc ou fem
\advisor{Prof.~Dr.}{Pedro Teixeira \textbf{LACAVA}}{ITA}

% Coorientador (Caso não haja coorientador, colocar ambas as variáveis \coadvisorgender e \coadvisor comentadas, com um % na frente)
\coadvisorgender{masc}
\coadvisor{}{}  % {Prof.~Dr.}{Carlos Alberto Rocha \textbf{PIMENTEL}}{ITA}

% Pró-reitor da Pós-graduação
\bossgender{fem}
\boss{Prof.~Dra.}{Emília \textbf{VILLANI}}

%Coordenador do curso no caso de TG
\bosscoursegender{masc}
\bosscourse{Prof.~Dr.}{ - } % masc ou fem

% Palavras-Chaves informadas pela Biblioteca -> utilizada na CIP
\kwcip{Engine taxonomy}
\kwcip{Emission classification}
\kwcip{Inteligência Artificial}
\kwcip{Artificial Intelligence}
\kwcip{Python}

% membros da banca examinadora
\examiner{Prof. Dr.}{Pedro Teixeira \textbf{LACAVA}}{President}{ITA}
\examiner{Prof. Dr.}{}{Internal Member}{ITA} % Vinícius \textbf{MALATESTA}}{Internal Member}{ITA}
\examiner{Prof. Dr.}{}{Internal Member}{ITA} % {Elisan dos Santos \textbf{MAGALHÃES}}{Internal Member}{ITA}
\examiner{Prof. Dr.}{}{Internal Member}{ITA} % {Carlos Alberto Rocha \textbf{PIMENTEL}}{External Member}{UFABC}
\examiner{Prof. Dr.}{}{Internal Member}{ITA} % {\textbf{KARL} Peter Burr}{External Member}{UFABC}

% Data da defesa (mês em maiúsculo, se trabalho em inglês, e minúsculo se trabalho em português) 
\date{14}{JULHO}{2025}

% Número CDU - (somente para TG)
\cdu{621.38}

% Folha de Registro do Documento
\newcommand{\advisorTitle}{Prof.~Dr.}
\newcommand{\advisorName}{Pedro Teixeira \textbf{LACAVA}}

\newcommand{\coadvisorTitle}{Prof.~Dr.}
\newcommand{\coadvisorName}{Carlos Alberto Rocha \textbf{PIMENTEL}}

\newcommand{\dateDay}{14}
\newcommand{\dateMonth}{JULHO}
\newcommand{\dateYear}{2025}

% Data de publicação e defesa (usada na folha de registro)
\newcommand{\frdDataPublicacao}{04 de agosto de 2025}
\newcommand{\frdNumeroRegistro}{DCTA/ITA/TD-029/2020}
\newcommand{\frdPalavrasAutor}{Big Data; Engine Emission; Emission Databank; Clustering; Artificial Intelligence}
\newcommand{\frdPalavrasResultado}{Taxonomia; Emissão de Motores; Planejamento de Dados; Inteligência Artificial; Processamento de Dados}
\newcommand{\frdApresentacao}{
	ITA, São José dos Campos. Curso de Mestrado. Programa de Pós-Graduação em Engenharia Aeronáutica e Mecânica. Área de Propulsão Aeroespacial e Energia.\\
	Orientador: \advisorTitle~\advisorName.\\
	Coorientador: \coadvisorTitle~\coadvisorName.\\
	Defesa em \dateDay/\dateMonth/\dateYear. Publicada em \frdDataPublicacao.
}

\newcommand{\PorcentagemDaLargura}{0.80}

\begin{document}
	
	% Folha de Registro do Documento
	% Valores dos campos do formulario
	\FRDitadata{\frdDataPublicacao}
	\FRDitadocnro{\frdNumeroRegistro} %(o número de registro você solicita a biblioteca)
	\FRDitaorgaointerno{Instituto Tecnológico de Aeronáutica -- ITA}
	%Exemplo no caso de pós-graduação: Instituto Tecnol{\'o}gico de Aeron{\'a}utica -- ITA
	\FRDitapalavrasautor{\frdPalavrasAutor}
	\FRDitapalavrasresult{\frdPalavrasResultado}
	%Exemplo no caso de graduação (TG):
	%\FRDitapalavraapresentacao{Trabalho de Graduação, ITA, São José dos Campos, 2015. \NumPenultimaPagina\ páginas.}
	%Exemplo no caso de pós-graduação (msc, dsc):
	\FRDitapalavraapresentacao{ITA, São José dos Campos. Curso de Doutorado. Programa de Pós-Graduação em Engenharia Eletrônica e de Computação. Área de Informática. Orientador: Prof.~Dr. Nei Yoshihiro Soma. Coorientador: Prof.~Dr. Rubens Junqueira Magalhães Afonso. Defesa em 30/11/2020. Publicada em 11/01/2021.}
	%\FRDitaresumo{\input{Cap0/resumo}}
	\FRDitaresumo{\documentclass[12pt]{article} % Tamanho de letra 12, papel A4, Estilo de artigo

\usepackage[utf8]{inputenc} % Pacote para permitir o uso de acentos
% Ajustes de hifenização
\tolerance=1
\emergencystretch=\maxdimen
\hyphenpenalty=10000
\hbadness=10000
\usepackage{lipsum} % Pacote para gerar texto aleatório
\usepackage{ragged2e} % Pacote para mais opções de alinhamento
\usepackage{graphicx} % Pacote para inserir figuras jpeg, pdf e png
\usepackage{ae}
\usepackage{epsfig}
\usepackage{mathtools}
\usepackage[export]{adjustbox}
\usepackage{lineno}
\usepackage{amsthm}
\usepackage{commath}
\usepackage{color, colortbl}
\usepackage{bbm}
\usepackage{setspace} % Pacote para ajustar espaçamento entre linhas

% Instructions to prevent LaTeX from hyphenating words
\tolerance=1
\emergencystretch=\maxdimen
\hyphenpenalty=10000
\hbadness=10000

\graphicspath{{figures/}} % Indica o diretório no qual as imagens estão sendo salvas, evitando que precisamos digitar parte do caminho da pasta que contém as imagens
\usepackage[rightcaption]{sidecap} % Pacote para mais opções de figura, colocar, título da legenda da figura ao lado direito
\usepackage{url} % Pacote para citar referências de sites
\usepackage{makeidx} % Pacote para criar índice
\makeindex % Habilitar a criação do índice
\usepackage{longtable} % Pacote para usar tabelas longas

% ----------
% Pacotes para dar quebra de linha dentro da tabela:
%\usepackage{fourier} 
%\usepackage{array}
\usepackage{makecell}

\renewcommand\theadalign{bc}
\renewcommand\theadfont{\bfseries}
\renewcommand\theadgape{\Gape[4pt]}
\renewcommand\cellgape{\Gape[4pt]}
% ----------

\usepackage{wrapfig} % Pacote para alinhar figuras à esquerda ou à direita do texto
\usepackage{ae}
\usepackage{epsfig}
\usepackage{mathtools}
\usepackage[export]{adjustbox}
\usepackage{lineno}
\usepackage{amsfonts}
\usepackage{amsthm}
\usepackage{commath}
\usepackage{color, colortbl}
\usepackage{bbm}

% ----------
% Pacote para inserir cabeçalho e rodapé:
% Atenção! Se remover o cabeçalho e rodapé, o LaTeX coloca como cabeçalho o padrão dele, que é bem interessante também!
% \usepackage{fancyhdr} % Pacote para inserir cabeçalho e rodapé
% \pagestyle{fancy} % Estilo do cabeçalho e rodapé, neste caso apenas uma linha no cabeçalho
%\lhead{{\footnotesize Prandiano\textregistered}} % Inserir texto no canto esquerdo do cabeçalho
%\chead{centro} % Inserir texto no centro do cabeçalho
%\rhead{{\footnotesize Curso de Python}} % Inserir texto no canto direito do cabeçalho
%\lfoot{{\footnotesize Copyright \textcopyright 2018 Prandiano}} % Inserir texto no canto esquerdo do rodapé
%\cfoot{centro} % Inserir texto no centro do rodapé
%\rfoot{{\footnotesize Reprodução proibida: Lei 9.610 de 19/02/1998}} % Inserir texto no canto direito do rodapé
% O comando \renewcommand re-nova o parâmetro especificado
%\renewcommand{\footrulewidth}{0.4pt} % Inserir uma linha acima do rodapé
% ----------

% ----------
% Pacote para inserir múltiplas figuras:
\usepackage{subfigure} % Pacote para inserir múltiplas figuras
% \usepackage{subcaption} % Pacote incompatível com subfigure
%\usepackage{caption}
% ----------

\usepackage{multicol} % Pacote que permite usar multicolunas
\usepackage{float} % Pacote para usar o "H" do figure
\usepackage[dvipsnames]{xcolor} % Pacote de cores para textos e possivelmente outros
\usepackage{amsmath} % Pacote para usar o cases
\usepackage[makeroom]{cancel} % Pacote para usar o cancelar (em e embaixo)
\usepackage{enumitem} % Pacote para enumerar itens
\usepackage{amssymb} % Pacote para inserir os símbolos dos conjuntos numéricos N, Z, Q, R, C
\usepackage[T1]{fontenc}
\usepackage{enumitem}
\usepackage{textcomp}
\usepackage{pdfpages} % Pacote para inserir página em PDF
\usepackage{courier} % Pacote para usar a fonte courier
%\usepackage{matlab-prettifier}
\usepackage{lscape} % Pacote para mudar sentido da página
\newtheorem{theorem}{Teorema} % Pacote para inserir Teorema
\newtheorem{corollary}{Corollary} % Pacote para inserir Corolário
\newtheorem{remark}{Remark}
\newtheorem{mytheorem}{Definition}
\usepackage{hyperref} % Pacote para inserir hyperref
\usepackage{multirow} % Pacote para usar os comandos \multirow
\usepackage{pdfpages} % Pacote para incluir PDF's
\usepackage{listingsutf8} % Pacote para inserir arquivos externos tipo .txt
\usepackage{etoolbox}
\AtBeginEnvironment{verbatim}{\footnotesize}

\usepackage{standalone} % Pacote para inserir sub-arquivos LaTeX
\usepackage{bookmark} % Pacote para NÃO tratar o Índice como Apêndice e colocar letra

% Nome da listagem em português
\renewcommand{\lstlistingname}{Listagem}

% REFERÊNCIAS: ----------
%\usepackage[numbers, compress]{natbib} % Pacote para usar as referências bibliográficas
%\usepackage{usebib}
% \bibliographystyle{references} % Estilo das referências
%\setcitestyle{authoryear,aysep={},open={(},close={)}}
%\newcommand{\myand}{e}

% ----------


% PARÂMETROS DA FUNÇÃO LISTINGS: ----------
\lstset
{
    basicstyle=\small\ttfamily,
    commentstyle=\color{Green},
    keywordstyle=\color{Cerulean},
    frame=single,
    language=Python,
    morekeywords={True, False},
    numbers=left,
    numbersep=10pt,
    numberstyle=\footnotesize\color{Gray},
    showstringspaces=false,
    stringstyle=\color{Mulberry},
    tabsize=3,
    breaklines=true, % Este comando é para quebrar uma linha comprida
    postbreak=\mbox{\textcolor{red}{$\hookrightarrow$}\space}, % Este comando é para indicar com uma seta vermelha que a linha continua embaixo 
    extendedchars=true,
    literate=
        {á}{{\'a}}1 {Á}{{\'A}}1 {à}{{\`a}}1 {À}{{\`A}}1 {ã}{{\~a}}1 {Ã}{{\~A}}1 {â}{{\^a}}1 {Â}{{\^A}}1
        {ç}{{\c{c}}}1 {Ç}{{\c{C}}}1 {é}{{\'e}}1 {É}{{\'E}}1 {ê}{{\^e}}1 {Ê}{{\^E}}1 {í}{{\'\i}}1 {Í}{{\'I}}1
        {ó}{{\'o}}1 {Ó}{{\'O}}1 {ô}{{\^o}}1 {Ô}{{\^O}}1 {õ}{{\~o}}1 {Õ}{{\~O}}1 {ú}{{\'u}}1 {Ú}{{\'U}}1
        {ü}{{\"u}}1 {Ü}{{\"U}}1 {“}{{``}}1 {”}{{''}}1 {’}{{'}}1 {—}{{--}}1 {–}{{--}}1 {−}{{-}}1
        {±}{{$\pm$}}1 {σ}{{$\sigma$}}1 {→}{{$\to$}}1 {…}{{\ldots}}1
}

\lstset{inputencoding=utf8/latin1}

\lstdefinelanguage{Matlab}{
	language=Matlab,
	basicstyle=\ttfamily\small,
	keywordstyle=\color{blue}\bfseries,
	stringstyle=\color{red},
	commentstyle=\color{gray}\itshape,
	morekeywords={function, end, if, else, for, while, switch, case, otherwise},
	sensitive=true
}
% ----------

% ----------
\usepackage{nomencl} % Pacote para criar nomenclatura e lista de nomenclatura
\makenomenclature % Criar a nomenclatura
\renewcommand{\nomname}{Lista de Nomenclaturas} % Renomear o título da nomenclatura
% ----------

\usepackage{amsmath} % Pacote para inserir matrix com chaves, entre outras coisas

\usepackage{booktabs} % Pacote para usar \toprule, \midrule e \bottomrule


% EXECUTAR CÁLCULOS A PARTIR DA LINHA DE COMANDO: ---------
%\usepackage{siunitx, xfp} % Pacote para executar cálculos a partir da linha de comando
%\usepackage{l3regex}%unless you have very new expl3
%%\sisetup{locale=DE}
%%\let\oldnum\num
%%\ExplSyntaxOn
%\def\num#1{\def\tmp{#1}\regex_replace_all:nnN{,([0-9])}{.\1}{\tmp}\oldnum{\tmp}}
%\ExplSyntaxOff
% ----------


%\pagecolor{black} % Alterar a cor da página
%\color{white} % Alterar a cor do texto

% ----------
% Pacotes para inserir algoritmos e pseudo-códigos:
\setlength{\marginparwidth}{2.5cm}
\usepackage[colorinlistoftodos]{todonotes}
\usepackage{algorithm} % Pacote para inserir algoritmos
\usepackage{algpseudocode} % Pacote para inserir algoritmos e pseudo-códigos
\floatname{algorithm}{Algoritmo} % Renomear a palavra Algorithm para Procedimento
\renewcommand{\algorithmicrequire}{\textbf{entrada:}} % Renomear a palavra Input para Entrada
\renewcommand{\algorithmicensure}{\textbf{saída:}} % Renomear a palavra Output para Saída
\renewcommand{\algorithmicwhile}{\textbf{enquanto}} % Renomear a palavra While para Enquanto
\renewcommand{\algorithmicif}{\textbf{se}} % Renomear a palavra If para Se
\renewcommand{\algorithmicend}{\textbf{fim}} % Renomear a palavra End Output para Fim do
\renewcommand{\algorithmicprocedure}{\textbf{procedimento}} % Renomear a palavra Procedure para Procedimento
\renewcommand{\algorithmicfor}{\textbf{para}} % Renomear a palavra For para Para
\renewcommand{\algorithmicelse}{\textbf{senão}} % Renomear a palavra Else para Senão
\renewcommand{\algorithmicfunction}{\textbf{função}} % Renomear a palavra Function para Função 
\renewcommand{\algorithmicthen}{\textbf{então}} % Renomear a palavra Then para Então
\renewcommand{\algorithmicdo}{\textbf{fazer}} % Renomear a palavra Do para Fazer
% ----------

\newcommand{\quotes}[1]{``#1''} % Comando para inserir um texto dentro de aspas-dupla com espaço após o final

% FLUXOGRAMA: ----------
\usepackage{tikz}
\usetikzlibrary{positioning, shapes.geometric}
\tikzstyle{line}=[draw] % here

\tikzstyle{startstop} = [rectangle, rounded corners, minimum width=3cm,
minimum height=1cm,text centered, draw=black, fill=white!30]
\tikzstyle{arrow} = [thick, ->, >=stealth]
\tikzstyle{process} = [rectangle, minimum width=3cm, minimum height=1cm, text centered, draw=black, fill=white!30]
\tikzstyle{io} = [trapezium, trapezium left angle=70, trapezium right angle=110, minimum width=3cm, minimum height=1cm, text centered, draw=black, fill=white!30]
\tikzstyle{connector} = [draw, circle, fill=white!20]
\tikzstyle{decision} = [diamond, minimum width=3cm, minimum height=1cm, text centered, draw=black, fill=white!30, aspect=3]
% ----------


% CAIXA ENTORNO DAS FÓRMULAS: ----------
\usepackage{dashbox}
\newcommand\dboxed[1]{\dbox{\ensuremath{#1}}}

% ----------

%\usepackage{tabularx} % para colunas com largura ajustável

% -------- CONFIGURAÇÃO DE MARGENS --------
% \usepackage[top=2cm,bottom=2cm,left=3cm,right=2cm]{geometry}

% -------- PACOTE fancyhdr --------
\usepackage{fancyhdr}

\setlength{\marginparwidth}{2cm}

% Ajuste do tamanho do cabeçalho para acompanhar exatamente \textwidth
\setlength{\headwidth}{\textwidth}

% Define um headheight mínimo (para evitar avisos)
\setlength{\headheight}{15pt} 

% Limpa configurações padrão de cabeçalho/rodapé
\pagestyle{fancy}
\fancyhf{} 

% Marcações automáticas de seção para o cabeçalho (opcional)
\renewcommand{\sectionmark}[1]{%
	\markboth{\thesection\ #1}{}%
}

% Texto do cabeçalho à esquerda: seção e título
% Usamos \leftmark pois o \sectionmark acima grava no \leftmark
\fancyhead[L]{\small \leftmark}

% Texto do cabeçalho à direita: número de página
\fancyhead[R]{\small \thepage}

% Linha horizontal no cabeçalho (espessura 0.4pt)
\renewcommand{\headrulewidth}{0.4pt}

% Sobrescreve o título padrão “References” para “Referências”
%\renewcommand{\refname}{Referências}
% Tradução de autoref para o português
\def\figureautorefname{Figura}
\def\tableautorefname{Tabela}
\def\equationautorefname{Equa\c{c}\~ao}

% \hypersetup{backref=true}

\ifdefined\backref
        \renewcommand*{\backref}[1]{}
        \renewcommand*{\backrefentrycount}[2]{#1\ifnum#2>1 ~(#2)\fi}
\fi

\DeclarePairedDelimiter{\ceil}{\lceil}{\rceil}

\onehalfspacing

\usepackage{pgfgantt}

%++++++++++++++++++++++++++++++++++++++++++++++++++++++++++++++++++++++++++++++
% Identificacoes (se o trabalho for em inglês, insira os dados em inglês)
% Para entradas abreviadas de Professora (Profa.) em português escreva: Prof$^\textnormal{a}$.
%++++++++++++++++++++++++++++++++++++++++++++++++++++++++++++++++++++++++++++++
\course{Engenharia Aeronáutica e Mecânica} % Programa de PG ou Curso de Graduação
\area{Propulsão Aeroespacial e Energia} % Área de concentração na PG (Não utilizado no caso de TG)

% Autor do trabalho: Nome Sobrenome
\authorgender{masc} %sexo: masc ou fem
\author{\textbf{EDEN DENIS} Ferreira da Silva Lopes}{Santos}
\itaauthoraddress{Rua João Fonseca dos Santos, 20, Apto. 93}{12.230-088}{São José dos Campos--SP}

% Titulo da Tese/Dissertação
\title{Aplicação de Inteligência Artificial (AI) em Banco de Dados de Emissão de Motores}

% Orientador
\advisorgender{masc} % masc ou fem
\advisor{Prof.~Dr.}{Pedro Teixeira \textbf{LACAVA}}{ITA}

% Coorientador (Caso não haja coorientador, colocar ambas as variáveis \coadvisorgender e \coadvisor comentadas, com um % na frente)
\coadvisorgender{masc}
\coadvisor{}{}  % {Prof.~Dr.}{Carlos Alberto Rocha \textbf{PIMENTEL}}{ITA}

% Pró-reitor da Pós-graduação
\bossgender{fem}
\boss{Prof.~Dra.}{Emília \textbf{VILLANI}}

%Coordenador do curso no caso de TG
\bosscoursegender{masc}
\bosscourse{Prof.~Dr.}{ - } % masc ou fem

% Palavras-Chaves informadas pela Biblioteca -> utilizada na CIP
\kwcip{Engine taxonomy}
\kwcip{Emission classification}
\kwcip{Inteligência Artificial}
\kwcip{Artificial Intelligence}
\kwcip{Python}

% membros da banca examinadora
\examiner{Prof. Dr.}{Pedro Teixeira \textbf{LACAVA}}{President}{ITA}
\examiner{Prof. Dr.}{}{Internal Member}{ITA} % Vinícius \textbf{MALATESTA}}{Internal Member}{ITA}
\examiner{Prof. Dr.}{}{Internal Member}{ITA} % {Elisan dos Santos \textbf{MAGALHÃES}}{Internal Member}{ITA}
\examiner{Prof. Dr.}{}{Internal Member}{ITA} % {Carlos Alberto Rocha \textbf{PIMENTEL}}{External Member}{UFABC}
\examiner{Prof. Dr.}{}{Internal Member}{ITA} % {\textbf{KARL} Peter Burr}{External Member}{UFABC}

% Data da defesa (mês em maiúsculo, se trabalho em inglês, e minúsculo se trabalho em português) 
\date{14}{JULHO}{2025}

% Número CDU - (somente para TG)
\cdu{621.38}

% Folha de Registro do Documento
\newcommand{\advisorTitle}{Prof.~Dr.}
\newcommand{\advisorName}{Pedro Teixeira \textbf{LACAVA}}

\newcommand{\coadvisorTitle}{Prof.~Dr.}
\newcommand{\coadvisorName}{Carlos Alberto Rocha \textbf{PIMENTEL}}

\newcommand{\dateDay}{14}
\newcommand{\dateMonth}{JULHO}
\newcommand{\dateYear}{2025}

% Data de publicação e defesa (usada na folha de registro)
\newcommand{\frdDataPublicacao}{04 de agosto de 2025}
\newcommand{\frdNumeroRegistro}{DCTA/ITA/TD-029/2020}
\newcommand{\frdPalavrasAutor}{Big Data; Engine Emission; Emission Databank; Clustering; Artificial Intelligence}
\newcommand{\frdPalavrasResultado}{Taxonomia; Emissão de Motores; Planejamento de Dados; Inteligência Artificial; Processamento de Dados}
\newcommand{\frdApresentacao}{
	ITA, São José dos Campos. Curso de Mestrado. Programa de Pós-Graduação em Engenharia Aeronáutica e Mecânica. Área de Propulsão Aeroespacial e Energia.\\
	Orientador: \advisorTitle~\advisorName.\\
	Coorientador: \coadvisorTitle~\coadvisorName.\\
	Defesa em \dateDay/\dateMonth/\dateYear. Publicada em \frdDataPublicacao.
}

\newcommand{\PorcentagemDaLargura}{0.80}

\begin{document}
% \chapter{Abstract}
	
	\begin{englishabstract}
		\noindent
		
                This work presents a data-driven methodology to classify aircraft engines based on design and thermodynamic characteristics related to efficiency and emissions. The proposed approach aims to identify similarities among different engine configurations using unsupervised machine learning techniques, enabling a systematic comparison of propulsion technologies and supporting technology assessment studies. The analysis is based on publicly available engine data provided by the European Union Aviation Safety Agency (EASA). A set of key design and performance parameters, including overall pressure ratio, turbine inlet temperature, bypass ratio, and other relevant variables, is selected to describe each engine in a multidimensional feature space. Prior to clustering, the dataset is preprocessed through normalization, consistency checks, and the treatment of incomplete entries in order to ensure robustness and comparability among samples. The K-means clustering algorithm is applied to group engines with similar characteristics. Several clustering configurations are evaluated by varying the number of clusters and applying standard validation metrics to identify an adequate partition of the dataset. The resulting clusters are examined to identify technological patterns, distinguish different design philosophies, and highlight engine configurations that exhibit relatively higher efficiency or more favorable emission-related characteristics within their respective groups. The clustering results indicate that engines with comparable thermodynamic and design characteristics tend to form well-defined groups, reflecting similarities in performance level and technological maturity. This classification allows a clearer comparison between engine families and supports the identification of representative configurations within each cluster. In addition, the analysis provides insights into the relative importance of the selected variables in the clustering process, contributing to a better understanding of how design parameters influence engine grouping. The proposed methodology does not rely on predefined classifications and allows patterns to emerge directly from the data. As such, it provides a flexible and reproducible framework for the analysis of aircraft engine technologies. The results may support future studies on emissions assessment, engine benchmarking, and the selection of representative parameters for the development of more efficient and environmentally sustainable aircraft propulsion systems.
	\end{englishabstract}
	
    \textbf{Keywords:} aircraft engine taxonomy, emission classification, Big Data, dendrogram, K-means, Python.
	
\end{document}
}
	%  Primeiro Parametro: Nacional ou Internacional -- N/I
	%  Segundo parametro: Ostensivo, Reservado, Confidencial ou Secreto -- O/R/C/S
	\FRDitaOpcoes{N}{O}
	% Cria o formulario
	\itaFRD
	
\end{document}
	
\end{document}
% Fim do Documento. O massacre acabou!!! :-)
